% =====================================================================
%  R. Nielsen, OM SAMFUNDETS VILKAAR.  Sex folkelige Foredrag.
%  København: Gyldendalske Boghandels Forlag (F. Hegel & Søn); trykt hos J. Jørgensen & Co., 1877.  103 pp.  Antiqua.
%
%  TRANSCRIPTION -- GENERATED; DO NOT EDIT BY HAND.  Rebuild from the repo root (the second --apply places the dashes
%  whose keys are renumbered by the first, p. 15, 65, 73, 102):
%    python3 ebook-method/om-samfundets-vilkaar/draft.py --emit --force \
%      && python3 ebook-method/collate/check.py om-samfundets-vilkaar --italics --apply \
%      && python3 ebook-method/collate/check.py om-samfundets-vilkaar --apply \
%      && python3 ebook-method/om-samfundets-vilkaar/draft.py --post
%
%  STATUS: COMPLETE (2026-10-12): title page, the print's Indhold, and the whole text, pp. 1-103, with its three notes
%  (p. 29: two, p. 98: one).
%
%  EDITION.  The text of record is Google Books' scan of the Andover-Harvard copy (Google id KbQVAAAAYAAJ;
%  ~/bibliotek/Nielsen, Rasmus/samfundets_vilkaar.pdf, SCANS.tsv), 600 ppi bitonal, with Google's OCR as text layer.  The
%  only printing; no e-book and no earlier transcription exists (checked 2026-10-12: the repo, the web, the Internet
%  Archive).  The Internet Archive's copy (omsamfundetsvil00nielgoog) is the same scan; its ABBYY FineReader 8 OCR served
%  as a third reading (abbyy.py), never as a vote.
%
%  ACCURACY (2026-10-12).  The words are those of Google's layer (draft.py), collated word by word with an independent
%  tesseract reading of the images (collate/check.py, witness 'ocr'): 449 disagreements in 25,282 words at the first
%  collation.  check.py's own rules settled 53; rules.py settled 133 more, each rule switched on only after a BLIND PILOT
%  of 102 crops (reader A, 22/22 controls) had measured it on this book: tesseract's glyph faults 8/8, its „ read as a
%  comma 8/8, its » for the print's „ 8/8, accents it adds 8/8, aa read an 5/5, stops before a lower-case word 3/3, the
%  layer's glyph faults 2/2, ø/æ 2/2, case 1/1, § read as $ 3/3; the third reading's two rules 8/8 and 8/8; tesseract's
%  long dashes 7/8 and its lone marks 7/8, each narrowed to exclude its miss; the word-division rule 3/8 stayed off.
%  The rest were read at the image BLIND (sheets.py: 255 word crops with two planted controls a sheet, 48 of 48 controls
%  read right; Sonnet readers), and about 70 doubtful spots in the main session from line crops (zoom.py, 600-dpi crops).
%  STRUCTURE read at the image: the lecture heads and first lines (h1, 21 crops), the title page and Indhold; four places
%  where the layer put words out of line order (pp. 2, 8, 14, 54); words broken at a PAGE end, whose hyphen Google's
%  layer loses (15 pages: draft.py takes the hyphen from tesseract's last line); the notes, set at their marks.
%  LETTERSPACING (the print's emphasis, \emph): candidate lines by the white columns inside words (spacedlines.py), 97
%  lines read blind (14/14 controls) and all seven sheets read again in the main session, which added whole spaced
%  sentences the readers had left roman (pp. 5, 26, 28, 35): 20 runs.  ITALICS by stroke slant: one word (Times, p. 50).
%  SWEEPS after OPEN 0, each spot read at the image: check.py --moved (a line put out of order, p. 2); the split sweep; a
%  review of reads.py's decisions; the quote-balance sweep per paragraph (the five open quotations, pp. 15, 50, 93, 99,
%  100, are the print's: the ABBYY reading counts the same marks); the unknown-word sweep (661 words in no other
%  transcription), which found the misprints and printer's defects below.  SAMPLED PAGE CHECK (pagecheck.py, seed
%  202610126): pp. 2, 24, 70 and the title page read whole and blind, 809 words: 0 errors (95 % upper bound ~4 per
%  1000 words).  TRANSLATION as a last reading: the six translators listed 39 odd spots (ebook-method/
%  om-samfundets-vilkaar/tr/odd-*.tsv and odd-01-*.tsv beside it), all read at the image (zoom zo-01, zo-02): 8 mended by
%  POST patch -- specks both OCRs read as stops (pp. 58, 65, 74, 80, 102), a dash the collation doubled (p. 26), „IIdens“
%  read „Ildens“ (p. 31), a broken m (p. 41) -- and 7 misprints marked; the rest are the print's.
%  What this cannot see: a misprint both OCRs normalise into a real word -- the sweeps find only non-words and nonsense.
%  PRINTED AS IS (each seen at the image; a % comment at the site): „Verdenhistorie“ (p. 4), „sammemkaldte“ (p. 17),
%  „ag“ (p. 21, for og), „var.“ (p. 42) and „Hæder.“ (p. 52) and „Læsere.“ (p. 47) and „Doctriner.“ (p. 58), stops before
%  a lower-case word, „Bourbornerne“ (p. 45), „eg“ (p. 69, for og), „Fremgangmaaden“ and „igjennnem“ (p. 80),
%  „Vilkaarligked“ (p. 85), „Brod“ (p. 86, for Blod), „Forvanksning“ and no stop after „Riger“ (p. 89), „Paaatand“
%  (p. 94), „Deus machina“ (p. 98); also, unmarked, the print's „sig;“ (p. 14), „naturligviis“, „Sverrige“,
%  „Folkevellet“, „Sjouerarbejde“, „kiærest“, and the title page's „København“.
%  PRINTER'S DEFECTS (sorts that did not ink or broke; the letters supplied, a % comment at the site where the word is
%  not otherwise clear): „styrende“ (p. 20, „styr.nde“), „rørig“ (p. 40), „elektrisk“ and „igjennem“ (p. 41,
%  „elcktrisk“, „igjer nem“), „om“ (p. 41, „o n“), „Person“ (p. 44), „Selvbeherskelse“ (p. 48, „Sclvbcherskelse“).
%  Thin word spaces set as spaces: „af Dannelsestrin“ (p. 61), „af Mistillid“ (p. 85), „i Historien“ (p. 97).
%  The ø rule (Hans 2026-10-08) found nothing.
%  Not transcribed: signature marks, library stamps, specks, Google's notice; the large initials of the lecture
%  openings are set as plain letters; the rule under each lecture heading is drawn by \lecture.  The Indhold's dot
%  leaders are set by \dotfill.  Quotation marks as printed („ “).  Notes marked *) and **), as printed.
%
%  PAGE MAP.  printed p = PDF page - 10 (p. 1 = PDF 11, p. 103 = PDF 113); title PDF 7, Indhold PDF 9.
%  Each lecture begins on a new page: I p. 1, II p. 18, III p. 36, IV p. 54, V p. 72, VI p. 89.
%
%  TOOLS (ebook-method/om-samfundets-vilkaar/, copied from liv-og-haab): draft.py, ocr.py, abbyy.py, rules.py (with
%  --classify for the pilot), sheets.py + reads.py + READER-BRIEF.md, spacedlines.py + spaced.tsv, spotsheets.py,
%  zoom.py, pagecheck.py + PAGECHECK-BRIEF.md, patches.py (PATCHES, POST), decisions.tsv, and the sweeps splitsweep.py,
%  linecheck.py, parcheck.py, unknown.py.
%
% =====================================================================
\documentclass[12pt,a4paper,openany]{book}
\usepackage[utf8]{inputenc}
\usepackage[T1]{fontenc}
\usepackage{libertinus}
\usepackage{libertinust1math}
\usepackage[danish]{babel}
\usepackage[protrusion=true,expansion=true]{microtype}
\usepackage{setspace}
\usepackage[margin=1.3in]{geometry}
\usepackage{fancyhdr}
\usepackage{hyperref}
\hypersetup{pdftitle={Om Samfundets Vilkaar}, pdfauthor={Rasmus Nielsen}, hidelinks}
\pagestyle{fancy}
\fancyhf{}
\fancyhead[LE,RO]{\thepage}
\fancyhead[C]{\textit{R. Nielsen: Om Samfundets Vilkaar (1877)}}
\setstretch{1.3}
\usepackage{marginnote}
\renewcommand{\thefootnote}{\ifcase\value{footnote}\or*)\or**)\or***)\fi}   % the print marks a page's notes *) and **); the counter restarts at each page
\newcommand{\tocline}[3]{\noindent\makebox[2.2em][r]{#1}\hspace{0.6em}#2\dotfill\ #3\par}
\newcommand{\opage}[1]{\marginnote{\footnotesize\textit{[#1]}}}
% a lecture opens a new page: its numeral and title, centred, then a short rule; the print's large initial is not reproduced.
\newcommand{\lecture}[2]{\clearpage\phantomsection\addcontentsline{toc}{chapter}{#1 #2}\begin{center}{\large #1}\\[0.6em]{\large #2}\\[0.4em]\rule{3em}{0.4pt}\end{center}\medskip\noindent\ignorespaces}

\begin{document}

% --- front matter, read at the image (front7, front9: reader F, 2026-10-12): PDF 7, title; PDF 9, Indhold.  The title
% page's „København“ is as printed (and as the layer reads it).
\begin{titlepage}
\begin{center}
\vspace*{2em}
{\Huge Om Samfundets Vilkaar.}\\[1.5em]
{\large Sex folkelige Foredrag}\\[1em]
{\small af}\\[1em]
{\large R. Nielsen.}\\[0.3em]
\rule{2em}{0.4pt}\\[2em]
[publisher's device]\\[4em]
{\large København.}\\[0.5em]
Gyldendalske Boghandels Forlag (F. Hegel \& Søn).\\[0.5em]
{\small Trykt hos J. Jørgensen \& Co.}\\[0.5em]
{\small 1877.}
\end{center}
\end{titlepage}

% Indhold (PDF 9): the dot leaders and the stop after each entry cannot be told apart in the scan: the leaders are set by \dotfill.
\begin{center}{\large\emph{Indhold}.}\\[0.4em]\rule{3em}{0.4pt}\end{center}
\hfill{\small Side}\par
\tocline{I.}{Vor Grundlovs Forhold til Øieblikkets politiske Bevægelser (holdt i Kolding)}{1.}
\tocline{II.}{Folkelig Selvstyrelse: Folk og Konge (holdt i Kolding)}{18.}
\tocline{III.}{Politiske Erfaringer: Parlamentarisme (holdt i Kolding)}{36.}
\tocline{IV.}{Fri Forfatning og Frihed (holdt i Middelfart)}{54.}
\tocline{V.}{Aristokratisme og Demokratisme (holdt i Viborg)}{72.}
\tocline{VI.}{Politik og Samvittighed: Macchiavellisme (holdt i Fredericia)}{89.}
\clearpage


% ===== TEXT (pp. 1--103) =====
% --- p. 1 ---
\setcounter{footnote}{0}%
\lecture{I.} {Vor Grundlovs Forhold til Øieblikkets politiske Bevægelser.}
\opage{1}En gammel Fabel melder, at der engang var udbrudt
en almindelig Krig mellem de firføddede Dyr og Fuglene.
Dyreverdenen havde deelt sig i to kæmpende
Hære, som rykkede frem: det gjaldt om Seier eller
Nederlag. Hvert levende Væsen havde taget sit Parti;
dog var der een Skabning, som ikke vilde erklære sig,
men holde sig svævende over begge Partier, det var
Flagermusen. Den tænkte som saa: i urolige Tider
maa man være forsigtig, ikke overile sig, men see Tingene
an og holde gode Miner med Alle. Ingen kan vide,
til hvilken Side Seieren vil helde; den Kloge er forberedt
paa Alt: seire de Firføddede, saa hører jeg i
deres Tal, thi ogsaa jeg er firføddet; men skulde det
skee, at Fuglene fløi af med Seieren, velan, saa har
jeg jo Vinger og kan meget godt gjælde for en Fugl.
Der er desuden noget Ophøiet og Overlegent i, saadan
at hæve sig over Partierne; man er saa uhildet, saa
alsidig, saa hæderlig uinteresseret. Og naar Lidenskaberne
lægge sig, hvad saa? Saa bliver jeg Fredsmægler,
anbefaler gjensidig Imødekommen, indleder
Overeenskomster og er bestandig ovenpaa. Saaledes
tænkte nu Flagermusen; thi Flagermusen var en Po%
% --- p. 2 ---
\setcounter{footnote}{0}%
\opage{2}litiker, en diplomatisk fiin Politiker. Men de stridende
Parter lode sig ikke narre; de mærkede snart
det tvetydige Væsen, og saa uenige de end indbyrdes
vare, bleve de dog uden videre enige om at slaae den
Tvetydige ihjel.

Som det gaaer Flagermusen i Fablen, saaledes
--- mene de meget Ængstelige --- maa det vist gaae hver
den, der i en bevæget Tid vil træde frem for Folket
med et veiledende Ord, hvis han, vel at mærke, ikke
selv vil tage Parti, men indskrænke sig til Forsøg paa
at tyde Bevægelserne og kaste Streiflys over Partiernes
Tumleplads. --- Dog, saa mørke ere Udsigterne ikke.
Vistnok kjender Historien Partikampe, i hvilke Lidenskaberne
skrige saa høit, at fornuftige Grunde maae
tie; men det er Uveirsrædsler, samfundsomvæltende
Jordskjælv. Slige Uveirsrædsler have hidtil, lykkeligviis,
holdt sig udenfor vort Fædrelands Grændser. Vi
vaagne af en gammel Døs, men derfor ville vi dog
ikke fare op i Lidenskab og give os til Priis for
uklare Stemningers blinde Indskydelser. Vi ville tænke,
før vi handle, og prøve, før vi dømme. Endnu staaer
Vidnesbyrdet fast, Historiens hæderlige Vidnesbyrd, at
det danske Folk er et roligt, sindigt, alsidig betænksomt
Folk.

Vi lytte til deres Ord, som ville stifte Fred; men
vi ynde ikke de politiske Flagermuus: de undvige det
klare Lys og sværme helst i Tusmørke. Vi agte paa
Besindighedens Stemme, men vi kunne ikke beundre
de tvetydige Mellemhandlere, hvis Færdighed det er
at bryde Spidsen af ethvert Princip og fordunkle, hvad
der ligger paa Bunden. Vi sætte Priis paa Forlig og
Maadehold, men ere paa vor Post mod tjenstvillige
% --- p. 3 ---
\setcounter{footnote}{0}%
\opage{3}Forligelseskommissærer, der møde med Oplag af allehaande
Meninger, indrettede efter Omstændighederne,
og skifte Farve som det Firbeen man kalder Kamæleon.
Vi have en Forestilling om, at de trods deres tilsyneladende
Upartiskhed sidde dybt indviklede i Partiinteressernes
Væv, at de, reent ud sagt, have Noget
bag Øret. Et Folk sætter aldrig Lid til dem, der
gjøre Indtryk af at have Noget bag Øret. Der er en
folkelig Dyd, der hedder Ærlighed; til Ærlighed svarer
Aabenhed og til Aabenhed den Frimodighed, der i
sømmelig Selvbegrændsning rykker ud med Sproget og
giver os Sagen, som den er, hvad enten den for
Øvrigt er os behagelig eller ubehagelig.

Det er vel efter menneskelige Vilkaar umuligt, at
et levende Folk kan røre sig frit under en fri Forfatning,
uden at der danner sig Partier. Dygtige Partiførere
ere uundværlige Størrelser i Samfundet; ethvert
Parti maa saa staae og falde med den Sag, hvorfor
det kæmper. Men det folkelige Foredrag kan ikke
tage Parti; det ligger i Sagens Natur, det er umuligt.
Med hvem skulde det vel tage Parti? Med Folket! ---
Men Folket som Folk er intet Parti. Det folkelige
Væsen rører sig, mere eller mindre fordunklet, i alle
Partier --- til Høire saa vel som til Venstre. Det
folkelige Væsen omslutter det Hele, forgrener sig i
Alt og gjennemtrænger os alle. Det folkelige Element
indgaaer i de allerforskjelligste Blandinger, det gjælder
om at faae det udskilt af disse Blandinger saa reent
som muligt. Naar to eller flere Partier strides om en
for hele Folket vigtig Sag, og den folkelige Maalestok
lægges an, da er det en Selvfølge, at et Parti kan
blive stillet i et fordeelagtigere Lys end et andet; men
% --- p. 4 ---
\setcounter{footnote}{0}%
\opage{4}heri er ingen Partiskhed. Det, hvorpaa det kommer
an, er alene det: at Maalestokken er rigtig, og Partiet
rigtig bedømt. At fremstille et Parti enten værre eller
bedre end det vitterlig er, er partisk; men uden alle
Bihensyn at ville charakterisere det som det er, er
ikke partisk.

Det vil koste Arbeide, inden den store Almeenhed
naaer saa vidt, at den, uhildet af Partihensyn, kan
skjelne skarpt mellem det Folkelige og det Ufolkelige.
Et glædeligt Tegn er dog allerede givet. Det menige
Folk har begyndt at faae Øinene op for en meget
vigtig Adskillelse, Adskillelsen mellem praktisk Politik
og almeenfattelig Belysning af politiske Tilstande. Et
lille Exempel. Ingen skal med Grund kunne sige, at
Befolkningen i denne Egn mangler praktisk Sands for
Øieblikkets brændende Spørgsmaal. Og Ingen skal
med Grund kunne sige, at der i de Foredrag, jeg forrige
Aar fra dette Sted havde den Ære at holde for
en talrig Kreds af velvillige Tilhørere, blev taget Parti
enten for Høire eller for Venstre. Naar der nu desuagtet
er yttret Ønske om, at de almindelige Synspunkter,
som da bleve fremsatte, maatte ved en videre
Fortsættelse finde nærmere Opklaring og tydeligere
Anvendelse paa givne Tilfælde, da tør jeg vel
uden Anmasselse ansee den hæderlig gjentagne Indbydelse
som et Beviis paa, at den nævnte Adskillelse
begynder at trænge igjennem.

Det folkelige Foredrag kan gaae ad meget forskjellige
Veie, plukke Frugter af forskjellige Træer og
øse af mange Kilder. Det kan hente sine Emner fra
Mythologie, Sagnhistorie, almindelig Verdenhistorie % print: „Verdenhistorie“ (for Verdenshistorie)
o. s. v.; men med alt dette er der dog Et, der aldrig
% --- p. 5 ---
\setcounter{footnote}{0}%
\opage{5}bør tabes af Sigte. Folket lever i det Nærværende.
Det nærværende, virkelige Liv har en Mangfoldighed
af Tilstande, Bryderier, Forviklinger og Spørgsmaal,
som daglig melde sig med stor Paatrængenhed; til
Besvarelse af disse Spørgsmaal maa det folkelige Foredrag
først og sidst, om end kun middelbart, pege hen.
Det folkelige Foredrag skal ikke slaae ned i det Øieblikkelige
med Paastand om, at Dette eller Hiint er
det Eneste, der bør gjøres; det skal meget mere lægge
an paa at opklare de stedfindende Forhold saaledes,
at hver Enkelt i Folket og derved det hele Folk efterhaanden
kan hæve sig til det Trin af Overskuelse og
Indsigt, at Fleertallet og, om muligt, Alle sættes i
Stand til at fælde en fornuftig, selvstændig og begrundet
Dom.

Hvorledes dette efter min Opfattelse, bliver at
forstaae, vil kunne sees, idet jeg nu gaaer over til at
behandle nærværende Foredrags anmeldte Emne: vor
\emph{Grundlovs Forhold til Øieblikkets politiske
Bevægelser}.

Den lille Bog med Overskriften: \emph{Danmarks Riges
gjennemseete Grundlov af 5te Juni 1849},
er vel i Omfang noget nær den mindste, men efter
Indhold den største i vor Litteratur. Det er en Folkebog
saa almeenfattelig, at den kan læses og forstaaes
uden juridisk Lærdom. Heri ligger da en Opfordring
for Hvermand, der vil danne sig en selvstændig Dom
om vor Forfatning, til ikke at lade sig nøie med udrevne
Sætninger, men at kjende Loven i dens Sammenhæng
og gjøre sig fortrolig med dens Aand og Bogstav.

At Grundloven deler den samlede Statsmagt i tre
Magter, den lovgivende, den udøvende og den døm%
% --- p. 6 ---
\setcounter{footnote}{0}%
\opage{6}mende, saaledes, at (§ 2) den lovgivende Magt er hos
Kongen og Rigsdagen i Forening, giver os i denne
Sammenhæng ingen Anledning til særlig Overveielse.
Men Rigsdagens Tvedeling vækker vor Opmærksomhed.
„Rigsdagen“, hedder det (§ 29), „bestaaer af Folkethinget
og Landsthinget.“ Hvad betyder denne Tvedeling?
Hvorfor kunde det ikke være nok med Folkethinget
alene? Den eneste Oplysning vi herom kunne
erholde, maa søges i Valgmaaden (§§ 30---40). Folkethingets
Medlemmer vælges umiddelbart af de i Befolkningen
Valgberettigede „efter Forholdet af et Medlem
til 16000 Indvaanere“. Valget af Landsthingets
Medlemmer er derimod mere begrændset og middelbart
(§ 34); det hele Antal er kun 66, af hvilke
Kongen udnævner 12, o. s. v.

At Grundloven ved denne dobbelte Valgmaade
har villet sikkre sig en saa alsidig og fuldstændig
Repræsentation af den hele Folkevillie, som muligt, er
udenfor al Tvivl. Den har i dette Øiemed opfattet
Folkevillien fra to forskjellige, men dog uadskillelige
Hovedsider, nemlig som Fleertalsvillie og som en i
det Heles Vel begrundet Fornuftvillie. Den har vel
antaget, at Statsstyrelsens fornuftige Element lige saa
fuldt vilde være tilstede i det ene Thing som i det
andet, men derhos tillige antaget, at den gjennem
Fleertallet gaaende Fornuftvillie bedst kunde blive
prøvet, klaret og begrundet, naar der ved en forskjellig
Valgmaade blev sikkret de forskjellige Anskuelser
en behørig Indflydelse paa Forhandlingernes
Gang. Ved at give de umiddelbare Valg en forholdsviis
vid Udstrækning har Loven sørget for, at hele Befolkningen
kan, saa at sige, komme Regjeringen i Tale og
% --- p. 7 ---
\setcounter{footnote}{0}%
\opage{7}igjennem sine valgte Ordførere faae en afgjørende
Stemme med i Lovgivningsvæsenet. Det levende Folk
med dets Stemninger, Tanker og Meninger, dets vexlende
Trang og Fordringer har da hvert Aar sine
Talsmænd paa Rigsdagen: det er, almindelig seet,
Folkethingets Betydning.

Men efter Grundlovens Mening kan Folkethinget ikke
siges at repræsentere Folket i Et og Alt; dette Thing er,
ifølge Valgmaaden, et Udtryk for Folkevillien, nærmest
opfattet som umiddelbar bevægelig Fleertalsvillie. Grundloven
gaaer ud fra, at et Folk ikke altid er tjent med
at faae alt det øieblikkelig sat igjennem, som et til
en given Tid herskende Fleertal fordrer; Grundloven
gaaer ud fra, at Vælgerfolket, der staaer bagved, om
det i Grunden altid vil det Rette, dog ikke altid har
den fulde, klare, tilstrækkelige Indsigt i, hvad der i
enhver Sag er det Rette, og desaarsag, naar Sindene
sættes i Røre, og Partiinteresser ere med i Spillet,
kan lade sig henrive til at ivre for Noget, som nøiere
beseet, er til Skade for Land og Rige. Grundloven
gaaer ud fra, at et opvakt Fleertal, just fordi det er
levende, driftigt og opvakt, kan forledes til at gjøre
Skridt, som Folket bag efter, naar Følgerne have viist
sig, bitterlig maa fortryde.

I Betragtning heraf har Grundloven til Modvægt
mod det rørige, fremaddrivende Element i Folkethinget
ved en egen middelbar Valgmaade oprettet et andet
Thing, Landsthinget, hvis Betydning er: at afgive det
mere paaholdende, controllerende Element og ved sit
Fleertal træde mæglende op mellem Regjeringsvillien
og Fleertalsvillien.

% --- p. 8 ---
\setcounter{footnote}{0}%
\opage{8}Mellem begge Thing staaer den kongelige Regjering,
virkeliggjort i et ansvarligt Ministerium.

Her er da i dette Regjeringshele tre Midtpunkter,
tre Myndigheder, tre Villier; men hvor er den herskende
Eenhed? Hvilken af de tre Myndigheder har
her den oprindelige Overvægt? Hvor ligger det egenlige
Tyngdepunkt? Det er paa Forhaand ikke let at
sige; thi hver af de tre har sin egen Selvgyldighed.
At Folkethinget ingenlunde tør betragtes som underordnet
Led, fremlyser tydelig af § 48, som siger:
„Forslag til Finantsloven og Tillægsbevillingslovene
behandles først i Folkethinget.“ --- Men at Landsthinget i Et og Alt er ligestillet med Folkethinget,
kan udtrykkelig sees af enhver Paragraph, hvori
Thingenes særlige Virksomhed, hvert for sig, og deres
gjensidige Forhold omtales (§§ 50---54). Kan Folkethinget
ændre Forslaget til Finantsloven, saa kan
Landsthinget, hvis det finder Grund dertil, sige Nei
til de foretagne Ændringer. I modsat Fald vilde det
jo være Folkethinget alene, der havde Skattebevillingsret.
Landsthinget maatte da være at betragte som
en Levning af „de raadgivende Stænder;“ men at
Landsthinget, hvad det Finantsielle angaaer, lige saa
fuldt er lovgivende som Folkethinget, er tydeligt.
Dersom Landsthinget, hvad Finantsloven angaaer, blot
havde en raadgivende, ikke en lovgivende Indflydelse,
saa maatte det ogsaa i Henseende til de øvrige Regjeringshandlinger
staae et heelt Trin lavere end
Folkethinget; dets Myndighed vilde da hermed ophøre
at være virkelig lovgivende Myndighed. Men at denne
Opfattelse staaer i aabenbar Strid med Grundlovens
Aand og Bogstav, falder i Øinene, naar vi sammen%
% --- p. 9 ---
\setcounter{footnote}{0}%
\opage{9}ligne § 2 „Den lovgivende Magt er hos Kongen og
Rigsdagen i Forening“, med § 29: „Rigsdagen bestaaer
af Folkethinget og Landsthinget.“

Tvedelingen er altsaa, naar vi see paa det hele
Regjeringssystem, blevet til en Tredeling: Folkethinget,
Landsthinget og mellem begge Kongens ansvarlige
Ministerium. Om denne Tredeling nu i Et og Alt er
hensigtssvarende, eller om der maaskee heri for Eenhedens
og Enighedens Skyld burde skee en Forandring,
er et Spørgsmaal for sig. Grundloven, een
Gang gjennemseet, kan atter gjennemsees; den udgiver
sig jo ikke for ufeilbar. Det hedder (§ 95): „Forslag
til Forandring i eller Tillæg til nærværende Grundlov
kunne fremsættes saa vel paa ordenlig som paa overordenlig
Rigsdag.“ Men --- dette staaer fast --- saa
længe Grundloven er, maa den gjælde ubetinget; den
er for Tiden vor høieste Maalestok for, hvad der er
lovligt og hvad der er ulovligt; vi have ingen anden.
At underlægge Grundloven en anden Mening end den,
som vitterlig svarer til dens Aand og Bogstav, er ufolkeligt,
thi det er uærligt. At benytte sig af Ordklang
og give det Udseende af, at Folkethinget i Et og Alt
repræsenterer den hele fuldstændige Folkevillie, er
ufolkeligt, thi det er uærligt. At rive en enkelt Sætning
ud af dens Sammenhæng og fordreie den til en
saadan Mening, at Landsthinget, hvad Finantslovens
endelige Vedtagelse angaaer, ikke skulde have lige
Myndighed med Folkethinget, er ufolkeligt, thi det er
uærligt. Og Eet endnu: at ville stille det ansvarlige
Ministerium i en saa eensidig Afhængighed af et
stemmegivende Fleertal, at den kongelige Villie svin%
% --- p. 10 ---
\setcounter{footnote}{0}%
\opage{10}der ind til en tom Form, et intetsigende Symbol, er
ogsaa, som vi nu skulle see, ufolkeligt og uærligt.

Hvad lærer Grundloven os om Kongens personlige
Stilling i Regjeringssystemet og om den kongelige Villies
Magt og Betydning? „Kongen er ansvarsfri, hans Person
er hellig og ukrænkelig. Ministrene ere ansvarlige
for Regjeringens Førelse“ (§ 12). Naar man nu heraf
udleder, at man ikke paa upassende Maade bør blande
Kongens Navn ind i Partiernes Strid, og at Kongen
kun igjennem ansvarlige Ministre udøver sine Regjeringshandlinger
o. s. v., da er dette jo ganske i sin
Orden. Men gaaer man videre og udleder heraf, at
Kongens personlige Villie i Virkelighed ingen Indflydelse
har paa Styrelsens Gang, da er dette saa
stridende mod Grundlovens Aand og Bogstav, som
det vel er muligt.

Ved personlig Villie maae vi her, vel at mærke,
ingenlunde tænke paa private Ønsker eller individuelle
Tilbøieligheder. En i Sandhed personlig Villie gaaer
altid ud paa noget Væsenligt, i personlig Forstand Almeengyldigt.
Den kongelige Villie gaaer personlig ud
paa at overholde, hævde og gjennemføre Danmarks
Riges gjennemseete Grundlov med alle deraf flydende
Bestemmelser: dette vil Kongen personlig; dette og
kun dette kan han som Konge personlig ville. „Forinden
Kongen tiltræder Regjeringen, afgiver han skriftlig
i Statsraadet den edelige Forsikkring, ubrødelig at
holde Rigets Grundlov“ (§ 7). Under denne Forudsætning
bliver Kongen personlig høieste Fuldbyrder af
alle Regjeringshandlinger: uden den kongelige Villie er
Loven død. Ministrenes Ansvarlighed er desaarsag
dobbelsidig: de ere ansvarlige for Rigsdagen, og de
% --- p. 11 ---
\setcounter{footnote}{0}%
\opage{11}ere ansvarlige for Kongen. Hvad siger Grundloven?
„Kongen udnævner og afskediger sine Ministre. Han
bestemmer deres Antal og Forretningernes Fordeling
imellem dem.“ (§ 13). Mener et Ministerium ikke at
kunne forsvare for Rigsdagen, at det bliver paa sin
Plads, da er det unegtelig beføiet til at træde af.
Men Kongen har, ifølge Grundloven, jo ogsaa fra sin
Side Ret til at afskedige Ministre, der ikke træffe,
hvad han til Rigets Bedste anseer for det Hensigtssvarende,
og vælge sig et Ministerium, hvor et saadant
kan findes, som ved alle lovlige Midler vil arbeide
paa at gjennemføre Hans Majestæts personlige
Opfattelse af den rette, lovlige Gang i Styrelsen. ---

Hvad lærer Grundloven om Kongens Forhold til
Rigsdagen? „Uden Kongens Samtykke kan Rigsdagen
ikke forblive længere sammen end to Maaneder“. (§ 13).
„Kongen kan sammenkalde Rigsdagen til overordenlige
Sammenkomster“. (§ 20). „Kongen kan opløse
enten hele Rigsdagen eller een af dens Afdelinger“.
(§ 22). „Har Kongen ikke stadfæstet et af Rigsdagen
vedtaget Lovforslag inden næste Rigsdags Samling,
ansees det som bortfaldet“. (§ 24). Det fremlyser da
heraf, at den kongelige Villies Functioner ere mere
end symbolske, at det er Regjeringshandlinger af gjennemgribende
Betydning. Grundloven har ved at indrømme
den kongelige Villie en saa væsenlig overordnet
Myndighed villet sikkre sig en den hele Regjering
besjælende og styrende Eenhed. Det er fra dette
Synspunkt af Vigtighed at gjøre sig Rede for, hvorledes
det constitutionelle Monarchie i Principet adskiller
sig paa den ene Side fra Enevoldsmagten (det
absolute Monarchie), paa den anden Side fra Fleer%
% --- p. 12 ---
\setcounter{footnote}{0}%
\opage{12}talsmagten (Demokratiet) --- særlig med Hensyn paa
Folkevillien.

I det absolute Monarchie gaaer Folkevillien under
i Kongevillien, thi Kongen staaer over Loven: den enevældige
Hersker magter vilkaarlig det Hele. I Demokratiet
gaaer Folkevillien til Grunde i Fleertalsvillien:
det er ikke Folkets sunde Fornuft, for saa vidt den
forfatningsmæssig arbeider sig frem, som her er Hovedsagen;
nei, det er Befolkningens i de fleste Stemmer
udtalte Stemning, som ophøies til øverste, styrende
Magt og Myndighed. I det constitutionelle Monarchie
staaer den kongelige Villie fast paa Lovens Grund.
Folkevillien faaer fra denne Side sit Udtryk i Form af
Eenhedsvillie, udeelbar Enkeltvillie. Denne Enkeltvillie
er ikke vilkaarlig, tvertimod! den er i allerhøieste Forstand
folkelig. Den kongelige Villie vil, hvad Folket
vil, fordi den vil alene, hvad Loven vil; den kan ikke
ville Andet, thi i modsat Fald vilde den jo, ifølge Forfatningen,
ophøre at være kongelig. I det constitutionelle
Monarchie faaer Folkevillien tillige sit rette, begrænsede
Udtryk i Fleertalsvilliens Form: det er den deelbare,
ved Delingen repræsenterede Folkevillie, som hos os
kommer til virkelig Udtalelse paa Rigsdagen. Mellemledet
mellem Kongevillien, den udeelbare Enkeltvillie,
og den deelbare, i de to Thing deelviis repræsenterede
Fleertalsvillie er Ministeriet med sin dobbelte Ansvarlighed.

Saa længe nu de af Ministeriet fremlagte Forslag
finde en saadan Tilslutning i begge Thing, at Enighed
kan tilveiebringes, saa længe vedbliver Flertalsvillien
at være i Overeensstemmelse med sig selv, og den herved
tilfredsstillede Kongevillie finder for saa vidt ingen
Anledning til særlig Indgriben. Ministeriets Ansvar%
% --- p. 13 ---
\setcounter{footnote}{0}%
\opage{13}lighed for Rigsdagen er i fuld Harmonie med dets Ansvarlighed
for Kongen. Men i Tilfælde af, at Rigsdagen
bliver splidagtig med sig selv, faaer Sagen et andet
Udseende. Vinder et Forslag Bifald i Landsthinget,
men forkastes af Folkethinget, kan det ikke sættes
igjennem, thi Folkethinget har en ubestridelig Ret til
at gjøre Modstand. Vinder et Forslag Folkethingets
Bifald, men forkastes af Landsthinget, kan det heller
ikke sættes igjennem, thi ogsaa Landsthinget har en
ubestridelig Ret til Modstand. Uenighed mellem Thingene
indbyrdes er en Uenighed i den ved sine Repræsentanter
deelbare Folkevillie; det er Folkevillien, som
under Form af Fleertalsvillie er blevet splidagtig med
sig selv.

Men en med sig selv splidagtig Villie kan ikke
handle. Det er forgjæves, at det ene Thing vil beraabe
sig paa, at det repræsenterer et større Fleertal end det
andet, at det har hele Vælgerfolket i Ryggen o. desl.;
slige Stemningsforklaringer have ingen Hjemmel i Grundloven.
Grundloven tager alene Hensyn til det Forhold,
hvori de to Thing, hvert med sine lovlig valgte Medlemmer,
stille sig til hinanden indbyrdes. I Tilfælde
af Strid bliver det altsaa nødvendigt, at Folkevillien
under den anden Form, Eenhedsvillien, den kongelige
Villie, griber ind og træffer Foranstaltninger til, at
Striden kan udlignes.

For at tilveiebringe Enighed, kan Kongevillien,
som bekjendt, vælge en af to Veie: enten lade Ministeriet
træde af eller opløse Rigsdagen. Vælges det
første Middel, da maa dets egenlige Betydning jo aabenbart
ligge deri, at der virkelig skeer et Systemskifte.
Et saadant Systemskifte som, forstaaer sig, maa være
% --- p. 14 ---
\setcounter{footnote}{0}%
\opage{14}i fuld Overeensstemmelse med Grundlovens Aand og Bogstav,
kan for Øvrigt helde mere til Stemningen i det ene
Thing end i det andet; men det bør dog, saa sandt
det skal være folkeligt, være hævet over eensidige Partihensyn.
Vælges det andet Middel, har H. M. Kongen den
fast begrundede Overbeviisning, at et Systemskifte vilde
være til Skade for Land og Rige, saa har den kongelige
Villie, paa Folkets Vegne, jo Magt til at opløse enten
hele Rigsdagen eller een af dens Afdelinger.

Men --- spørger man --- hvis dette sidste Middel
nu heller ikke hjælper, hvis det ved gjentagne Omvalg
viser sig; at det opponerende Parti i det ene Thing
bliver større og større, hvad saa? Vil det ikke heraf
blive en nødvendig Følge, at der maa foregaae en Systemforandring
til Fordeel for Fleertallet og dets
Vælgere? Vi svare: Jo! det vil ganske vist blive nødvendigt.
Hvis den hele Forfatning ikke skal sprænges,
maa det omsider komme dertil, at Vælgerfolket faaer
sin Villie respecteret. Vilde Enevillien forsøge paa at
tage Fleertalsvillien med Magt, saa var Grundloven
brudt: vi havde da ikke længere et constitutionelt, men
et absolut Monarchie. Men --- lad os ikke springe
Mellemleddene over. Den udtrykkelige Forskjel, Grundloven
gjør mellem at opløse „den hele Rigsdag“ og at
opløse „een af dens Afdelinger“, maa dog vel have
Noget at betyde. Betydningen maa være denne: naar
een af dens Afdelinger opløses, saa er det kun Medlemmerne
af denne Afdeling, der underkastes Omvalg;
opløses derimod hele Rigsdagen, saa maae alle Medlemmerne
(Landsthingets lige saa vel som Folkethingets)
vælges paa Ny. At gjentage Opløsningen af blot den
ene Afdeling (Folkethinget f. Ex.) og stille provisoriske
% --- p. 15 ---
\setcounter{footnote}{0}%
\opage{15}Foranstaltninger i Udsigt, for at frembringe Enighed,
vilde jo være en partisk, altsaa ufolkelig Fremgangsmaade.
Forandre nu Valgene sig --- vel at mærke
efter gjentagne Opløsninger af hele Rigsdagen --- saaledes,
at Landsthinget omsider bliver enigt med Folkethinget,
da er dette et uimodsigeligt Beviis for, at
Folkethinget i det Hele har Vælgerfolkets Stemning
for sig. Den kongelige Eenhedsvillie har da fra sin
Side hjulpet Fleertalsvillien til lovbestemt Enighed med
sig selv og derved givet den dens tabte Handlekraft
tilbage. Men dersom det Modsatte bliver Tilfældet,
dersom det efter gjentagne Omvalg viser sig, at et
Fleertal i Landsthinget --- ligegyldigt, om det er større
eller mindre --- bestandig vedbliver at staae stivt imod
et Fleertal i Folkethinget, da har Folkethingets Førere
aldeles ingen Ret til at beraabe sig paa Vælgerfolkets
Samstemning i det Hele; thi Vælgerne til Landsthinget
--- H. M. Kongen ikke at forglemme --- tilhøre jo, i
Følge Grundlovens utvetydige Bestemmelse, Vælgerfolket,
lige saa fuldt som Vælgerne til Folkethinget.

At Modstand, Modstand og Modstand ikke, naar
det da skal gaae lovlig til, formaaer at fremtvinge
Overeenskomst, er indlysende ved sig selv. Vi kunne,
uden ubeføiet Indblanding i den praktiske Politik, kaste
et Blik paa Partiernes Stilling. Fra den ene Side høre
vi Erklæringer, der lyde omtrent saaledes: „Det maa
omsider blive vitterligt, at hvad her strides om, er mere
end dette eller hiint. Vi kæmpe imod et Regjeringssystem,
der ikke længere passer til Tiden. Personerne
vexle, men det udslidte System kommer stadig igjen:
lidt Ændring i Miner og Holdning, men altid det gamle
Ansigt! Hvad vi ville, hvad vi stride for, er noget
% --- p. 16 ---
\setcounter{footnote}{0}%
\opage{16}Nyt og eiendommelig Frisk, der arbeider sig frem i
Folket. Her er ikke Strid om nogle Millioner alene,
ikke om Skatternes Fordeling alene, ikke om disse eller
hine Forsvarsmidler alene: Striden er om Retning mod
Retning, om System mod System, om Aand imod Aand.
Hvad Luther sagde til sine Modstandere ved Forhandlingerne
i Marburg, det sige vi Mænd af det menige
Folk med god Samvittighed til vore Modstandere paa
Rigsdagen: „I have en anden Aand end vi.“ --- Fra den
anden Side lyder Forklaringen saaledes: „Her er slet
ikke Strid om System mod System, men Spørgsmaalet
er om System mod Parti. Spørgsmaalet er, om et
Fleertalsparti, af misforstaaet Iver for hvad det selv
stempler som folkeligt, skal have Ret til at slaae
Partiaanden løs, rive Stemningen med sig, kaste Vrag
paa Sagkyndiges Dom og appellere til Mængden. Her
gjøres Modstand alene for Modstandens Skyld. Ved at
slaae de mest uligeartede Sager i Bunke og filtre Betingelserne
sammen i en Haardknude gjør man enhver
Regjering umulig. Er det til Landets Bedste, er det
folkeligt, er det fornuftigt, er det Ret?“

Alt som Knuden strammes og Spændingen voxer,
nærmer man sig det kritiske Vendepunkt, da provisoriske
Foranstaltninger kunne blive nødvendige, navnlig
en provisorisk Finantslov. Man har freidig afviist Muligheden
af en saadan Finantslov og trøstet sig med Grundlovens
Ord (§ 49): „Forinden Finantsloven er vedtagen,
maa Skatterne ei opkræves. Ingen Udgift maa afholdes,
som ikke har Hjemmel i Finantsloven eller i en Tillægsbevillingslov.“
Er Udstedelsen af en provisorisk Finantslov
da ligefrem lovlig? Vi svare: Nei! Men at beraabe sig paa
det Lovlige, naar man selv har været med til at gjøre det
% --- p. 17 ---
\setcounter{footnote}{0}%
\opage{17}Lovlige umuligt, er en Selvmodsigelse, der ingen Trøst
indeholder. Den er heller ikke ligefrem ulovlig. „I
særdeles paatrængende Tilfælde kan Kongen, naar
Rigsdagen ikke er samlet, udstede foreløbige Love.“
(§ 25). Hele Eftertrykket falder paa det „særdeles
paatrængende.“ Men enten en Sammenkaldelse af
Rigsdagen bliver umulig, fordi Fjenden er i Landet,
eller den sammemkaldte Rigsdag paa Grund af Partiernes % print: „sammemkaldte“ (misprint for sammenkaldte)
indbyrdes Modstand gjør sig selv umulig, er
Tilfældet jo lige paatrængende. Istedenfor at trøste
sig med det Forslidte: det har ingen Fare, var det dog
vist bedre at see lidt nærmere paa § 95, at den Vei,
ad hvilken Fleertalsvillien med fuld Bevidsthed om sin
lovbegrændsede Ret skrider frem til fuld Overeensstemmelse
med den kongelige Eenhedsvillie, kan komme
til at ligge aaben og klar for den hele Befolkning.

Det er vor Forfatnings Grunde vi ere i Begreb
med at uddybe. Der er Banker og Skjær i disse Farvande:
vi maae lodde, sætte Sømærker ud og sikkre
Seiladsen. Nogle Vindstød forskrække os ikke; om
der ogsaa kommer en Kuling med Braadsøer, vort Skib
vil dog ikke strande: agter om Stormast staaer sund
Fornuft, og Besindighed sidder ved Roret. Endnu er
Farten god, og Faren kun en Sky i Horizonten; endnu
er intet Ulovligt skeet; endnu staaer Vidnesbyrdet fast,
Historiens hæderlige Vidnesbyrd, at det danske Folk,
end i Spændingens Øieblik, er et roligt, sindigt, alsidig
betænksomt Folk.

% --- p. 18 ---
\setcounter{footnote}{0}%
\lecture{II.} {Folkelig Selvstyrelse: Folk og Konge.}
\opage{18}En kraftig, sund og alsidig Udvikling af Folkelivet er
betinget af folkelig Selvstyrelse. Naar et Folk er naaet
saa vidt, at det har hævet sig op paa Frihedens Trin
og er blevet myndigt, kan det ikke længere taale nogen
udvortes Herre eller Formynder over sig. Det
myndige Folk vil have sin Villie; det vil raade sig
selv, styre sig selv og tage Følgerne af, hvad det selv
har besluttet.

Denne Grundsætning har for Mænd af den gamle
Skole noget høist Anstødeligt, næsten Gruopvækkende.
„Hvad skal det blive til“, siger man, „naar et Folk bliver
saa selvraadigt, at det ikke vil lade sig regjere:
Frækhed, Uvorenhed, Tøilesløshed maa jo tage Overhaand
Dag for Dag.“ --- Man fatter let, at der her
skjuler sig en Misforstaaelse. Det ligger i Begrebet
Styrelse, at der maa være Nogle, der styre, Nogle, der
styres, og en Lov, hvorefter der styres; men det ligger
ingenlunde i Begrebet Styrelse, at det Styrende og det,
der styres, maae holdes adskilte ved en uoverstigelig
Skranke. Dersom det Sidste var ubetinget nødvendigt,
da vilde al Selvstyrelse være umulig: intet Væsen, hverken
i Himlen eller paa Jorden, vilde da være i Stand
% --- p. 19 ---
\setcounter{footnote}{0}%
\opage{19}til at styre sig selv. Det er en overfladisk Forestilling,
at de Styrende ere frie og kunne gjøre hvad de
ville, medens de, der maa lade sig styre, aldrig faae
deres Villie, men bestandig ere under Tvang. Regjeringen
--- siger man --- er fri, Folket er ufrit; Øvrigheden
er fri, Undersaatterne ufrie; Formynderen er
fri, Myndlingen ufri. Dersom dette Forhold stod uforanderlig
fast, saa kunde der unegtelig aldrig være Tale
om et frit Folk; men saa kunde en Myndling jo heller
aldrig blive myndig. Skal Myndlingen blive myndig,
saa maa Formynderen træde af. Naar kan Formynderen
træde af? Vi svare: naar det Styrende, som
hidtil var i ham, altsaa udenfor Myndlingen, lægges
ind i Myndlingen selv. Denne har da ikke længere
Formynderskabet udenfor sig, men inden i sig: Formynderskabet
hører ikke op, men det forflyttes. Netop
ved at have optaget Formynderen i sit Indre er Myndlingen
blevet myndig. Han er i samme Grad myndig,
som han, aandelig forstaaet, har spiist, fortæret og fordøiet
sin Formynder.

Det Samme gjælder om et heelt Folk. Det begynder
i Umyndighed med at have Styrelsen, Regjeringen,
udenfor sig som et Fremmed, et udvortes Overordnet.
Overgangen til Selvstyrelse skeer ikke ved, at
et Folk kaster Regjeringen af, omtrent som en Hest
slaaer Rytteren af; nei derved, at det inderlig optager
det styrende Element i sig selv og forlægger Regjeringens
Sæde til sit eget Midtpunkt. Formynderskabet,
den udvortes, udenfor Folket staaende Myndighed, hører
op, naar Folket er blevet myndigt, naar det, med
andre Ord, har spiist, fortæret og fordøiet sin Formynder.


% --- p. 20 ---
\setcounter{footnote}{0}%
\opage{20}Den væsenlige Lighed mellem Selvstyrelsen hos det
enkelte Individ og hos Folket har imidlertid en virkelig
Ulighed, vi ikke maae oversee. Naar den Enkelte
skal til at styre sig selv, erfarer han strax, at der i
hans Natur er et lavere Element, som skal bringes til
Lydighed under et høiere. Der er sandselige Tilbøieligheder,
Drifttilskyndelser og Lidenskaber, som skulle
beherskes; kun den af Indsigt ledede Fornuftvillie, i
endnu høiere Forstand, Aandsvillien bør herske. Men
et Menneskes Villie bevæges ikke alene af Hensyn til
det Fornuftige; der er en Naturvillie i Kjød og Blod,
en Driftvillie, der grener sig i mange Tilbøieligheder:
Hovedet vil en Vei, Hjertet en anden, Egenkjærligheden
en Vei, Pligtfølelsen en anden. Idet Forstanden tager
Parti snart med den ene, snart med den anden af disse
Tilbøieligheder, opstaaer der Partikampe, ofte stormende
Valgkampe i Menneskets Indre; det skal da vise sig,
hvor vidt der er en herskende Villieseenhed eller ikke.
Uden en høiere Villieseenhed kan ingen Charakteer
danne sig. Individet med de mange Villier bliver
splidagtigt med sig selv, rives hen snart af en Stemning,
snart af en anden, svimler, svingler og raver og
ender med at løbe vild i sig selv. Skal der danne sig
en Charakteer, da maa der være en overordnet, styrende
Villie, som magter det Splidagtige, en kongelig
Villie, som er i Stand til at forlige de indbyrdes kivagtige
Smaavillier. Der er i enhver gjennemarbeidet
Charakteer en constitutionel-monarchisk Forfatning.
Den styrende Enevillie kan ikke ved et vilkaarligt
Magtsprog bringe nogen som helst af Driftvilliens
Yttringer til Taushed, heller ikke bringer den Orden
og Fred ved at lade de deelte Villier afslibe hver%
% --- p. 21 ---
\setcounter{footnote}{0}%
\opage{21}andre i fortsat Strid, og holde sig bag de villende Tilbøieligheders
Fleertal som et tomt Symbol, et Lovens
Skyggebillede. Charakteervillien er den levende Lov;
den gaaer ind paa den sjælelig-legemlige Constitutions
Vilkaar, den gjennemtrænger, renser, begrændser og
forliger de indbyrdes stridige Tilskyndelser: jo mere
constitutionel den herskende Eenhedsvillie er, desto
mere kongelig er den. Alt dette skeer paa indvortes
Maade, og kun naar dette virkelig skeer, har Mennesket
sig selv i sin Magt, og er selvstyrende.

Anvende vi dette paa den folkelige Selvstyrelse, da
er Væsenet ganske vist det samme; men dets Fremtræden
i sit Udvortes er en anden. Her er i Samfundet
ikke blot en Mangfoldighed af Villiesretninger; men
her er lige saa mange Villier som Individer. Folket er
i den hele Befolkning; Folkets Villie, seet fra denne
Side, er splittet og deelt: Mængdens Villie er en Samling,
en Sum af indbyrdes forskjellige Villier, en Villiesvrimmel,
som, for at danne en sand Folkevillie, maa
gaae sammen til Eenhed. Idet nu den styrende Villie,
for at tilveiebringe Overeensstemmelse og Eenhed, træder
udvortes op mod de mange selvraadige Villier,
fremkommer der en uheldsvanger ag faretruende Spænding % print: „ag“ (misprint for og)
mellem Eenhedsvillien og Mængdevillien. Saa
længe Mængdevillien, nærmere bestemt som Fleertalsvillie,
ved Paaberaabelse paa Fleertalsmagten stiller sig
skarpt imod en hvilken som helst herskende Villies
Enemagt, er Samfundet splidagtigt med sig selv, bestandig
svævende mellem Tyranni og Anarchie. Under
saadanne Omstændigheder kan der naturligviis ikke
være Tale om fornuftig Styrelse, endnu mindre om
folkelig Selvstyrelse. For Styrelsens Skyld maa der
% --- p. 22 ---
\setcounter{footnote}{0}%
\opage{22}tilveiebringes en lovbestemt, fornuftig Vexelvirkning
mellem Eenhedsvillie og Mængdevillie, der maa, med
andre Ord, gives en Constitution.

Hvad det har kostet at faae en nogenledes fornuftig
Constitution udarbeidet, lærer Historien. Tidlig har
man følt Nødvendigheden af en styrende Eenhedsvillie.
Allerede Homer synger: „Fleertal som Herskab er ikke
godt; kun Een være Hersker.“ Men tidlig har man
ogsaa maattet erfare Følgerne af Enevældens umenneskelige,
eller rettere, kun altfor menneskelige Overgreb:
Folkene have sukket under et tungt Aag, Mængdevillien
er efter vilde Oprør bestandig sunket tilbage i Afmagt.


I Middelalderen var Samfundsconstitutionen aristokratisk.
Kun i de store Adelsslægter rørte sig en Fleertalsvillie
til Fordeel for Friheden, vel at mærke: de
Adeliges Frihed; om Frihed for Menigmand var der
naturligviis ikke Tale. Med hvilken Iver Adelen stræbte
at gjøre Kongemagten til et Skin og forvandle den
kongelige Villie til et underskrivende Symbol paa Rigsraadets
Myndighed, derom kan Haandfæstningsperioden
i vort Fædrelands Historie --- fra Christopher II. til
Frederik III. --- afgive tilstrækkeligt Vidnesbyrd.

Ved at række Borgerstanden Haanden og give sig
Udseende af at ville ophjælpe det hele Folk, fik den
bundne Kongemagt Lænkerne rystede af sig; men hvad
skete nu? Enevælden kuede Adelen og lærte den at
bøie sig for Thronen. Men Bøiningen gjorde kun Stoltheden
smidig, Adelen holdt sig i Veiret som Hofadel,
og det menige Folk sank endnu dybere. Almeenvillien
gik op i Staten, og Staten gik op i den kongelige Villie.
Det var ikke forfængeligt Praleri, men virkelig Sandhed,
% --- p. 23 ---
\setcounter{footnote}{0}%
\opage{23}naar Ludvig XIV. sagde om sig selv: Jeg er Staten
(l'état c'est moi). Folket havde nu faaet een Formynder
isteden for mange; men Formynderen stod udenfor
Folket, han var med guddommelig Myndighed ophøiet
over Folket, og lod det føle, at det i Et og Alt var
umyndigt. Det gik indtil Revolutionen. Med Revolutionen
(1789) kom et Element til Gjennembrud, som i
vore Dage spiller en vigtig Rolle, nemlig det demokratiske.
For at nedbryde Enevælden blev det nødvendigt,
at Demokratiet væltede Thronen, jevnede Grunden
og banede Veien for videre Fremskridt. Saa fik Demokratiet
Strømmen med sig. Udtryk som: Mængdevillie,
Fleertalsvillie, have i den offenlige Mening faaet
en saadan Indgang, at de --- til ubodelig Skade for den
folkelige Sag --- uden videre antages for eensbetydende
med Folkevillien. Det er en Skævhed i Synsmaade,
som efterhaanden maa rettes. Demokratiets Styrke er,
at gjøre Fleertalsvillien gjældende, deri har det sin Ret
og Betydning. Men i sin Styrke har det ogsaa sin
Svaghed. Byggende paa Fleertalsvilliens Grundlag, maa
Demokratiet uafbrudt kæmpe for et Uopnaaeligt, thi
det kan aldrig naae til den herskende Villies Eenhed;
Fleertallet skifter, Parti staaer op mod Parti.
Den Eenhed, som søges i det Heles lovbestemte Sammenhæng,
er svag og vaklende: Sammenhængseenhed
uden Eenhedsvillie, uden gjennemgaaende Selveenhed,
mangler tilstrækkelig Grund. Hvilende paa demokratisk
Underlag er Statsbygningen løs i Fugerne, idelig
udsat for Orkanstød, naar Fleertalskræfterne røre sig i
Partielementernes Uveir. Man tænke paa Amerikas demokratiske
Partikampe og de Valgstorme, som rase ved
hvert Præsidentskifte; man tænke paa Frankrigs for%
% --- p. 24 ---
\setcounter{footnote}{0}%
\opage{24}viklede Tilstand; gamle og nye Dynastier intriguere
mod hverandre, medens den communistiske Republik,
„det røde Spøgelse“, truer med Demokratiets conseqvente
Gjennemførelse. Det demokratiske Element har
været, er endnu og vil altid være en uundværlig Factor
i den frie Samfundsudvikling; men Demokratiet er intet
fuldstændigt Princip. Ophøies Demokratiet til ene
gjældende Princip, da bliver det i sit Væsen lige saa
eensidigt, for Folkelivets Udvikling lige saa fordærveligt,
som det absolute Monarchie. Vi skulle, for at
bane os en Vei til fornuftig Indsigt i folkelig Selvstyrelse,
forsøge at belyse det nærmere.

Der er vistnok i vor Tid Mange, som have vanskeligt
ved at forlige sig med den Tanke, at det constitutionelle
Monarchie skulde betegne et heelt, i sig selv
sammenhængende Regjeringssystem; det maa tvertimod
forekomme dem, som var det en Art Flikværk, et efter
Tidsforholdene praktisk aftvunget Compromis, en Sammenstykning
af to indbyrdes uforligelige, men til gjensidig
Afpasning beklippede Systemer, nemlig det absolute
Monarchie og det rene Demokratie. Man gaaer herved
ud fra den aldeles feilagtige Forudsætning, at Enevælden
og Demokratiet, hver for sig, ere hele Principer. Det er
en fra Misforstaaelse af Historien hentet Fordom. Historisk
betragtet, gjør ethvert af de to Systemer unegtelig
Indtryk af at være et fuldstændigt Hele; men de historiske
Omvæltninger gjøre jo tillige Indtryk af, at Samfundet
lige saa lidt har kunnet slaae sig til Ro i det
ene System som i det andet. Der er i Europas nyere
Historie Intet, der i den Grad har vanskeliggjort en
fornuftig Brug af det constitutionelle Monarchie og
hæmmet Folkelivets Udvikling, som den falske Fore%
% --- p. 25 ---
\setcounter{footnote}{0}%
\opage{25}stilling, at dette Monarchie skulde være en Blanding
af uforligelige Bestanddele. Allerede Sammensmeltningen
af Benævnelser som: det „indskrænkede“ og det
„constitutionelle“ Monarchie bidrager Sit til at nære
Misforstaaelsen.

Da Revolutionen var endt med Napoleons Fald,
dukkede Forestillingen om den gamle patriarchalske
Enevælde, Kongedømmet af Guds Naade, frem med
fornyet Styrke. Ruslands, Østerrigs og Preussens Eneherskere
sluttede --- uden Medunderskrift af deres Ministre ---
i Paris den \emph{saakaldte hellige Alliance}
(26 Septbr. 1815). I dette Dokument hedder det, at
de tre contraherende Magthavere, i Overeensstemmelse
med Skriftens Ord, ville betragte hverandre som Brødre,
hjælpe hverandre i mødende Tilfælde og lede deres
Undersaatter og Hære, til hvis Familiefædre Forsynet
havde sat dem, i den samme Broderlighedens Aand;
de erkjende at „den christelige Nation, til hvilken de
og deres Folk høre, i Virkeligheden ikke har anden
Souverain end ham, hvem Magten alene tilkommer,
nemlig Gud og vor guddommelige Forløser, Jesus
Christus, den Allerhøiestes Ord, Livets Ord.“ For en
trohjertet Fromheds lyseblaa Blik maatte en saa
christelig-faderlig Optræden af Jordens Mægtige unegtelig
tage sig ud som et Tegn fra Himlen, en Forjættelse
om Fred og Lyksalighed for Europas Folk
efter lange og haarde Lidelser. Det kunde jo synes,
som var den selvherskende Enevillie ved en forhøiet
Temperatur af Religiøsitet og Kjærlighed omsider blevet
saa gjennemvarmet og udvidet, at den nu virkelig formaaede
at stille Fleertalsvillien til Ro og nære Folkefriheden
ved sin Barm. Men Statsmændene smilte; de
% --- p. 26 ---
\setcounter{footnote}{0}%
\opage{26}som kjendte lidt til Verdens Gang og mindedes, at
Traktaten om Polens Deling netop begynder: „\emph{i den
hellige Treenigheds Navn}“, følte sig betagne af
uhyggelige Anelser og skimtede allerede de Kløer, paa
hvilke man kjender Løven.

Skulde Folkefriheden ikke gaae til Grunde, da var
det paa Tide, at Demokraterne rykkede i Marken med
Fleertalsvillien. Det var en Optræden, der krævede
Udholdenhed og Mod. Ethvert Forsøg paa enten at
opnaae eller endog bevare saadanne Rettigheder (Valgret,
Skattebevillingsret, Pressefrihed, Forsamlingsfrihed o.s.v.),
som ere nødvendige, dersom et Folk skal kunne røre
sig frit, ethvert Forsøg i frisindet Retning blev af de
absolutistiske Regjeringer betragtet som et frækt Angreb
paa det Allerhelligste. At afsløre Misbrug var
ugudelig Ondskab, at henvende sig til Mængden var en
politisk Forbrydelse. Absolutismens fanatiske Tilhængere
saae i enhver friere Rørelse et nyt Hoved skyde
frem paa Revolutionens Hydra, og ilede med at hugge
det af. Her er et Vendepunkt, hvor demokratisk Borgersind,
Mod og Opoffrelse viser sig i sit fordeelagtigste
Lys. Revolutionen havde --- til Lykke for Humanitetens
Sag --- givet Demokratismen et saa mægtigt Opsving,
at den frie Aand ikke mere lod sig kue. I Perioden
fra 1815 til 1830 er det Demokratiets Førere, der træde
i Forgrunden og maae drage Læsset. Deres Uforfærdethed,
Iver og Udholdenhed har --- trods alle
demagogiske Udskeielser --- gyldigt Krav paa Efterverdenens
Tak og Paaskjønnelse.

Man vil heraf kunne see, hvorledes det er gaaet
til, at de to Halvprinciper, det monarchiske med Eenhedsvillien
og det demokratiske med Fleertalsvillien, have
% --- p. 27 ---
\setcounter{footnote}{0}%
\opage{27}udgivet sig, hvert især for et heelt, fuldstændigt Princip
og ere optraadte ei alene som midlertidige Modstandere,
men som fødte Fjender. Det vil heraf kunne forstaaes,
at de Indrømmelser, saadanne to Magter under fortsatte
Brydninger kunne aftvinge hinanden, ingenlunde egne
sig til at danne et Regjeringssystem med sand organiserende
Eenhed. Det constitutionelle Monarchie, Monarchiet
omgivet og begrændset af demokratiske (republikanske)
Institutioner, maa, naar det sees fra denne Side,
nødvendig lide af de Onder, som giftige, med Organismen
uforenelige Stoffer foraarsage i et for øvrigt livskraftigt
Legeme. Hvilke Krampetrækninger heraf kunne opstaae,
er historisk anskueliggjort i Stridighederne under Ludvig
XVIII og Carl X om Anvendelsen og Fortolkningen af
det franske Charte.

Intet er fordærveligere end at bygge paa en
Grundlov, der baade skal gjælde og dog ikke gjælde.
Opfattet som en Sammenstykning af to fjendtlige Principer,
er det constitutionelle Monarchie netop egnet til
at vække, nære og vedligeholde en uafladelig Tvedragt
mellem Folkerepræsentation og Regjering. Naar begge
Parter i Grunden ere utilfredse med Forfatningen, kunne
de ved Modstand, Intriguer og Overgreb hver fra sit
Stade chicanere hinanden i det Uendelige.

Er det Kongevillien, der i absolutistisk Selviskhed
føler Trang til at udvide sin Magt, vil Overgrebet fra
oven aldrig savne Midler. Her hjælper det ikke at
stole paa Ministeransvarligheden. Hvad der ikke kan
skee med Vold kan skee med List; tør et Ministerium
ikke udsætte sig for en Rigsrets-Dom, nu vel, saa tør
et andet: en Carl X vil altid finde en Polignac. Der
er hos Regjeringsmænd, som ikke tage det nøie med
% --- p. 28 ---
\setcounter{footnote}{0}%
\opage{28}det Moralske i Midlerne, en egen Snuhed til at sondere
Stemninger. Man har i al Hemmelighed et Parti i
Baghaanden, man gjør maskerede Bevægelser, der forvirre
Modstanderne, og holder sig med stolt Foragt for
den offenlige Mening til det ene afgjørende Spørgsmaal:
\emph{hvor meget kan man byde et Folk, før det gjør
Oprør}? Vel sandt, Beregningen kan slaae feil, men
man maa vove: Forsøget kan jo ogsaa lykkes. Hvad
der mislykkedes for en Carl X. (1830), faldt jo ved
Statskoupet (1851) heldigt ud for en Napoleon III.

Er det omvendt en demokratisksindet Folkerepræsentation,
ja, selv om det kun er en Afdeling af samme,
der er misfornøiet med Forfatningen, da vil Overgrebet
fra neden heller ikke savne Midler til snedig Afbenyttelse.
Først og fornemmelig kommer det da an paa
at benytte Mængdens eller dog Fleertallets Nærsynethed
til at faae den paa Afstand holdte Grundlov stillet i
en saa skæv Belysning, at det livagtig seer ud, som
om den ubetinget talte for Demokratismens Anskuelse.
Løsrevne Sætninger, Slagord og Talemaader ere vel
skikkede til at agitere Befolkningen. Det menige
Folk troer oprigtig paa slige Forklaringer --- især naar
det ikke kjender og paa det Eftertrykkeligste advares
imod at gjøre Bekjendtskab med andre Forklaringer,
eftersom disse paa Forhaand vides at være et Væv af
Skin og Bedrag --- hildes da let i den Indbildning,
at dets Rettigheder krænkes, at de demokratiske Førere,
ingenlunde af Begjærlighed efter Magt, men i sand
Selvfornegtelse og med megen Opoffrelse stride for et
sandt, et heelt og fuldstændigt Princip til Alles Gavn
og Bedste. Dernæst --- naar Agitationen er i fuld
Gang --- gjælder det om at imponere Regjeringen og
% --- p. 29 ---
\setcounter{footnote}{0}%
\opage{29}styrte det ene Ministerium efter det andet. Det sidste
Middel er lige saa populært som det første. Historien
lærer, at det menige Folk, saa længe det endnu ikke
er blevet myndigt, gjerne fryder sig og klapper i Hænderne,
naar en Regjering falder. Der er en Række
Exempler paa, at det endogsaa er skeet i saadanne
Tilfælde, hvor den faldne Regjering var Folkesagens
egenlige Forsvarer, og det var Folkets Fjender, der ved
Omvæltningen kom til Roret. Hvis en Regjering under
saadanne Omstændigheder lader sig imponere af Fleertalsføreres
Fripostighed, kan det ganske vist ende med en
Kuldkastelse af Forfatningen. Hvis en Regjering derimod,
ved at holde sig urokkelig paa Lovens Grundlag,
forvandler de demokratiske Stormløb til Rodomontader\footnote{Storpralerier.}
og blind Allarm, vil det kunne ende med, at den almene
Sands for store Interesser efterhaanden sløves, og Tumulten
afløses af en døsig Ro. Det Sidste er, i folkelig
Forstand, lige saa sørgeligt som det Første.

Vi have hidtil forudsat, at det kun var en af Parterne,
der ved Overgreb vilde krænke Forfatningen;
men et saadant Tilfælde hører til det Sjældnere. Det
langt Sædvanligere er, at Overgreb fra den ene Side
ægger til Overgreb fra den anden. Vi have et Exempel
herpaa i Sveriges Historie i Perioden fra Karl XIItes
Død til Revolutionen under Gustav III. „Karl XIItes
Død og hans Søster Ulrikke Eleonores Valg til Dronning
(1719)“ --- hedder det\footnote{Verdenshistorie efter Cæsare Cantù, frit bearbeidet af E. Holm og Ph. Weilbach. S. 498 flgd.} --- „havde gjort Sverige
fra en enevældig Stat til et næsten republikansk
% --- p. 30 ---
\setcounter{footnote}{0}%
\opage{30}Rige. Ulrikke Eleonore fratraadte allerede Aaret efter
sin Thronbestigelse Regjeringen, som hun fik overdraget
til sin Mand Frederik af Hessen, en Fyrste, der ingenlunde
fattedes Evner, men var overfladisk, nydelsessyg
og stedse stod som en Fremmed i sit Rige. Had imod
Enevælden og Uvillie over de Krogveie, ad hvilke Frederik
og hans Dronning stræbte at faae saa stor Magt
som muligt, fremkaldte paa flere Rigsdage (1720 og 1723)
en Række saadanne Forfatningsbestemmelser, at man
skal have Ondt ved at opvise noget andet Kongedømme,
hvor Kongen har havt mindre at sige eller været stillet
paa en Maade, der var mere trykkende for hans personlige
Værdighed\dots{} Et stærkt bevæget politisk Liv udviklede
sig i Sverige under denne republikanske Forfatning.
Der skabtes en Vækkelse og en Øvelse i kritisk
Drøftelse, som forplantedes paa andre Omraader
til Gavn for Udviklingen af Selvstændighedsfølelsen og
den frie Rørelse. Man vandt Kjendskab til de parlamentariske
Former og gjorde ved de talrige Brydninger
en Række Erfaringer, der ere blevne til Gavn for det
politiske Liv senere i Sverrige. Tanken om, at Folket
havde Ret til at opstille lovmæssig fastsatte Skranker
imod Kongemagten, slog Rod, for ikke mere at forsvinde.
Men den Eensidighed, hvormed Vægten var lagt hos
Stænderne, lod i Forbindelse med andre uheldige Forhold
efterhaanden saa stærke Skyggesider træde frem
ved „Frihedstiden“, at man kan forstaae den Uvillie,
hvormed de fleste Svenskere see tilbage paa den.
--- Da Frederik af Hessen var død, besteg efter den fastslaaede
Thronfølgeorden Adolf Frederik Thronen (1751).
Den Ligegyldighed og Sløvhed, hvormed hans Forgænger
næsten havde ladet Kongemagten gaae i Forglemmelse,
% --- p. 31 ---
\setcounter{footnote}{0}%
\opage{31}afløstes nu af Forsøg paa at udvide Kongens Indflydelse,
hvoraf Følgen var dennes endnu dybere Fald. Adolf
Frederiks milde, godmodige Natur og stærkt begrændsede
Aandsevner syntes at gjøre ham vel skikket til at nøies
med det beskedne Raaderum, som Forfatningen hjemlede
ham; men hans Dronning, den preussiske Konge Frederik
IIdens Søster Louise Ulrikke, der fuldstændig
havde Herredømmet over ham, var lige saa ærgjerrig
som begavet“. --- I hvilken Grad Dronningens Plan til
en Revolution dog mislykkedes, sees af Følgende. „Et
Par underordnede Deeltageres Ubesindighed voldte, at
Planen blev bekjendt, og Straffen fulgte Opdagelsen i
Hælene. Kongen, der Intet havde vidst om Planen,
brød selv i en Skrivelse til Raadet Staven over de
Mænd, der havde vovet Livet for ham, men frelste ikke
derved sig selv og Dronningen for nye Ydmygelser.
Medens otte af Sammensværgelsens Deeltagere bleve
henrettede, gav Præstestanden den stærkt compromitterede
Dronning en skarp Irettesættelse, hvori den lod
hende høre at hun havde tilsidesat sin Pligt imod Gud,
Kongen og Sveriges Rige. Kongen selv maatte i en
Skrivelse fra Stænderne see udtalt, at han havde brudt
sin Ed, saa at Rigsdagen, om den vilde, havde Ret til
at erklære det Baand for løst, der knyttede Undersaatterne
til ham.“ --- Hvad der ikke lykkedes Adolf
Frederik og Louise Ulrikke, lykkedes deres Søn, Gustav
III., som efter Faderens Død (1771) besteg Thronen.
Da alle lemfældigere Forsøg paa at udjævne Partiernes
(Hattenes og Huernes) lidenskabelige Stridigheder sloge
feil, modnedes hans Beslutning om at vove en Revolution.
Uden at oppebie Sprængtportens og de finske
Soldaters Komme, forstod han „ved sin raske begeistrende
% --- p. 32 ---
\setcounter{footnote}{0}%
\opage{32}Optræden at vinde de Tropper for sig, der vare i Stokholm.
Borgerne sluttede sig til ham, medens Rigsraadet
og sekreta utskottet var raadvildt, saa at Revolutionen
neppe var begyndt, førend den var seierrig
endt. Efter at al Modstand var forstummet, holdt
Kongen en Tordentale til de forsamlede Stænder om
de Ulykker, som Frihedstiden havde bragt over Landet.
Derpaa lod han en ny Forfatning oplæse, og da han
fra Thronen, med Gustav Adolfs Sølvhammer i Haanden,
spurgte Stænderne, om de samtykkede deri, lød et tre
Gange gjentaget Ja som Svar. Den samme Rigsdag,
der for faa Dage siden havde viist det hensynsløseste
Overmod, dristede sig nu ikke til den svageste Modstand
for at værne om den Folkefrihed, som hvis Bærer
og Forsvarer den havde pleiet med Bram at optræde.
Udenfor Stokholm sluttede Folket sig med Glæde til
den nye Tingenes Orden, saa at Frihedstidens Forfatningsbygning
styrtede sammen som et Korthuus, der
vælter ved det første Aandepust.“

Fra et folkeligt Synspunkt stiller det constitutionelle
Monarchie sig i en heelt anden Belysning. Det
Monarchiske er ikke en beklippet Rest af Enevældens
Absolutisme, saa lidt som det Constitutionelle er et
Produkt af en halvtæmmet Demokratisme; men begge
ere Sider og Frembringelser af Folkevilliens eget Væsen.
Eenhedsvillien, repræsenteret og levendegjort i Kongens
Person, er et ligesaa organisk-nødvendigt Led i en
folkelig fri Forfatning som den af Befolkningens Ombud
repræsenterede Fleertalsvillie. Grundloven maa desaarsag
være affattet saaledes, at Folkevillien under de
to Hovedformer kan træde kraftig frem og holde sig i
stadig lovbestemt og organisk Vexelvirkning med sig selv.
% --- p. 33 ---
\setcounter{footnote}{0}%
\opage{33}Der er Fare paa Færde, naar den gjentagne, i lovlig,
utvetydig Form udtalte Fleertalsvillie krænkes; men
der er ogsaa Fare for Folkets Sag, naar den lovlige
Form for Eenhedsvilliens Afgjørelser ikke respecteres.

Saa længe et Folk endnu ikke er blevet myndigt,
vil det bestandig leve hen i den Indbildning, at Konge,
Regjering, Embedsmænd, Intelligens og hvad derunder
hører, tilsammen udgjøre en fremmed Magt, der stadig,
om ikke altid aabenbart, saa dog hemmeligt, lægger an
paa at kue og trykke Folket saa meget som muligt;
det vil sige: Folket er i sin Umyndighedstilstand ikke
naaet videre end til at fatte og føle sig som Almue.
Denne Umyndighed veed Demokraten ypperlig at benytte,
naar han træder op som Demagog. Hvad er en
Mand af Folket? Det er en Mand af Almuen; Enhver,
som enten ved Fødsel eller ved Formue eller ved Intelligens
hæver sig lidt over Almuens Trin staaer allerede
udenfor Folket og er en mistænkelig Person. Den demagogiske
Vei til folkelig Selvstyrelse er tydelig herved
betegnet: ved at slaae Formynderen ned skal den
Umyndige styre sig selv.

Jo mere et Folk arbeider sig frem til Myndighed
og Mands Modenhed, desto klarere vil det indsee, at
den styrende Magt, der nødvendig maa staae over Almuens
Trin, ikke falder udenfor, men indenfor dets
egen Villies Omraade. Et myndigt Folk vil lige saa
lidt oversee Almuens Ret til virksom, lovbestemt
Deeltagelse i den hele Samfundsudvikling som det i
Et og Alt vil reducere sig selv til Almue og gjøre Statsstyrelsen
til en slet og ret Almuesag. Et myndigt Folk
vil sørge for, at Regjeringen og dens Embedsmænd ikke
handler egenmægtig i en mod Folkevellet stridende,
% --- p. 34 ---
\setcounter{footnote}{0}%
\opage{34}bureaukratisk Aand; det vil sikkre sig, at Staten ikke
ansætter eller lønner sine Embedsmænd paa en Maade,
som staaer i Misforhold til det Maal af aandelige Kræfter,
de anvende i Samfundets Tjeneste; men naar de
behørige Sikkerhedsmidler i denne Henseende ere anvendte,
da vil Folket ogsaa vide at holde Regjeringsmyndigheden
og dens Tjenere i Hævd; thi Statens
Embedsmænd ere jo Folkets egne Embedsmænd, Folkevilliens
Organer. Et Folk, der veed, hvad det gjør,
vil lige saa vist finde det fornærmeligt, ja oprørende,
hvis Nogen tilraader det at svække dets egen Villies
Organer, som en fornuftig Mand vil finde det oprørende,
hvis Nogen er saa taabelig at raade ham til at svække
de Muskler og overskjære de Nerver, ved hvilke han
dog alene er i Stand til at bevæge sine Lemmer. For
et myndigt Folk staaer det klart, at ogsaa den udøvende
Magt bør have et frit Omraade, at de Styrendes
Autoritet er Folkets egen Autoritet, Lærernes Intelligens
dets egen Intelligens o. s. v.

Den constitutionelle Medvirkning mellem Eenhedsvillien
og Fleertalsvillien maa i Grundloven være saa
tilstrækkelig bestemt, at ethvert Overgreb saavel fra
den ene som fra den anden Side kan forfatningsmæssig
bekæmpes og tilbagevises. Dette er i kort Begreb
Betydningen af \emph{folkelig Selvstyrelse}.

Kongens Villie og Folkets Villie ere ikke to, men een
Villie; Folkets Majestæt og Kongens Majestæt ere ikke
to, men een Majestæt. I den kongelige Majestæt ærer
Folket sin egen Souverainetet; det anskuer i Kongens
Person sin egen personlige Eenhed. Kongens Frihed
og Folkets Frihed indskrænke ikke hinanden; det er
% --- p. 35 ---
\setcounter{footnote}{0}%
\opage{35}een og samme folkelig constitutionel Frihed. At Kongen
staaer over Folket, vil sige, at Loven staaer over Folket.
Kongens Villie er Lovens Fuldbyrdelse: Kongen er den
levende Lov; i denne Mening forstaae vi Ordet: „\emph{en fri
Konge over et frit Folk}.“

% --- p. 36 ---
\setcounter{footnote}{0}%
\lecture{III.} {Politiske Erfaringer: Parlamentarisme.}
\opage{36}Erfaring, siger man, er en god Læremester, og det er
sandt; men Læremesteren er undertiden noget dyr, ifølge
det gamle Ord: af Skade bliver man klog, men sjældent
rig. Dersom man saa endda altid blev klog, saa fik
det at gaae med Skaden; men det er ikke sagt, at man
altid bliver klog, man kan desværre ogsaa blive dum
af lutter Erfaring.

Det er ikke saa ligefrem en Sag at gjøre Erfaringer,
som Mange bilde sig ind, at sige: naar Erfaringerne
skulle due noget. For at overvinde Vanskeligheder,
maa man kjende dem. Som Indledning til en nærmere
Overveielse af de Vanskeligheder, som stille sig i Veien
for dem, der lægge Vind paa politiske Erfaringer, kunne
vi begynde med nogle Erfaringer, hentede fra andre
Kredse.

Naturlæren er en Erfaringslære; her er Kjendsgjerninger
nok, og nok at erfare. Men hvad har det
ikke kostet, inden Forskningen naaede saa langt frem,
at den kunde fastsætte den rette Fremgangsmaade og
sikkre sig mod falske Slutninger. Det er en gammel
Erfaring, at der undertiden kort før Solens Opgang
% --- p. 37 ---
\setcounter{footnote}{0}%
\opage{37}viser sig en klar Stjerne paa den østlige Himmel, d. e.
Morgenstjernen, og til andre Tider umiddelbart efter
Solens Nedgang ligeledes en klar Stjerne paa den vestlige
Himmel, nemlig Aftenstjernen. Den Mening laa da
nær, at det maatte være to forskjellige Stjerner; men
endskjøndt det seer ud som en virkelig Erfaring, er det
dog kun en falsk Slutning. Oldtiden blev overrasket,
da En af dens Viismænd opdagede, at Morgen- og
Aftenstjernen var een og samme Stjerne. Det er en
gammel Erfaring, at Stene falde, og at Røg stiger i
Veiret. Den græske Philosoph Aristoteles byggede paa
slige Erfaringer den Lære, at nogle Legemer af Naturen
bevægede sig nedefter, andre opefter, og at der
maatte gjøres en væsenlig Adskillelse mellem Legemer
med og uden Tyngde. Men Physiken lærer os, at Forestillingen
om Legemer uden Tyngde beroer paa en
falsk Erfaring. Erfaringer i den sandselige Verden
gjøres ikke blot ved at høre, see og føle; det er nødvendigt
at kunne tænke, dømme og slutte.

Er det i Naturriget, hvor dog Alt gaaer efter nødvendige
Love, vanskeligt at hitte Rede i Erfaringernes
Mangfoldighed, hvor vanskeligt maa det da ikke være
i Frihedens Verden, i Menneskelivet, hvor et Valg finder
Sted, og hvor Villiens Valg ikke altid følger Fornuftens
Bud, men hellere lader sig lede af Drifttilskyndelser
og uberegnelige Tilbøieligheder. Her
kommer det ikke alene an paa, hvad der erfares, men
først og sidst, hvorledes det erfares. Det er en bekjendt
Sag, at to Mennesker kunne opleve samme Begivenhed
og endda faae det Allerforskjelligste ud deraf.
Erfaringsslutninger ere ofte i høi Grad overilede og
skæve, uagtet Vedkommende, der bilder sig ind at have
% --- p. 38 ---
\setcounter{footnote}{0}%
\opage{38}Kjendsgjerninger for sig, selv anseer dem for ubedragelige.
Det er en Erfaring, at Mænd kunne bedrages af
Kvinder, og Kvinder af Mænd. Den Mand, som er saa
ulykkelig at blive bedragen een Gang, to Gange, ja
maaskee tre Gange, faaer, hvis det ret gaaer ham til
Hjerte, Erfaringens Attest for, at tre Gange ere alle
Gange, og bliver Fruentimmerhader. Kvinder, hvis
finere Nerver ikke taale nogen Gjentagelse, kunne have
nok i en eneste Erfaring for at blive Mandfolkehadere.
Ja, hvis man tør fæste Lid til Digteren --- og Digteren
maa jo forstaae sig paa Sligt --- skal der endog gives
kvindelige Individualiteter med en saa fiint mærkende
Sands, at de, just fordi det ikke er lykkets dem selv
at gjøre Erfaringer, af Erfaring vide, at alle Mandfolk
ere Bedragere. Om den erfarne Tante, der havde paataget
sig at vaage over „Alhambras Rose“, bemærker
Washington Irving: „Den aarvaagne Fredegunde var
den forsigtigste gamle Skabning, man kunde tænke sig.
Med Hensyn til „det fjendtlige Kjøn“, som hun benævnede
det, nærede hun en agtværdig Rædsel og Mistillid,
som ved Forlængelsen af hendes ugifte Stand efterhaanden
var tiltagen. Selv havde den gode Dame rigtignok
aldrig lidt ved Kjønnets Rænker, idet Naturen i
hendes Ansigt havde givet hende en Beskyttelse mod
alle Efterstræbelser; men Kvindfolk, som have mindst
Aarsag til at frygte for deres eget Vedkommende, ere
netop mest redebonne til at vaage over deres mere
fristende Medskabninger.“

Hvad Børneopdragelse angaaer, bliver man ikke
sjældent Vidne til, hvor indbyrdes uenige Erfaringer kunne
være. Den ene Part vil have Erfaring for, at Barnet
maa under streng Tugt, naar det ikke skal blive uvorent
% --- p. 39 ---
\setcounter{footnote}{0}%
\opage{39}og egensindigt, den anden Part kan af Erfaring bevise,
at det kjære Barn kun bliver trodsigt, naar man tager
det med det Onde, men altid er artigt, naar man føier
sig lidt efter det og tager det med det Gode.

Er Talen om Folkeoplysning og lærd Underviisning,
kan det undertiden, især ved livlige Forhandlinger, gjøre
En urolig at see, hvorledes Erfaring staaer op mod
Erfaring. Vi holde os forvissede om, at Folkets sunde
Sands er aaben for Livets høieste Interesser, at det
menige Folk har en hurtig, uvilkaarlig sikker Fatteevne,
en naturlig skarp Forstand, der kun tager feil,
naar den hildes i uklare Stemninger og ledes vild af
ubegrundede Forestillinger. Det er, efter vor Synsmaade,
en paatrængende Nødvendighed, at Almuen,
den menige Befolknings store Masse, paavirkes af Lysets
og Sandhedens Ord og faaer Øinene op; det er
vort Haab, at Sagen vil gaae, naar der er Lærere nok,
som vide Besked og tale reent ud, uden at smigre
Mængden og fiske efter Fleertal. Men kan dette ligefrem
bevises af Erfaring? Ikke saa ganske. Der er
grundlærde Mænd, som mene at have Erfaring for, at
den saakaldte Folkeoplysning skader istedenfor at gjøre
Gavn. „De ulærde Overblik ere kun Overhalinger, der
mangle Grundighed. Og hvorfor? Fordi de mangle
Kritik, Archæologie, Philologie og Grammatik. Løse
Antydninger og spredte Lysglimt sætte Sindene i Gjæring
og gjøre Menigmand indbildske. Heste, der arbeide
strængt, bør ikke fodres med Græs: de maae have
Havre. Underviisningens Havrekjerne ligger i Archæologie
og Philologie; den „folkevækkende“ Oplysning er
en Grønfodring, der blæser op, men ikke giver Kræfter.
Kjærnen i Alt, hvad det menige Folk skal lære, er kort
% --- p. 40 ---
\setcounter{footnote}{0}%
\opage{40}og eenfoldig dette: det skal lære at see op til Lærdommen
og vise Respekt for de Sagkyndige. Intet er
mere fordærveligt end „det levende Ord“, hvormed man
rører omkring i Alt. Hvad er dette „levende Ord“?
Det er rørig Snak. Men rørig Snak er som rørig Mad;
det levende Ord puster op som Kaal og kan saa let
give Trommesyge. De oppustede, trommesyge Ordførere
ere halvstuderede Røvere, der tale med om Alt, men
ikke vide Besked med Noget.“ --- See, det var nu lærde
Erfaringer, og det er i Sandhed ingen Spøg at give sig
i Kast med lærde Erfaringer.

Mod de lærde Erfaringer stille sig de ulærde. Vi
agte Lærdom og Videnskab, paa hvilke Omraader Undersøgelserne
end bevæge sig, det være nu i det Archæologiske,
det Philologiske eller det Geologiske. Men der
er folkelig Opvakte, som i den Grad have erfaret, hvor
huul, ufrugtbar og tom den lærde Forskning er, at de
paa staaende Fod kunne læse den Texten. „Hvad hjælper
det“, sige de, „at grave i Jordlag og rode i gamle
Dokumenter? Kan man grave Livet op, som man opgraver
Lig? Samlinger fra Steenalderen, Broncealderen,
Jernalderen o. s. v. kunne jo være underholdende nok
for dem, der lade sig begrave i et Musæum, men vi,
der aande under Guds frie Himmel og ikke føle Trang
til at blive levende begravne, hvad Udbytte have vi
deraf for vort nærværende Liv? Skal man af det Liv,
Oldfædrene levede i Steenalderen, regne ud, hvorledes
vi Nutidsmennesker skulle bære os ad med at
leve i Politikalderen, da behøves der mere end een
Skjæppe Ærter. Man gnaver i de mange Sprog og de
mange Grammatikker: Chinesisk, Persisk, Tamulisk,
Malebarisk, Sanskrit, Prakrit, Sprog og Sprog, Gram%
% --- p. 41 ---
\setcounter{footnote}{0}%
\opage{41}matik og Grammatik! Hvad hjælper det Alt, naar man
dog til Syvende og Sidst ikke kan gjøre levende Brug
af sit eget levende Sprog? Det store Oldtidsskramleri er
en stor „Kjøkkenmødding“, og det endeløse Sprogforskeri
en uhyre Grammatik. Man lægge den archæologiske
Mødding paa Grammatikken, om man hellere vil, Grammatikken % defect: „o n“ (a broken m): the letter supplied
paa Møddingen; man kalde Physiken til Hjælp
og lade en elektrisk Strøm gaae igjennem --- fra Grammatikken % defect: „elcktrisk“ (the bar of the e did not ink), „igjer\,nem“ (a broken n): letters supplied
til Møddingen, fra Møddingen til Grammatikken
--- der kommer dog ikke Glimt af Liv, ikke Gnist af
Aand: det Døde er dødt og bliver dødt, hvor længe
man end forsker og graver“. --- Det er ulærde Erfaringer,
og ulærde Erfaringer er det ingenlunde let at bekæmpe.

Tage vi det Hellige med og besinde os paa de religiøse
Erfaringer, da er Striden endnu mere gjennemgaaende.
De Mange, der samlede sig om Christus, og
vare Øienvidner ved hans Mirakler, gik bort med høist
forskjellige Erfaringer. Nogle erfarede, at det i Sandhed
var Guds Gjerninger; men Andre erfarede, at det
var blindt Sværmeri, og det Hele et Oprør mod Gud
og Keiseren. Hvor meget erfarede vel Pilatus, da han
spurgte: \emph{hvad er Sandhed}? Der fortælles i Apostlernes
Gjerninger (Cap. 25), at da Kong Agrippa og
Berenice kom til Cæsarea for at hilse paa Festus, var
denne i Forlegenhed og vidste ikke, hvad han skulde
dømme om Jødernes Anklage mod Paulus. Det Eneste,
han havde erfaret, var, at de havde nogle Tvistigheder
om deres Gudsdyrkelse og „om en Jesus, som er død,
hvilken Paulus sagde at leve.“ Det er et charakteristisk
Træk. Den hele Begivenhed med Christi Død og
Opstandelse stod for Festus som noget saa Ubetydeligt
og Forbigaaende, at det var uværdigt for en Romer at
% --- p. 42 ---
\setcounter{footnote}{0}%
\opage{42}sætte sig nærmere ind i Sagen. Saaledes gik det i
mere end to Aarhundreder med de romerske Erfaringer
og Christendommen. Christuspartiet blev paa Forhaand
betragtet som et vanvittigt Almueparti; om et saadant
Almueparti var. det da tilladt at troe det Urimeligste, % print: „var.“ (a stop before a lower-case word)
det Allertaabeligste, der kunde opdigtes: man havde jo
Erfaringen for sig.

Vi vende os nu til Emnet om de politiske Erfaringer.
Disse Erfaringer ere behæftede med alle de
Vanskeligheder, vi i det Foregaaende have mødt. De
skulle rette sig efter almeengyldige Love; men Frihedens
Love ere saa bevægelige, at Tilbøieligheder og
Stemninger friste til overilede Slutninger og kunne give
os en aldeles skæv Opfattelse af de virkelige Forhold.
Til at gjøre sande Erfaringer kræves et uhildet Sind,
der formaaer at frigjøre sig fra alle forudfattede Meninger;
men hvor personlige Interesser ere med i Spillet,
ville forudfattede Meninger let sætte sig fast som
ubedragelige Sandheder, ufravigelige Grundsætninger.
For at gjøre sande politiske Erfaringer er det ikke
alene nødvendigt, at de Paagjældende i al Almindelighed
have Agtelse og Sands for Sandhed og Ret: her
fordres endnu noget Mere. Her fordres Selvovervindelse,
for at sætte sig i Modstanderens Sted og fra Grunden
af forstaae hans Anskuelse. En saadan Selvovervindelse
falder ikke i Alles Smag. De Kyndige ere lige saa
fristede til at forblive i deres een Gang afsluttede Forestillingskreds,
som de Ukyndige i deres. Det falder
saa behageligt at leve med lutter Ligestemte og høre
Omgivelsens Echo af sine egne Meninger. Deraf kan
det forstaaes, at mægtige Partier kunne kæmpe en heel
Tidsalder og med stor Skarpsindighed udspeide hin%
% --- p. 43 ---
\setcounter{footnote}{0}%
\opage{43}andens Vrangsider uden dog at bringe det til at gjennemskue
hinandens Væsen. Vi ville betragte Vanskelighederne
først fra oven, dernæst fra neden.

See vi ovenfra nedad, da er det jo en Selvfølge,
at en Konge paa sin Throne maa gjøre politiske Erfaringer.
Men Erfaringerne blive desværre farvede,
naar de sees igjennem farvet Glas. Der er en Atmosphære,
som kaldes Hofatmosphære: den har en
egen lysbrydende Virkning. Det er bekjendt, at Lyset
fra en Stjerne bøies ved at gaae igjennem forskjellige
Luftlag, saa at det tager sig ud, som om Stjernen stod
noget høiere end den i Virkeligheden staaer. En saadan
Tilsyneladelse beroer paa, hvad man i Astronomien
kalder Refraction. Jo nærmere Stjernerne ere ved
Horizonten, desto mere virker Tilsyneladelsen, saa de
synes at staae høiere. Astronomen veed, at han maa
beregne Refractionen. For Monarkens Iagttagelse vil
det, paa Grund af Hofatmosphærens lysbrydende Evne,
kun altfor let tage sig ud, som stode de lavere Samfundsklasser
høiere, med andre Ord: som stod det
bedre til med Undersaatterne i hans Rige, end det
virkelig gjør; det er ham ingenlunde altid bekvemt at
beregne Refractionen. Da Katharina II. havde erobret
Krim, og 30,000 Tartarer vare blevne nedsablede, vilde
Potemkin vise sin keiserlige Elskerinde et storartet
Skue af Pragt og Herlighed. Keiserinden gjorde en
Reise for at besøge sine nyerobrede Lande. „Potemkin
samlede ved Dniper en stor Hær og ved Hjælp af dygtige
Decorationsmalere opdigtede han, saa at sige, en
blomstrende Tilstand af Landet, som det aldrig havde
kjendt. Bredderne vare overstrøede med Landsbyer,
men de vare malede paa Lærred: man saae Kirker
% --- p. 44 ---
\setcounter{footnote}{0}%
\opage{44}under Bygning, Skibe og Huse. Med Knutten vare
Tartarer langveis fra drevne ned til Flodbredderne, for
at disse kunde synes befolkede; det unge Græs blev
afgnavet af Kvæg, der var hentet i 400 Versters Omkreds.
Der havde neppe været større Udgifter ved at
faae noget Nyttigt indrettet. Men Katharine lod sig
skuffe for at kunne skuffe Europa med Hensyn til Rigets
Kræfter og sin egen Virksomhed.“ --- Hvad der
blænder en Konges Øine er vel i sjældnere Tilfælde
lave Smigrerier; thi i saa Henseende vil Erfaringen nok
lære ham at trække fra. Men hvis der er en Kreds af
fortrolige Raadgivere, paa hvis Ufeilbarhed han, uden
selv at have tilstrækkelig Indsigt i Sagernes Gang,
ubetinget vil stole, da kan en saadan blive ham farlig.
Naar Himlen er overtrukken med et Slør af lette Skyer,
danner der sig jevnligt om Solen Ringe og Bisole;
saaledes vil der i Hoffets Dunstkreds ogsaa om den
kongelige Sol let danne sig Bisole; en Skriftefader, en
Førsteminister eller hvad Oraklet nu hedder, vil sørge
for, at Majestæten aldrig seer den nøgne Virkelighed,
men Billeder i et Speil, der kaster Glandsen af hans
eget Lys tilbage. Selv om Virkeligheden bliver saa
nærgaaende, at den ved en Omvæltning river den fyrstelige
Person ud af Tryllekredsen og lader ham gjennemgaae
de store Prøvelsers strenge Skole, saa er det
dog dermed ikke givet, at han gjør sande Erfaringer.
Man tænke, for at minde om Stuarternes Historie, paa
Karl II. Hvad havde denne Fyrste ikke oplevet. Hans
kongelige Fader maatte lægge sit Hoved paa Skaffottet;
selv havde han i flere Aars Landflygtighed maattet
gjennemgaae de Ydmygelser, som altid møde den faldne
Storhed. Og da han saa omsider --- efter Cromwells
% --- p. 45 ---
\setcounter{footnote}{0}%
\opage{45}Død --- blev kaldt tilbage og drog ind i sin Hovedstad
under Folkets Jubel, hvilke Erfaringer bragte han da
med sig op paa Thronen? „At hver Kvindes Ære var
tilfals, og hver Mands Troskab at faae for Penge.
Forskjellen var kun i Prisen: Nogle solgte sig for billigt
Kjøb, Andre lode sig dyrt betale.“ Man læse
Macaulays storstilede Fremstilling af Carl II.'s sidste
Levedage, og man vil see, hvad det var for Livserfaringer,
denne Konge havde høstet. Om Bourbornerne har % print: „Bourbornerne“ (misprint for Bourbonerne)
Napoleon træffende sagt: „de have Intet glemt og de
have Intet lært.“ De have altsaa, trods Revolution og
Keiserdømme, kun dreiet sig i deres egne Indbildningers
Kreds, i gamle vindskæve Erfaringer.

See vi nedenfra opad, mon det da skulde være
lettere for Folkene at gjøre sande politiske Erfaringer,
end det er for Fyrsterne? Ingenlunde. Dersom Magtstillinger
i Samfundet lode sig gruppere som Træer i
Skoven, dersom Styrelsens Gang i Op og Ned var anskuelig
som Sivets Sveien ved Bredden af en Sø eller
som Kornets Bølgen paa Ageren: da skulde den sunde
Forstand let bane sig Vei og skjelne mellem Ret og
Uret. Men Striden om Magt er en indviklet Strid;
ingen Forførelseshistorie er saa vidtløftig som den, der
handler om gjensidig politisk Utroskab mellem Folkene
og deres Førere.

Som oftest er det vel Føreren, der misbruger Folkets
Tillid og svigter den gode Sag, naar uforudseete
Vendinger indtræde. Dog gaaer det i Reglen ikke saaledes
til, at Vedkommende træder op i det fuldt bevidste
Øiemed at ville bedrage, at ville indsmigre sig
og vinde Folkegunst blot for selv at komme til Magten.
Vi maae til Ære for Menneskeslægten tvertimod an%
% --- p. 46 ---
\setcounter{footnote}{0}%
\opage{46}tage, at de Fleste blandt dem, der henvende sig til
Folket, begynde i ærlig Sympathie med Folket. Men
derfor kan det ligefuldt ende med et Bedrag. Den
moralske Selvfølelse, som vækkes hos Føreren derved,
at Mængden slutter sig til ham, er udsat for store Fristelser.
Den Feirede føler sig alt som Magthaver, og
denne Følelse er meget kildrende; den driver ham frem:
han kan ikke nøies med moralsk Indflydelse, han trænger
til Nævekræfter; han maa nødvendig kommandere
et Parti og vælge Midlerne, som de falde. Saa tirres
han af Modstandere, saa ægges han af Tilbedere, saa
bliver han Massehelt, Kommandør af Fleertalsordenen,
og Despoten i ham vaagner. Det kan da skee paa den
allernaturligste Maade, at den, der begynder med Folkets
Sag, kan gaae over til at tjene sin egen Sag og
ende med at forraade Folket.

Men Forholdet kan ogsaa vende sig saaledes, at
det er Folket, der forraader Føreren. Der er hos det
menige Folk --- saa længe det endnu befinder sig paa
Umyndighedens Trin --- en eiendommelig Vexlen mellem
Godtroenhed og Mistro. Af Naturen er et ufordærvet
Folk godtroende; det gjælder som Regel. Enhver, der
træder op og vil antage sig dets Sag, finder let en Tilslutning.
Men jo blindere den Tillid og Begeistring er,
hvormed det følger sin Fører, desto brattere er igjen
det Omslag, hvorved det opgiver ham, naar dets Mistanke
vaagner. Grækenlands og Roms Historie have
herpaa talende Exempler. Trang til en Omordning af
Jordbesiddelsen og Striden om Agerloven var for Romerne,
ved noget over et Aarhundrede før Christi Fødsel,
hvad Spørgsmaalene om Pauperisme og Fattigunderstøttelse
ere i vore Dage. Brødrene Tiberius og
% --- p. 47 ---
\setcounter{footnote}{0}%
\opage{47}Cajus Gracchus, begavede Mænd af ædel Byrd, med
sjældne Dyder og de redeligste Hensigter, opbøde Alt,
for at hjælpe det stakkels forarmede Folk paa Fode;
men deres Bestræbelser vare frugtesløse, de maatte see
sig stemplede som Forbrydere og bøde med Livet.
Adelen med Senatet forstod fortræffelig at føre Hoben
bag Lyset: den uvidende Almue var dum nok til at
lade sig blænde, feig nok til at svigte sine Beskyttere.

En nødvendig Betingelse for, at et Folk kan gjøre
Fremskridt i politiske Erfaringer, er blandt Andet dette,
at Forhandlingerne om de offenlige Anliggender føres
besindig, saa at Grundene for og imod kunne komme
til Almeenhedens Kundskab; men det er netop en Betingelse,
som vanskelig opfyldes. Naar Partierne komme
i Lidenskab, faaer man ikke Tid til at tænke paa Sagen;
man tumler med det Personlige: Enhver har travlt
med at anfalde, bagvaske og sværte sin Modstander.
Læsere. der ere vante til Sværten, kunne gjerne vide % print: „Læsere.“ (a stop before a lower-case word; misprint for a comma)
af Erfaring, at nærgaaende Beskyldninger just ikke
altid skulle tages efter Bogstaven; men de, der speculere
i Ophidselser, vide ogsaa af Erfaring, at en dristig
Forvrængning sjældent forfeiler sin øieblikkelige Virkning.
Det demokratiske Amerika har Greb paa at anvende
Methoden. Washington tog det med Ro. Skjøndt
han blev voldsomt angreben af det demokratiske Parti,
vrededes han dog ikke; han skrev kun til dets Leder,
Jefferson: „Jeg vilde aldrig have tænkt mig, jeg vil ikke
sige, at det var sandsynligt, men blot, at det var muligt,
at, medens jeg brugte mine bedste Kræfter til at
skabe en folkelig Politik, som ganske var vor egen, og
at sikkre Landet mod Krigens Rædsler, alle mine Handlinger
da skulde blive fordreiede, fremstillede paa en
% --- p. 48 ---
\setcounter{footnote}{0}%
\opage{48}plump og fornærmelig Maade, i saa overdrevne og
usømmelige Udtryk, at de neppe nok kunde anvendes
paa en Nero, paa en navnkundig Bandit, men snarere
paa en dagligdags Slyngel. Men det er nok, jeg er
allerede gaaet for vidt i at udtale mine Følelser.“

Vil Du gjøre \emph{sande Erfaringer, saa behersk
dine Lidenskaber og brug din Forstand}. Denne
Washingtonske Grundsætning er kun lidet tiltalende
for den store Hob, og heri stikker Knuden. Det er
ikke Mangel paa Kundskaber, thi derpaa kan der bødes,
men det er Mangel paa Selvbeherskelse, der forvirrer % defect: „Sclvbcherskelse“ (the bars of two e did not ink): letters supplied
Folkets Sands og bereder den folkelige Sag saa mange
Hindringer. Der kan naturligviis Intet indvendes imod,
at Førerne for et Parti fremstille Sagen fra deres Synspunkt
og søge at støtte deres Opfattelse ved saa stærke,
indtrængende, overbevisende Grunde som muligt; men
at fortie Modpartens Grunde og fremstille hans Udtalelser
som simple, uanstændige Jammerligheder, er
det ærligt, er det folkeligt? At man vil holde et Speil
op for Folkets Øine, er i sin Orden; men det er, som
bekjendt, ikke alle Speile, hvis Billeder svare til Gjenstandene.
Man har i Haverne de saakaldte Kuglespeile.
Om de Spadserende end ere nok saa velskabte, naar
de komme nær til et saadant Speil, undergaae de dog
en uhyggelig Forandring: store Næser, tykke Maver,
smaa bitte Been! de see ud som lutter Nøddeknækkere.
Der er andre Speile, der forherlige Skikkelserne og vise
dem i overnaturlig Størrelse. Naar nu Partiførerne, de
Mænd, som udgive sig for at holde Sandhedsspeilet op
for Mængdens Øine, stadigt kombinere de to Slags
Speile saaledes, at de selv komme til at vise sig som
Helte, medens deres Modstandere staae som Nødde%
% --- p. 49 ---
\setcounter{footnote}{0}%
\opage{49}knækkere, er det saa at veilede Folket til politiske,
sande politiske Erfaringer?

Paa Fremskridt i politiske Erfaringer beroer Indflydelsen
af den saakaldte Parlamentarisme. Vi sige:
Indflydelsen; thi Tilblivelsen falder sammen med Grundlægningen
af en hvilken som helst friere Forfatning.
Allerede med Indførelsen af raadgivende Stænder følger
en vis Art af Parlamentarisme. Det første for Parlamentarismen
Charakteristiske er, at Styrelsens vigtigste
Anliggender gjøres til Gjenstand for en Forhandling
mellem Regjering og Folkerepræsentation. Folkerepræsentationens
Indflydelse kan gjennemgaae forskjellige
Grader; men først, naar Forhandlingernes Resultat afgjøres
ved Afstemning saaledes, at Fleertallets Afstemning
tillige bliver afgjørende for et Ministeriums Bestaaen
eller Ikke-Bestaaen, træder Parlamentarismen i
fuld Kraft. Det er til Parlamentarismens Gjennemførelse
nødvendigt, at Fleertallets ved Afstemningen udtalte
Anskuelse faaer en faktisk og juridisk Indflydelse
paa Sagernes Gang; men det er ingenlunde nødvendigt,
at Fleertallet skal være samlet i eet Kammer; det kan
jo ogsaa være fordeelt paa tvende. For at blive ved
vor egen Grundlov, da skjønnes det ikke, at den gjør
noget Skaar i Parlamentarismen, ved at lade Folkething
og Landsthing have hver sit Fleertal; thi det tør dog
vel forudsættes, at et Ministerium, der faaer begge
Things Fleertal imod sig, maa træde af. Noget ganske
Andet er det, at der endnu savnes Bestemmelser for,
hvad der lovlig skal foretages i Tilfælde af vedvarende
Uenighed mellem det tvedeelte Fleertal. Den ubestemte
Forestilling om Muligheden af provisoriske Foranstaltninger,
er uhyggelig ikke alene af den Grund, at der i
% --- p. 50 ---
\setcounter{footnote}{0}%
\opage{50}det, som skal vise sig, naar Loven ikke strækker til,
altid er noget Bussemandsagtigt; men navnlig af den
Grund, at Fleertalsvillien hverken kan komme til sin
fulde Ret ei heller blive lovformelig begrændset, naar
Grundloven ikke bestemt og tydelig angiver den Vei
den har at gaae, for at naae til Maalet. Som Forbillede
for en consekvent gjennemført Parlamentarisme
pleier man at anprise den engelske Forfatning. Men
at det engelske Underhuus lige saa lidt kan antages
at være et Mønster for vort Folkething, som Overhuset
„kan være et Mønster for vort Landsthing, er indlysende
for Enhver, der har endog blot et overfladisk Kjendskab
til det engelske Parlament.

Parlamentarismen i England har paa eiendommelig
Maade udviklet sig af de engelske Forholds særlige Krav
med historisk Nødvendighed: den har Aarhundreders
politiske Erfaringer til Grundlag. Opdraget til Selvstyrelse,
har det engelske Folk lært at holde sig strengt
til Lovens Bogstav og agte Kronens Ret. Og dog kan
Bølgegangen undertiden blive farlig nok, naar Folkehavet
kommer i Røre. Faren er imidlertid forbigaaende. Med
Undtagelse af nogle philanthropiske Paroxysmer, er det i
Reglen ikke store Ideer, men store Projecter, der beherske
Stemningen. Loven har et Bolværk i Kapitalen; thi
til at udføre store Projecter udkræves Penge. Det er da
ogsaa en Lykke for England, at det i bevægede Tider
har en Veileder saa praktisk som \textit{Times}, en Veirhane
over Børsen, saa politisk omskuende, at den, uden at
skifte Standpunkt, bestandig kan dreie sig efter Vinden.

Paa Fastlandet har Parlamentarismen i de sidste
Aar just ikke høstet Laurbær. Det var i 1870 vistnok
for en god Deel den gjenopvakte Parlamentarisme, som
% --- p. 51 ---
\setcounter{footnote}{0}%
\opage{51}i Forbindelse med Opinionens moralske Paatrængenhed
ved hidsende Midler rev Frankrig ind i Krigen med
Tydskland.

En nærmere Betragtning af Parlamentarismens
Væsen maa føre til den Erkjendelse, at dens Virksomhed
er betinget af to beslægtede, men dog uligeartede
Factorers Samvirken; den ene er den moralske, den
anden den juridiske. Den moralske Factor ligger i
Forhandlingerne; thi hvad skulde Discussionen og Forhandlingerne
vel sigte til, dersom Hensigten ikke var
at faae de Grunde, der tale for og imod en foreliggende
Sag, samvittighedsfuldt og omhyggelig prøvede. Den
Anskuelse, der støttes af de mest overbevisende Grunde,
maa da siges at have den moralske Overvægt. Den
juridiske Factor ligger i Afstemningens Resultat. Naar
Afstemning har fundet Sted, maa den Anskuelse gaae
igjennem, som har de fleste Stemmer for sig; det er
juridisk Ret: her hjælper ingen Ræsonneren. Dersom
de to Factorer nu altid samvirkede i fuld Overensstemmelse,
saa vilde en consekvent gjennemført Parlamentarisme
unægtelig være en politisk Velsignelse for Land
og Folk. Men desværre, Erfaringen lærer, at det Modsatte
ofte er Tilfældet. Den moralske Factor er undertiden
belemret med umoralske Bestanddele.

Der er jo Partianskuelser, der staae saa fast, at
de ikke lade sig rokke ved Modgrunde. Selv om den
evige Fornuft vilde nedlade sig til at tage Menneskeskikkelse
paa og deeltage i Forhandlingerne, den skulde
dog ikke formaae at gjøre Partiet usikkert. Forhandlingerne
føres kun for et Skins Skyld, Sagen er allerede
afgjort i Forveien: det er saa uparlamentarisk
som muligt. Om ogsaa Misbrugen ikke gaaer til en
% --- p. 52 ---
\setcounter{footnote}{0}%
\opage{52}saadan Yderlighed, om der end i Forsamlingen er
Sands og Modtagelighed for Grunde, saa er der dog ---
dersom de sande politiske Erfaringer mangle --- hverken
i den moralske eller den juridiske Factor Sikkerhed
for, at det bliver til mere end Skingrunde. Det er saa
fristende for en begavet Parlamentstaler at anvende
Midler, som blænde, at glimre ved Vid og Humor, at
saare sin Modstander med Satirens Giftbraad og rive
Forsamlingen med sig. Om det han forsvarer, er i den
bedste Overeensstemmelse med hans egen Fortid og vil
holde sig i en kommende Tid, bekymrer ham ikke:
han trylles af Dagens Hæder. og stoler paa, at Publikum % print: „Hæder.“ (a stop before a lower-case word)
er glemsomt. Ja, Publikum er glemsomt; det
har, som Karker siger om Hakon Jarl, „saa mange Ting
i Hovedet“. Men et Folk med politiske Erfaringer er
ikke glemsomt; det lader sig ikke bedaare af et Øiebliks
Fyrværkeri, men husker netop det, hvorpaa det
kommer an, husker, hvem det var, der talede, og hvad
der blev sagt, og husker det længe.

At stole paa den moralske Factor uden at sikkre
sig den juridiske (statsretslige) er upraktisk; men at
stole paa den juridiske alene og lade haant om den
moralske, er i Sandhed ogsaa upraktisk. De, som føre
Ordet i en parlamentarisk Forsamling, maae vide, at de
ikke blot have Tilhørere, men ogsaa Dommere, alvorlige
Dommere, der i Tidens Løb ville kalde dem til
Regnskab for hvert utilbørligt Ord, de have talet.
Først naar det kommer saa vidt, at en Parlamentstaler,
der har vundet en øieblikkelig Seier, om hvilken han
veed, at den, var den end nok saa glimrende, hviler
paa hule Grunde, maa sige til sig selv, hvad en græsk
% --- p. 53 ---
\setcounter{footnote}{0}%
\opage{53}Fyrste sagde til sine Lykønskende: „vinde vi een saadan
Seier til, saa er det ude med os“, først da er
Parlamentarismen gjennemført. Men da vil det ogsaa
vise sig, at den beroer paa mere end en juridisk Afstemningsmechanik,
at dens Grundlag er en suurt erhvervet
Skat af sande politiske Erfaringer.

% --- p. 54 ---
\setcounter{footnote}{0}%
\lecture{IV.} {Fri Forfatning og Frihed.}
\opage{54}Det er nu 28 Aar, altsaa henved en Menneskealder
siden, da begyndte der en ny Tid for Folkene i Europa.
Det var i 48, at Livsaanderne rørte sig: Morgenen
lyste, Dagen brød frem, og Frihedens Sol gik op over
jublende Millioner. Fra Paris, fra Wien, fra Berlin,
fra Kjøbenhavn, overalt lød det samme Budskab. Hvor
de gamle Baand ikke, som hos os, kunde løses med
Lempe, der maatte de give sig for Magt, der maatte
de briste. Hvad der længe havde gjæret i Mulighed, i
skuffede Forventningers, i dunkle, uopfyldte Ønskers
Skjød, kom med eet Slag til Virkelighed; og Virkeligheden
var overraskende. Man vidste paa det ene Sted
Intet om, hvad der foregik paa det andet, og dog skete
Alt som efter hemmelig Aftale. Morgenen kom, og see!
den nye Friheds Bygning stod for Alles Blik fuldfærdig
som Aladdins Slot, et Værk af Lampens Aander, tryllet
frem i en eneste Nat.

Men de frie Forfatningers Trylleslotte have desværre
ikke kunnet tilfredsstille Alle i Længden, navnlig
ikke dem, som uden at bekymre sig om Stilen i det
% --- p. 55 ---
\setcounter{footnote}{0}%
\opage{55}Hele, kun tænke paa at bringe deres private Huusholdningssager
i Orden. Der er ved de nye Samfundsbygninger
Adskilligt, som ikke er behageligt. De ere
prægtige nok at see til, men der er megen falsk Forgyldning;
de ere luftige nok, men der er saa megen
Trækvind; de ere rummelige, ja altfor rummelige. Der
er en Støi og Uro, en Tummel og Færdsel, en Skubben,
Trykken og Løben op ad Trapperne og ned ad Trapperne,
ikke alene ad Hovedtrapperne, men ad Bagtrapperne,
ja, hvad der især maa genere dem, der
holde af Fred i deres eget Kjøkken, endogsaa ad
Kjøkkentrapperne. Grunden er ikke fast, Pillerne ikke
stærke. Gulvet gynger; Loftet slaaer Revner; man befinder
sig som i et bestandigt Jordskjælv. Hvert Øjeblik
er der Noget der dumper; man gaaer i en dødelig
Angst for at faae en Bjælke, en Rigsdagsmand eller
en Minister i Hovedet. Man trøster sig med, at Bygningen
er saa ny, at det holder nok; en tarvelig Trøst!
Man har jo Exempler paa, at nye Bygninger falde,
endog før de ere blevne beboede.

Nei, da var der anderledes Orden i Statsbygningen
under de gamle enevældige Konger. Hver Stand stod
i sin Højde; man vidste at gjøre Forskjel mellem
Herskabet i Beletagen og Tjenerskabet i Messaninen.
Nu løber det Nederste sammen med det Øverste:
Fruen paa første Sal og Fruen i Kjælderen, Født og
Ufødt, Hat og Kaskjet, Uniform og Kofte er Hips som
Haps. Og saa spotter man over Rangforordningen og
gjør Væsen af Grundloven. Hvor bagvendt! Rangforordningen
er den egenlige Grundlov, den eneste
Lov, der kan adskille Klasserne og holde Orden i Samfundet:
naar Rangforordningen afskaffes, maa Staten
% --- p. 56 ---
\setcounter{footnote}{0}%
\opage{56}gaae til Grunde. Det er vigtigt for et civiliseret Samfund,
at det kjender sine store Mænd; men hvor er
det muligt at kjende dem, naar der ingen Rangforordning
er? Naar man i ældre Dage mødte en Mand med
et hvidt Baand og en Stjerne, betoges man af uvilkaarlig
Ærefrygt, man følte, at det var en stor Mand,
en af Staten fortjent Mand; hvori hans Fortjeneste
bestod, behøvede man ikke at vide, man havde jo
Rangforordningen. Selv et lille Kors var en Naadegave,
altsaa noget andet end en slet og ret Naturgave.
Distinctioner ere lige saa uundværlige for Meningen i
Selskabet, som Komma, Semicolon, Punktum og Udraabstegn
ere nødvendige for Meningen i en Bog. Men
de liberale Forfatninger oversee Distinctionen. Nuomstunder
gaae fortjente Mænd omkring og see ud
som andre Mænd; det giver Anledning til Confusioner.
Vil man kjende den fortjente Mand, skal man selv
have den Uleilighed at undersøge, hvori hans Fortjeneste
bestaaer, og saa kan Efterforskningens Udbytte endda
være tyndt nok. Paa mundtlige Overleveringer tør man
ikke stole: man maa gaae til skrevne Kilder, det
vil sige til Bladene; men Bladene? Det ene Blad
foreviser ham med en Straalekrands om Panden, det
andet med store Blækklatter i Ansigtet. Hvad skal
man tænke, hvad skal man troe? Maatte man ikke
ønske, at Folkerørets politiske Eventyr vilde ende med
Knald som et Eventyr, at Skæbnen eller Tilfældet eller
Himlen eller Lampen vilde sende os en statsklog
Nureddin, der kunde flytte Aladdinsslottet, dette Frihedens
Luftkastel med dets forlorne Prindsesse, over
til Afrika.

% --- p. 57 ---
\setcounter{footnote}{0}%
\opage{57}Slige Yttringer af uvilkaarlig Misstemning vilde
ganske vist have Noget for sig, dersom det Synspunkt,
hvorfra de udgaae, ikke var aldeles forfeilet. Men
Sagen er, at de frie Forfatninger, hvorpaa Nutiden
arbeider, ikke ere Frembringelser af Mængdens Vilkaarlighed,
men et Værk af historisk Nødvendighed.
De store Udviklingstrin i Menneskehedens Historie
kunne sammenlignes med Klasserne i en stor Skole.
Naar Eleverne have siddet en vis Tid i en af de lavere
Klasser, maae de enten flyttes op i en høiere eller gaae
ud af Skolen. De Nationer, som efterat have gjennemgaaet
den europæiske Civilisations Mellemstadier i det
17de og 18de Aarhundrede, endnu ikke ere modne til
at tage fat paa Frihedens Gjerning, maae finde sig i
at gaae ud af Historien. Der staaer over Scenen i
vort kongelige Theater en Indskrift: \emph{Ei blot til Lyst};
denne Indskrift bør staae med store gyldne Bogstaver
over vor ny Samfundsbygning.

Et Skuespilhuus, dets Stiil og Udseende være for
Øvrigt nok saa fortræffeligt, er ikke opført for dets egen
Skyld, men for Skuespillets Skyld. Var der ingen
Skuespildigtere og ingen Skuespillere, saa vilde den
prægtige Bygning staae som et tomt Monument, et
Minde om Savn og aandelig Afmagt. Dette gjælder i
endnu høiere Forstand om den paa Frihed anlagte
Forfatning. Er der ingen Frihedsideer og ingen Mænd,
som forstaae at bringe dem frem, saa staaer Bygningen
der som et Gudehuus uden Guder og seer haanlig ned
paa en vanartet Slægt, hvis aandelige Marv er fortæret.
Den frie Forfatning er til for Frihedens Skyld; den er
et Frihedens Grundlag med ydre Betingelser, men ikke
Friheden selv. En fri Forfatning behøver ikke altid
% --- p. 58 ---
\setcounter{footnote}{0}%
\opage{58}at kjøbes dyrt, den kan jo faaes som en Foræring: vi
takke Kong Frederik VII for hans Gave. Men Friheden
selv kan ikke faaes i Foræring; den maa kjøbes til
den dyre Priis, den maa erhverves i Kamp, om end i
fredelig, ublodig Kamp.

At Frihed maa erhverves, ligger i Sagens Natur;
thi hvad er Frihed? Frihedens Første er Selvbestemmelse,
og Selvbestemmelse er Virksomhed, i inderste
og dybeste Forstand Selvvirksomhed. Det Menneske,
der ikke føler Evne og Kraft til Selvbestemmelse, er
ufrit, hvor lykkelige dets ydre Vilkaar end ere, og det
Samme gjælder om et heelt Folk. Et Folk kan have
en fri Forfatning og endda befinde sig i en Tilstand
af Ufrihed, enten fordi det endnu ikke har lært at
bruge Friheden eller, hvad der omtrent kommer ud
paa det Samme, fordi det misbruger den. Til at gjøre
fornuftig Brug af Frihedens, særlig af Folkefrihedens
ydre Betingelser, udkræves der to Ting: Kræfter og
Indsigt. Kræfter uden Indsigt give kun Røre, og det
blinde, tankeløse Røre gjør Folket fortumlet. Indsigt
uden Kræfter give kun Doctriner. og tomme Doctriner % print: „Doctriner.“ (a stop before a lower-case word; misprint for a comma)
ere afmægtige overfor Virkeligheden.

Naar Frihedskriserne begynde, vaagner Folke-Aanden
og spørger om Kræfter. Der er dem, der
præke og sige, at „en Folkets Aand“ ikke behøver at
lede om Kræfter, eftersom den selv har Kræfterne i
sig, og kan, om vi end ere nok saa afkræftede, komme
og give os alle de Kræfter, vi behøve, hvad Dag den
vil; det er en Brouten med Aand og Ord, som vildleder
istedenfor at veilede. Skulde Aanden bringe
Kræfterne med sig fra Himlen, saa maatte det jo være
overnaturlige Kræfter, og Folkelivets Opsving beroe
% --- p. 59 ---
\setcounter{footnote}{0}%
\opage{59}paa et Mirakel. Det er Overtro. I en Handelskrise
siger Kjøbmanden: jeg maa have Penge, hvor jeg saa
skal faae dem fra, ellers gaaer jeg fallit. I en Frihedskrise
siger Aanden: jeg maa have Kræfter, hvor jeg
saa skal faae dem fra, ellers er det Hele forgjæves.
Det gaaer i denne Henseende lige saa naturlig til i en
Frihedskrise som i en Pengekrise. Dog er her unegtelig
ogsaa en Forskjel. Kapitalerne kunne skjule sig
saa snildt, at Kjøbmanden ikke er i Stand til at lede
dem op; men Kræfterne kunne ikke saaledes skjule sig
for Aanden; den skal nok vide at finde dem, hvor de
end ere. Det er ikke altid sagt, at de findes i Samfundets
øverste Lag; der er en Civilisation, som tærer
paa Kræfterne, der er en Kulturens Solbrand, som
udtørrer Sindet, saa at Livsspirerne visne. Finder
Aanden de øvre Lag forbrændte, arbeider den sig
dybere ned og finder Kræfterne i de lavere. Folke-Aanden
gjør ingen Mirakler, den skaber ikke af Intet;
men den bruger det Givne: dens Kraft bestaaer i at
vække de slumrende Kræfter og forhøie dem ved Beaandelsens
Tilvæxt.

Her er et Synspunkt, hvorfra Ordførerne for den
formelle Intelligens mene at kunne slaae et Slag til
Ære for hvad de forstaae ved Oplysning og Dannelse.
Skal Aanden, sige de, søge sine Kræfter i de nederste
Samfundslag, da maa den tage til Takke med raae
Kræfter; hvorledes forholder det sig nu med Kræfterne
og Raaheden? Antager man, at kun de raae Kræfter
ere usvækkede, saa maa man for Alting søge at bevare
Raaheden for at bevare Kræfterne. Er Folkefriheden
betinget af Kræfter, og Kræfterne betingede af Raaheden,
da maa Frihedstiden nødvendigviis blive en
% --- p. 60 ---
\setcounter{footnote}{0}%
\opage{60}Raahedens Tid. Antager man, at Kraften kan skilles
fra Raaheden, saa at Kræfterne kunne bevares, efter
at Raaheden er forsvundet, hvorledes vil man da kunne
lade haant om Civilisationens høieste Frembringelser,
Kunst, Poesi og Videnskab? Skal Folkefrihedens Tid
være en Oplysningens og Dannelsens Tid, da er Halvoplysningens
mythologiske Skumring ikke tilstrækkelig til
at fortrænge den raa Luft: Intelligensens Sol maa staae
lysende paa den klare Himmel. Men hvis man indrømmer,
at Kræfterne ikke svækkes, men forædles ved Oplysning,
saa maae de ædleste Kræfter jo søges i Samfundets
øverste, solbeskinnede Lag, og ingenlunde nederst i
Raahedens Mørke.

Det er den bekjendte Strid om Kulturintelligens
og Folkeoplysning, vi her berøre. At Kunst, Poesi og
Videnskab mægtig bidrage til at vække og klare Aandslivet,
og at Frihedskræfterne for saa vidt snarere maae
søges i de øvre Lag end i de nedre, er en Kjendsgjerning,
som Ingen, der kjender lidt til Historien, vil
kunne negte. Men at Kulturen ved Siden heraf udklækker
en Sværm af Intelligensinfusorier, der ræsonnere
paa Overfladen og tære paa Kræfterne, vil heller Ingen
kunne negte. Overfladisk Ræsonneren gaaer fra Tvivl
til Tvivl, og den endeløse Tvivl svækker Handlekraften.
Hvad gavner det en civiliseret Nutid at kjende alle
Fortidens Dyder og Lyder, dersom den ikke kjender
og forstaaer sig selv; og hvad hjælper det, om den
ogsaa betragter sig selv i Intelligensens Speil, naar det
bliver ved den golde Betragtning. Lys uden Kræfter
er et dødt Skin; men hvor Kræfterne ere, kan Aanden
finde dem og tænde Lys, endog i det dybe Mørke.
% --- p. 61 ---
\setcounter{footnote}{0}%
\opage{61}Sandhedens Ord siger ikke, at Lyset er Menneskets
Liv, men det siger, at Livet er Menneskets Lys.

Hvor stærke Kræfter afføde klare Tanker, fremkommer
der en Dannelse, hvorved det Fællesmenneskelige
og det for hvert dygtigt Menneske Eiendommelige
udvikles. I det folkelige Samfund vil Frihedsudviklingen
da udvise en Mangfoldighed af Dannelsestrin og Dannelseskredse,
og Dannelsen i enhver Kreds vil nøiagtig svare
til de Arbeider, der ere de særlige Klasser anviste.
Det er ganske sandt, at de haardføre Mænd, der færdes
paa Land og Hav, anderledes forstaae at haandtere
Virkeligheden end de, der sidde i et Studerekammer
og sysle med Tankerne; men derfor bør det dog ikke
glemmes, at der ogsaa er et Tankernes Arbeide, at
mangt et Lys, som har straalet over Samfundet og sat
en Verden af nye Kræfter i Røre, er blevet tændt i
det stille Kammer. Hovedsagen er, at de forskjellige
Arter af Dannelse indgaae i den rette Vexelvirkning,
idet de alle, hver med sit Maal af Kræfter og Lys,
træde i Livsudviklingens og Frihedens Tjeneste. Kjernen
i al Dannelse er Selvdannelse, Sindsdannelse, den Dannelse,
som lærer den Enkelte, at han i Sandhed kun
virker for sig selv ved at virke for det Hele. Paa
disse Betingelser kan Folke-Aanden røre sig, og Frihedens
Værk fremmes.

Kræfterne kunne virke uden Aand; men Aanden
kan ikke virke uden Kræfter. Hvor Kræfterne virke
uden Aand, vil Samfundet befinde sig snart i døsig
Orden, snart i uhyggelig Gjæring. Enten vil Virksomheden
være fordeelt saaledes, at hver Stand passer
Sit, Regjeringen Sit og Folket Sit; her mangler Vexelvirkning.
Eller de modsatte Bestræbelser ville geraade
% --- p. 62 ---
\setcounter{footnote}{0}%
\opage{62}i Strid med hverandre, det ene Parti vil bekæmpe det
andet og tabe det fælles Maal af Syne: her mangler
harmonisk Samvirken. Kjendemærket paa, at det er
Folke-Aanden, der sætter Kræfterne i Røre, er dette
dobbelte: at alle Bestræbelser virke sammen til eet
Maal, og at Alle, Høie og Lave, ere sig bevidste, at
de maae deltage i Virksomheden. Det sidste Punkt
bliver især af Vigtighed, naar den frie Forfatning har
afløst Enevælden.

I en Samfundstilstand, hvor Aanden endnu er
halvt slumrende, falder det naturligt, at kun nogle
Individer af enkelte Klasser give sig af med offentlige
Anliggender, medens den store Mængde holder sig til
det Private og lader Regjeringen sørge for Staten. Det
er en Synsmaade, som behager Regjeringen; den holder
ikke af de urolige Hoveder. Øvrigheden søger som
Folkets Formynder at slaae den Grundsætning fast, at
Undersaatterne have at underkaste sig deres Foranstaltninger
og for Øvrigt holde sig rolige.

Det ligger i Statens Natur at holde sig som
et Bestaaende; de Styrende holde desaarsag skarpt
Øie med de Urolige, men see naadigt og mildt
til de Stille i Landet. Statsmændene i Oldtiden vare
i denne Henseende allerede opmærksomme paa Forskjellen
mellem Arbeiderne i Stæderne og Landbefolkningen.
„Arbeideren“, siger Aristoteles, „er en Mand
af ringe Dyd; og kun i ufuldkomne Stater kan det
tilstedes Arbeiderne at deeltage i Lovgivning og Styrelse.
En Befolkning, der fornemmelig sysselsætter sig
med Agerdyrkning og Kvægavl, er i Reglen tilfreds,
naar den blot kan faae Ro til at blive ved sit Arbeide,
og lader sig gjerne nøie med en begrændset Andeel i
% --- p. 63 ---
\setcounter{footnote}{0}%
\opage{63}Statsstyrelsen. Langt uroligere er Massen af Haandværkere,
Handlende og Saadanne, der leve af Fortjenesten:
deres Næringsvei virker skadeligere paa
Charakteren, og i Staden, hvor de findes sammenhobede,
ere de tilbøielige til at støie i Forsamlingerne.“

Naar det skal blive til Alvor med Brugen af den
frie Forfatning, altsaa med Friheden, kunne vi --- med al
Agtelse for Ro og Orden --- ikke ubetinget lovprise de
Stille i Landet. Hvad er det for Folk? Skikkelige Folk;
i private, personlige Forhold hæderlige, agtværdige,
nyttige Medlemmer af Samfundet. Men deres Hang til
Fred og Iver for at passe deres Eget er saa overveiende,
at virksom Deeltagelse i det offentlige Liv er
dem næsten modbydelig. Det er forgjæves at minde
dem om Indskriften: Ei blot til Lyst, thi de have ikke
engang Lyst. Heller ikke nytter det stort at foreholde
dem den fælles Pligt, som Forfatningen paalægger os i
Frihedens Navn, thi hvad Friheden angaaer, have de
deres egen Forklaring: „I en fri Forfatning maa Enhver,
som ikke vil med i den offentlige Tummel, jo
have Frihed til at holde sig tilbage. Der er, sige de,
Folk nok, som have Lyst til at agere paa den politiske
Skueplads. Men der er ingen rigtig Fornøielse ved at
være Skuespiller, naar der ingen Tilskuere er. Vi, der
ikke kunne agere, tjene det Offentlige ved at være
Tilskuere“. De Stille i Landet ere Tilskuere. Hvilken
Glæde have disse Tilskuere nu af de folkepolitiske
Optrin? Glæden er ikke stor. De, der sidde i første
Parket og tydelig kunne see Alt gjennem Kritikens
Kikkert, gjøre ikke Andet end ærgre sig. Gaaer Stykket
flot, saa ærgre de sig, thi Friheden bliver kaad og
kaster dem Sand i Øinene. Gaaer Stykket trevent, saa
% --- p. 64 ---
\setcounter{footnote}{0}%
\opage{64}ærgre de sig ogsaa: der burde jo pibes; men tys! ---
for Rolighedens Skyld tør man neppe engang nyse,
endnu mindre pibe. De, der staae saa langt tilbage,
at de knap kunne skimte, hvad der foregaaer paa
Scenen, ansee det Hele for Noget, som i Grunden ikke
kommer dem ved; det Eneste, der leilighedsviis kan give
deres Rolighed et Stød, er Rygter om nogle forskrækkelige
Ting, som skulle foregaae bag Coulisserne.

Da de Stille i Landet ikke udtale sig offentlig,
kunne vi naturligviis ikke have nogen tilforladelig
Kundskab om deres Meninger i det Enkelte; men det
kan ogsaa være det Samme. Hvad det her kommer
an paa, er alene den Kjendsgjerning, at de ikke ansee
en virksom Deeltagelse i de store Samfundsanliggender
for en almindelig Borgerpligt. Kunde vi ligefrem henvende
Spørgsmaal til dem, vilde de uden Tvivl svare:
„Der kommer ikke Noget ud af al denne politiske Rumsteren.
Det kan ikke være hver Mands Pligt at forstaae
sig paa Regjeringssager. Vi forstaae os ikke
paa Regjeringssager og holde os for gode til at gjøre
Sjouerarbejde hos dem, som for at kunne bruge os,
maae begynde med at bilde os ind, at vi ere store
Politikere. En fornuftig Mand befatter sig aldrig med
det, han ikke forstaaer.“ --- Et saadant Svar vilde være
fyldestgjørende under Enevælden, men under en fri
Forfatning er det ingenlunde fyldestgjørende. Dersom
et Menneske uden videre kunde slippe for at befatte
sig med det, han ikke forstaaer, saa vilde det jo være
snildt at lægge Vind paa at forstaae saa Lidt, som
muligt. Der kan være Meget, vi ikke forstaae, og som
vi dog maae lære at forstaae, naar det vedkommer os
og vor Gjerning. Vistnok kan hvert andet Menneske i
% --- p. 65 ---
\setcounter{footnote}{0}%
\opage{65}en Befolkning ikke offre Tid og Flid paa at kjende
Regjeringssagerne i det Enkelte; men rette vi det
Spørgsmaal til de Tilbageholdne, om de da ville opgive
deres Valgret? faae vi uden Tvivl det Svar: „paa ingen
Maade! hvorfor skulde vi ikke have samme Ret som
alle de Andre.“ --- Men hvorledes kunne I sætte Priis
paa en Rettighed, som I dog, efter egen Tilstaaelse,
ikke forstaae at bruge? --- „Vi forstaae det ligesaa godt
som alle de Andre.“ --- Saa siger os da, hvorledes I
bære Eder ad paa en Valgdag? --- „Undertiden møde
vi, undertiden møde vi ikke. Men hver Gang vi møde,
komme vi strax ind i Sagen. Vi høre Candidaternes
belærende Udtalelser og mærke os de oplysende Svar,
naar de blive interpellerede. Vel kunne vi ikke altid
saa lige hitte Rede i det Enkelte, men Aanden er jo
ogsaa det Vigtigste.“ --- Hvordan bære I Eder ad med
at opfatte Aanden? --- „Det er den letteste Sag af
Verden. Man behøver for den Sags Skyld ikke engang
at forstaae, hvad det er, der siges, men blot at lægge
Mærke til, hvorledes det Sagte virker paa Forsamlingen.
Virker det Hør, Hør! og Ja, Ja! saa er Aanden der;
virker det Hys, Hys! Nei, Nei! saa er Aanden der ikke.
Finde vi nu efter omhyggelig Overveielse, at Aanden
er for en Høiremand, saa vælge vi naturligviis en Høiremand;
men vil Aanden absolut have en Venstremand,
nu vel! saa kunne vi jo ogsaa vælge en Venstremand.“

--- Altsaa: paa denne Maade tjene I Friheden ved
samvittighedsfuldt! --- at bruge Eders Valgret.

Valgretten er et farligt Vaaben; at betroe et farligt
Vaaben til Folk, som ikke engang gide erkyndige
sig om dets rette Brug, er ubetænksomt. Der klages
over, at vi have for mange Examener; man kunde dog
% --- p. 66 ---
\setcounter{footnote}{0}%
\opage{66}fristes til at foreslaae een Examen endnu, nemlig en
almindelig Vælgerexamen. Hvor mange Commissioner
der i den Anledning vilde blive at nedsætte, hvorledes
de skulde sammensættes, af hvilke og hvor mange
Medlemmer en saadan Examenscommission helst burde
bestaae, hvor store Diæter der i Forhold til Kornpriserne
eller Tobaksskatten eller den stigende Udførsel
af Æg kunde tilkomme ethvert Medlem i Commissionen,
det er praktiske Hensyn, vi ikke her kunne indlade os
paa; det er os nok at pege paa Tidens Trang, dens
dybtfølte Trang til en Vælgerexamen. Ved en saadan
Examen burde det da ikke nærmest komme an paa
Candidatens udvortes Stilling, om han er menig eller
fornem, fattig eller rig --- Persons Anseelse gjælder
ikke ved en offentlig Examen. --- Udfaldet beroer paa
Candidatens Svar, og at der holdes Orden med Charaktererne.
Kan en Candidat gjøre tilstrækkelig Rede for
sin Synsmaade og de Grunde, der lede ham, naar han
vælger, da bør det ikke komme ham til Skade, at han
for Resten er en uformuende Mand; selv en lille
Mellemregning med Fattigvæsenet skal ikke forhindre
ham i at faae Meget godt med Kryds. Melder sig
derimod en Candidat uden ordentlig Forberedelse, væver
han ud og ind om Conservative og Liberale, om Stilstandsparti
og Fremskridtsparti saa længe, indtil han
bogstavelig ikke kan skille Høire fra Venstre: da bør
Commissionen ikke lade sig blænde af hans ydre fordeelagtige
Stilling; selv om han eiede ti Fabriker og
tolv Herregaarde, han burde dog have en maadelig
Charakteer: jeg for min Deel vilde stemme for Nul.

Men --- indvender man maaskee --- et saadant Forslag
gaaer aldrig igjennem; det støder paa praktiske
% --- p. 67 ---
\setcounter{footnote}{0}%
\opage{67}Vanskeligheder. Kan nok være. Men derfor ville vi
dog ikke tabe Modet. Vi have i Virkeligheden, naar
vi see nærmere til, en Vælgerexamen. Tilværelsen selv
med sine Erfaringer er den store Examinator. Friheden
har sit Tugtens Riis i den strænge Nødvendighed; den
tugter os ved Følgerne af vore Handlinger. En fri
Mand, forstaaer sig, har Lov til at sove. Men i den
travle Tid, naar Sæden skal i Jorden eller naar Høsten
staaer paa, veed den frie Mand meget vel, hvilke Følger
det kan have at sove over sig. Hvad Agerdyrkning og
Industrie angaaer, har en fri Mand ganske vist Lov til
at staae stille og sige: jeg følger ikke med Tiden; hos
mig skal Alting blive ved det Gamle. Jeg handler som
min Bedstefader handlede, da der endnu hverken var
Dampskibe eller Jernveie. Jeg holder min Gaard i
samme Drift som mine Forfædre; de nymodens Opfindelser
komme ikke mig ved; jeg vil ikke ud af mine
gamle Vaner. See, det kan en fri Mand jo ganske vist
gjøre; thi for Afløsning af dovne Vaner vil den frie
Forfatning neppe udstede nogen Tvangslov; men Følgerne
udeblive ikke: Friheden vil tugte den frie Mand
med Nødvendighedens Riis. Saaledes, og ikke anderledes,
vil det i Tidens Løb gaae med Brugen af Valgretten;
man kan misbruge den og man kan forsømme
at bruge den, men Følgerne ville ikke udeblive.

De, der klage over, at det gaaer skævt med Friheden,
og ved deres Klager antyde, at de bedre vide,
hvorledes det burde gaae, tør ikke sidde hen som Tilskuere.
De maae frem paa Scenen, at Offentligheden
kan see det Skjoldmærke, de føre. Frihed er Virksomhed,
Selvvirksomhed, og Selvvirksomhed er nøie
beslægtet med Driftighed. Det driftige Parti i et Land
% --- p. 68 ---
\setcounter{footnote}{0}%
\opage{68}vil altid være et Hestehoved forud for det blot ræsonnerende.
Et Modparties Overgreb lade sig ikke afvise
ved blot at udpege dets Skygger og Skævheder, men
ved at organisere sig, slutte sig sammen og handle.
Ganske vist kan Friheden ogsaa fordre af sine Besluttede,
at de for en Tid staae stille, see Tingene an og lade
de iltre Modstandere løbe Linen ud. Men en saadan
Stillestaaen er da selv en Handling, og en saadan
Handling vidt forskjellig fra det at være enten muggen
Kritiker eller tankeløs Tilskuer.

At paakalde Folke-Aanden, medens man dog gjør
Alt for at kvæle den, er ganske vist en Modsigelse;
det er et sørgeligt Syn, naar et driftigt Fremskridtsparti
i forblindet Trods griber feil af Midlerne. Saalænge
der imidlertid er Tro paa Aand og Mod til at
kæmpe for Frihedens Sag, er der dog Haab; Feilgrebene
kunne jo rettes, og Udviklingen igjen komme ind i det
oprindelige Spor. Men de, der ere i Færd med at
opgive Aanden, de betænksomme Ordensmennesker, der
kun ere rolige, fordi de ere matte, og matte, ja dødtrætte,
fordi deres Livskapital er brugt op, fordi de
ikke have en eneste levedygtig og frugtbar Tanke
tilbage, hvad Haab er der for dem? Deres Løsen
er: at holde \emph{det gaaende}. Forstaaer sig, at
holde det gaaende kan ogsaa være Opgave for dem,
der ere i fuld Virksomhed med det Fremtidige. Men
naar det, at holde det gaaende, er Mathedens, Udlevethedens
og dog Seiglivethedens sidste Tilflugt, saa
falder der et Alderdomsskjær derover, som er uhyggeligt.
Der fortælles hos Snorre Sturlesøn om en Kong
An eller Øn, hvorledes han søgte at holde det gaaende.
Det hedder i Grundtvigs Oversættelse: „Efter Jørrun
% --- p. 69 ---
\setcounter{footnote}{0}%
\opage{69}blev hans Søn, An eller Øn, Konge i Sverrig; det var
en snild Mand og en stor Afgudsdyrker, men Krigen
lod han andre om, og havde Hjemmet kiærest. Til al
Uheld for ham var det anderledes fat med den danske
Kong Halvdan\dots{} Han gik til Sverrig med en Krigshær,
og sloges seierrig med Kong Øn, saa længe, til
han maatte tage Flugten ind i Vester-Gylland, og lade
Halvdan herske til sin Dødsdag i Upsal. Dengang
havde Øn alt regjeret i tyve Aar, eg Halvdan regjerede % print: „eg“ (misprint for og)
i andre tyve, men da han var lagt i Høien, kom Øn
tilbage til Upsal i sit tredsindstyvende Aar, og gjorde
strax et blodigt Offer, for at leve længe endnu. Det
var hans egen Søn, han offrede til Odin, og fik derfor
det Løfte, at han skulde blive ligesaa gammel til, som
han var, men Riget beholdt han heller ikke dette
Sæt længere end tyve Aar, thi saa kom Kong Fridlevs
Søn i Danmark, Ole Bravkarl, med sin Krigshær, slog
Øn baade sønder og sammen og nødte ham til atter
at tye til Vester-Gylland. Nu gik der igjen tyve Aar,
men saa myrdede Stærk-Odder hiin Gamle Ole Bravkarl,
og strax kom Øn tilbage til Upsal, og herskede
sine tyve Aar paa tredie Omgang. Han vilde dog leve
længer endda, og offrede Odin sin anden Søn, hvorfor
han fik Løfte paa at leve i ti Aar, og saa fremdeles ti
for hver Søn han gav Odin, dog paa det Vilkaar, at
han hvergang skulde give et Herred samme Nummer
som den Søn, han offrede. Nu blev han ved, og da
han offrede den syvende Søn, kunde han ikke længer
styde paa Benene, saa han maatte bæres i en Stoel,
men ikke desmindre offrede han den ottende med og
kunde i de følgende ti Aar ikke komme op af sin
% --- p. 70 ---
\setcounter{footnote}{0}%
\opage{70}Seng. Endda offrede han sin niende Søn, og maatte,
for at holde Livet, som et Pattebarn die af et Horn.
En eneste Søn havde han endnu, og ham vilde han
ogsaa have offret til Odin, med samt Upsal og Herredet
deromkring, som skulde hedt Tiende-Land, men det
formeente Svenskerne ham, saa der blev Intet af, og
An maatte finde sig i at døe, og blive høilagt ved
Upsal“.

Som der kan være en forunderlig Seiglivethed hos
gamle udlevede Folk, saaledes kan der ogsaa være en
næsten utrolig Seiglivethed i gamle, udlevede Tilstandsformer.


Men ikke blot Slapheden, Villieløsheden og Alderdomssvækkelsen
er en Tyngsel paa Friheden; ogsaa den
uvorne Driftighed, den selviske Paatrængenhed, kan
sætte den i Fare. --- En fremmed Magthaver skal en
Gang have sagt: hver Gang Danmark er ved at gaae
under, har Vorherre altid et lille Mirakel i Baghaanden.
Vi have ingen Ret til i dette Punkt at stole paa Mirakler,
men vi ere forpligtede til at stole paa Aandens
Bistand, agte paa Tidens Tegn og vedkjende os Frihedens
Fordring. Friske Kræfter, frugtbar Oplysning,
høisindet Mod og folkelig Dannelse ere ideale Magter,
med hvilke vi først og sidst maae søge Alliance. Det
er Allierede, som ei ville svigte i Farens Stund, naar
vi blot ikke svigte os selv. Men vi svigte os selv,
naar vi smaaligt kjævles om Magten; vi opgive os selv,
naar vi i lidenskabelig Forblindelse forvandle vor Forfatningsbygning
til et Spillehuus, hvor Partiaanden siger
va banque og sætter Friheden paa Spil.

Endnu staaer Danmarks gamle Huus, trods Rov og
Plyndring, paa sin egen Grund. Af Frihedstømmer
% --- p. 71 ---
\setcounter{footnote}{0}%
\opage{71}skal den lille Rest, der er tilbage, bygges op paa Ny.
Det bliver vel ikke et Aladdins Slot, men ere Arbeiderne
rede, og Bygmesterne enige om Stiil og Plan,
kan Huset reise sig med Ærens Krands, et Enighedsminde,
et Aandens Værk, en Folkeborg mod Tidens
Storme.

% --- p. 72 ---
\setcounter{footnote}{0}%
\lecture{V.} {Aristokratisme og Demokratisme.}
\opage{72}Den travle Nutid tør aldrig glemme, hvormeget den
skylder sin Forberedelsestid, sin hæderlig begravne
Fortid. Skulde den glemme det andetsteds; i denne
By vil den ikke glemme det: Fortidens Haand er her
udrakt og griber i det Nærværende. Ungdomsfrisk med
Fremtidshaab og Fremtidsblik kan Livet røre sig overalt,
i Borgernes Huse, i Arbeidernes Værksteder, paa
Torve og Gader, men over det rørige Nutidsliv breder
sig stille et Oldtidsskjær, et forklarende Skin fra det
Svundne: Viborg staaer paa historisk Grund med
gamle Tanker og med gamle Minder. Domkirken hist
med de runde Buer og de høie Taarne gjemmer paa
Minderne; Skolen overfor i dens Paasyn samler paa
Tankerne, og begge, Kirke og Skole, ere enige om at
Enhver, der tilgavns vil forstaae den forviklede Nutid,
ikke bør forsømme at leve sig ind i den store erfaringsrige
Fortid.

Fortid og Nutid, det Gamle og det Nye befinde
sig i uophørlig Brydning; en ustandselig Higen efter
Forandring er idelig ifærd med at omdanne det Bestaaende.
Det Gamle, hedder det, skal falde og det
% --- p. 73 ---
\setcounter{footnote}{0}%
\opage{73}Nye reise sig paa dets Ruiner, kun Fremskridt er Liv,
Stilstand er Død. I slige politiske Slagord er der
noget Sandt, men ogsaa noget Usandt, Skævt og Vildledende.
Spørgsmaalet er, hvorledes det i Samfundet
egenlig forholder sig med Overgangen fra det Gamle
til det Nye.

Naar et Huus bliver gammelt og forfalder, reparerer
man det. Man kan da bygge Adskilligt op,
tage Noget bort og øge Noget til; men, naar dette skeer
gjentagne Gange, kan ogsaa Bygningen efterhaanden blive
saa forandret, at den oprindelige Stiil ikke længere er til
at kjende: den deelvise Omdannelse ender med en barok
Sammenstykning. Den engelske Statsbygning med de
Reformer, den i Tidernes Løb er undergaaet, har Adskilligt,
som minder herom. En engelsk Parlamentstaler
yttrede engang ved Anledningen af den Mængde
Ændringer, der blev foretaget med en Reformbill:
denne Bill minder mig om Bomuldsstrømpen, der blev
stoppet saalænge med Silke, indtil den blev en Silkestrømpe.


Den modsatte Fremgangsmaade kommer til Anvendelse
for Exempel i Frankrig, hvor man rask væk
river det Gamle ned for at faae Plads til det Nye.
Naar det Bestaaende ikke længere tilfredsstiller, gjør
man Revolution. I Løbet af et Par Dage kan man,
naar man ret kommer i Øvelse, rive en heel Statsbygning
omkuld og indføre en splinterny Forfatning.

I første Tilfælde --- efter engelsk Mønster --- er
det det Gamle, der deelviis forvansker det Nye; i
sidste Tilfælde --- efter fransk Mønster --- er det det
Nye, der ødelægger det Gamle. Men just derved at
% --- p. 74 ---
\setcounter{footnote}{0}%
\opage{74}det Gamle ødelægges, undergraves Grundvolden for
det Nye. Det Sidste er værre end det Første.

Forholdet mellem det Gamle og det Nye stiller
sig heelt forskjelligt, eftersom det er en blot mekanisk
eller en organisk, en livløs eller en levende Bygning,
man har med at gjøre.

Den menneskelige Samfundsbygning er ikke en
død mekanisk Sammensætning, men en levende besjælet
Organisme.

I de levende Organismer, Planter og Dyr, foregaaer
der en uafbrudt Omdannelse af alle, selv de
mindste Dele. Der arbeides i Celler og Væv og utallige
Værksteder paa Omdannelsen, saaledes at Forandringen
bliver gjennemgaaende. Grundformen, hvad
vi kunde kalde Stilen, vedligeholdes. Periodeviis kan
vel ogsaa Grundformen ændres, --- der er, som bekjendt,
Planter og Dyr, der i forskjellige Livsperioder
have et aldeles forskjelligt Udseende --- men selv hvor
Forandringen synes at skee pludseligt, er den dog
længe, alsidig og paa indvortes Maade forberedt, saa
at der i alle Led og Dele viser sig den inderligste
Sammenhæng mellem det Gamle og det Nye.

Hvad Naturen har foretegnet os i Exempler af de
levende Organismer, skulde nu i høiere Maal og paa
høiere Maade gjentage sig i det menneskelige Samfund.
Men netop fordi Udviklingens Betingelser ere
af høiere Art, gjennemgaaer Udviklingen selv forskjellige
Svingninger.

Samfundslivet med dets mange Indretninger skyder
ikke frem som Skud af et spirende Frø, voxer ikke en
stille Væxt efter nødvendige uforanderlige Love; thi
Udviklingen har et dobbelt Grundlag, et naturligt og
% --- p. 75 ---
\setcounter{footnote}{0}%
\opage{75}et aandeligt. Samfundslovene, hvor naturnødvendige
de end kunne være, ere Frihedslove, med Præg af
Vilkaarlighed; Udviklingen er en Kamp mellem mange
Interesser. Bliver Kampen saa heftig at Samfundslagene
vælte sig over hverandre, at Klasse staaer op
imod Klasse, Parti mod Parti, da er den jevne Udvikling
afbrudt, og Spørgsmaalet bliver, hvorledes det
Nye vil stille sig til det Gamle. Ustraffet gjør intet
Folk et Spring i sin historiske Udvikling. Overleveringer,
Vaner, alle Eftervirkninger fra Fortiden fortsætte
sig med en Slutningshastighed, der er bestemmende
ikke blot for Bevægelsens Fart, men, hvorhen
det Nye end forresten stræber, ogsaa for dens Retning.
Overgangen fra en Forfatning til en anden er
altid mislig, og det i samme Grad, som de tilsvarende
Betingelser ere uligeartede. Et Folk, der i Aarhundreder
har levet under en monarkisk Regjering, er ilde
tjent med at blive kastet hovedkulds ind i en demokratisk
Republik. Efter de Erfaringer, Folk og Statsmænd
maae have gjort i sidste Halvdeel af forrige og
Begyndelsen af dette Aarhundrede, skulde man antage,
at den Tid dog nu maatte være forbi, da man bildte
sig ind at fuldfærdige Statstheorier kunde erholdes i
en politisk Boutik, ligesom færdigsyede Klæder i et
Skræddermagazin. Ikke desto mindre vil Parti-Aanden,
naar den finder sin Regning derved, uden Hensyn til
det paagjældende Folks organiske Livsbetingelser, naturligviis
ikke undlade at røre Trommen til indtrængende
Anbefaling af Demokratiet. Demokratiet er et
Slagord, der altid passer i Mængdens Kram og hvorfor
skulle Demagogerne saa ikke føre det i Munden?
Naar det hedder om en Mand, at han er Aristokrat,
% --- p. 76 ---
\setcounter{footnote}{0}%
\opage{76}saa veed Hoben strax, at det er en ufolkelig Mand,
en Mand, der overseer og ringeagter Folket. Idet det
Aristokratiske paa denne Maade bliver eenstydigt med
det Ufolkelige, bliver det Demokratiske naturligviis
eenstydigt med det Folkelige. Demokraten af det rene
Vand, den ægte Demokrat er Folkets Ven; det troer
Mængden og tager de ved denne Anledning brugelige
Slagord for gode Vare. Naturligt nok, naar den ikke
kjender den Blanding af Sandt og Falsk, der ligger i
saadanne Slagord, naar den med andre Ord ikke er
istand til at raadføre sig med Historien og føle
dens store Lærdom. Det er unegtelig sandt, at der
gives en Form af Aristokrati --- være sig Fødselsaristokrati
eller Embedsaristokrati --- som er uforenelig
med Folkelivets Fordringer. Men det Samme, det
Selvsamme gjælder om Demokratiet. Ogsaa i Demokratiet,
navnlig det radikale Demokrati, er der noget
Ufolkeligt, noget Udbrugt og Forslidt, der maa omdannes.
Vi kunne bedst overbevise os herom ved at
see paa Grundsætningerne.

Hvad er det vel for en Grundsætning, hvorpaa
Aristokratiet staaer fast, naar vi see bort fra de historiske
Misbrug? det er bogstavelig og efter Ordlyden den,
at de \emph{Bedste skulle regjere}. Dersom det Demokratiske
nu var, hvad det udgiver sig for, det Aristokratiske
stik modsat, saa maatte Demokratiets Grundsætning
jo aabenbart være den, at \emph{de Sletteste
skulle regjere}. Men en saadan Grundsætning vil
dog vel Ingen vedkjende sig. Vil Demokraten erklære,
at det hverken er de Bedste ei heller de Sletteste,
der skulle regjere, men Mængden, og at i Mængden
Alle ere lige gode: saa maae vi spørge ham, hvad
% --- p. 77 ---
\setcounter{footnote}{0}%
\opage{77}det da er Mængden seer paa, naar den vælger sig en
Fører. Dersom det virkelig er sandt, at Alle ere lige
gode, saa er det jo ikke nødvendigt at gjøre særegne
Ophævelser med Valget, man kan vælge iblinde eller
ved Lodtrækning eller hvordan man vil. Men endskjøndt
Mængden veed eller dog bilder sig ind at vide,
hvad der bør gjøres, saa er det dog, naar det kommer
til Stykket, ikke Enhver i Mængden, der veed, hvorledes
det bør gjøres, ja, seer man nærmere til, er det
kun saare Faa, der have nogenlunde klar Forestilling
om, hvad det egenlig er, der bør gjøres. Mængden
har af Instinkt en Fornemmelse af, at endskjøndt Alle
saadan i Almindelighed ere lige gode, saa ere dog
Alle, naar det virkelig gjælder om at udføre Noget,
ikke lige gode. Det er Mængden af Vigtighed at faae
dygtige Førere, thi uden Førere kan den Intet udrette.
Naar den har fundet en Fører, som ret er
efter dens Sind, hvad gjør den saa? saa følger den
ham blindthen; den lyder ham ubetinget, troer paa
ham --- at sige saa længe det varer --- og hvorfor?
naturligviis fordi den anseer ham for den Dueligste,
den Bedste. Den Bedste skal altsaa føre an, den
Bedste skal have Magten, den Bedste skal regjere.
Men at den Bedste skal regjere er jo en aristokratisk
Grundsætning. Mængden er altsaa uden at vide deraf
aristokratisk stemt.

Det er uberegneligt, hvad Folkelivet vilde vinde,
dersom denne simple, for Alle indlysende Grundsætning,
at de Bedste skulle regjere, fik Lov til at trænge
igjennem.

Lad os betragte Modsætningen fra en anden Side
og see ikke blot paa Personerne, men paa Sagen, ikke
% --- p. 78 ---
\setcounter{footnote}{0}%
\opage{78}alene paa de Bedste, men paa det Bedste, det Sande,
det Gode, det Rette, kort: det som er hele Samfundet
tjenligt. Hvorledes stiller Forholdet sig nu mellem
det Aristokratiske og det Demokratiske? Aristokraten
forudsætter, at det statsborgerlige Samfund, hvorledes
Forfatningen end er, underligger bestemte Udviklingslove,
at der er noget i og for sig Skadeligt og noget i
og for sig Gavnligt, og at Forskjellen mellem Skadeligt
og Gavnligt ikke kan afgjøres ved løst Skjøn eller
øieblikkelig Fleertalsstemning, men kun ved Sagundersøgelse
og Indsigt. Hans Valgsprog er, at det Sandeste
altid er det Bedste. Dersom han bildte sig ind,
at man blot behøver at være født med en Herregaard
eller et Fideicommis for at have Privilegium paa at
vide, hvad der for Samfundet er det Bedste, saa var
hans Aristokratisme ganske vist en Indbildning; men
den sande Aristokrat har en Agtelse for Sandhed og
Ret, han ikke kan opgive, uden at opgive sig selv.
--- --- Lad os saa see, hvorledes Demokraten bærer sig ad.
For at udfinde hvad der er til Samfundets Bedste, behøver
han ikke at plage sig med møisommelig Granskning;
han har blot at agte paa de Flestes Mening;
„Alles Mening“, og at vende sig til Mængden.
Mængden er hans Orakel, Mængden alene veed hvad
der er det Bedste. Ingen Støver har finere Næse til
at opspore Vildtet end Demokraten til at opspore
Mængdens Mening. Den Skarpsindighed, han herved
lægger for Dagen, har billigt Krav paa vor Beundring.
For det Første har han selv en Mening; men at vedkjende
sig den og træde offenlig frem med en stakkels
individuel og privat Mening, hvad kunde det føre til?
Den vilde jo blive underkastet en skarp Kritik; thi
% --- p. 79 ---
\setcounter{footnote}{0}%
\opage{79}med en individuel Mening gjør man ingen Complimenter;
ja Almeenheden var vel endog istand til at fordre
af ham det Uoverkommelige, at han for Alvor skulde
indlade sig paa at begrunde sin Mening. Nei Tak!
saa dum er han ikke. Han vælger en mere praktisk
Fremgangsmaade. Idet han træder frem for Mængden,
forstaaer han at tale saaledes, at det ved et akustisk
Bedrag, --- thi Demokraten er en øvet Bugtaler ---
kommer til at lyde, som om det var Mængdens egen
Mening, han udtalte. Mængden, der nu først rigtig
faaer at vide, at den har en Mening, lyser øieblikkelig
Meningen i Kuld og Kjøn, og, ligesaa glad overrasket
ved at høre sin egen Mening, som Konen i „Gjenboerne“
ved at see sine egne „Vindever“, jubler i Sky
af Henrykkelse. Nu staaer kun tilbage, hvad der er
den letteste Sag af Verden, at slaae Mængden og
Folket i Hartkorn, og hvilken Forvandling! Demokratens
Mening er bleven til Mængdens Mening, Mængdens
Mening til Folkets Mening, og Folkets Mening behøver
ingen Begrundelse.

Under disse Forudsætninger kan Demokraten forsvare
sig for hvilken Domstol, det skal være. Det er
forgjæves, at de Sagkyndige fremkomme med Indsigelser;
han svarer blot: Saadan vil Folket have det
og man maa bøie sig for Folkets Villie. Det er forgjæves
at bedre Underrettede foreslaae Noget, som
i sig selv er langt fornuftigere; kan Demokraten ikke
værge sig mod de slaaende Fornuftgrunde, der tale
for Modpartens Forslag, saa siger han: Som Doctrin
kan Sagen jo nok tage sig ud, især naar den
sees igjennem lærde Briller; men Folket, der helst seer
Tingen med sine egne naturlige Øine, har nu engang
% --- p. 80 ---
\setcounter{footnote}{0}%
\opage{80}faaet Mistillid til denne Sag, og at byde Folket Noget,
hvortil det har Mistillid, er upraktisk. Man seer her
det Fine i Fremgangmaaden: først vækker han det % print: „Fremgangmaaden“ (for Fremgangsmaaden)
lidet sagkyndige Folks Mistillid til den Sag, han ikke
vil have frem, og saa anvender han den vakte Mistillid
til at faae Sagen afviist, før Folket endnu er
kommet til at indsee sin Vildfarelse. Saaledes maa
den ægte Demokrat bære sig ad, naar han ret skal
handle i Overeensstemmelse med den demokratiske
Grundsætning.

Men, spørger man, ere alle Demokrater da virkelig
saadanne Bedragere? Vi svare nei, langtfra, de Fleste
iblandt dem ere uden Tvivl ærlige velmenende Personer,
der staae paa Frihedens Grund og virkelig
bilde sig ind at tjene den folkelige Sag. Hemmeligheden
er, at de lægge noget Andet og Mere ind i Grundsætningen,
end den egenlig indeholder. Grundsætningen
er, at Mængden veed hvad der er det Bedste, og
at Fleertalsvillien desaarsag bør være den eneste
raadende. Det heri skjulte Bedrag kommer først for
Dagen, naar man alvorlig gjør sig rede for \emph{Sandhed}
og Fleertal.

At komme til Sandhedserkjendelse --- dette gjælder
ogsaa om den politiske Sandhed --- er og har fra Arilds
Tid været en saare vanskelig Sag. Ikke Mængden, men
kun ganske Faa blev det givet at see saa klare Glimt
af Sandhedens Lys, at de fandt en Vei igjennnem Vildfarelsens % print: „igjennnem“ (misprint)
Udørken; endnu Færre blev det givet at trænge
igjennem og skaffe Sandheden Indgang. I Oldtiden blev
Sandheden anseet for en Hemmelighed, der kun turde betroes
til Indviede; at meddele den store Hob Sandheden
var at kaste Perler for Sviin, et livsfarligt Foretagende,
% --- p. 81 ---
\setcounter{footnote}{0}%
\opage{81}en Krænkelse af det Hellige. Det var en falsk Aristokratisme.
Vel er Sandhedens Opdagelse særlig forbeholdt
de Begavede, men den opdagede Sandhed vil
derfor ikke indespærres inden ægyptiske Tempelmure:
Sandheden er Lys og vil lyse for Alle. Det gaaer med
Livssandheden som med Skjønheden; der gives en naturlig
Sands for det Skjønne, der gives ogsaa en naturlig
og oprindelig Sands for det Sande. Det menneskelige
Øie maa efterhaanden vænne sig til Lysskjæret.
En pludselig Overgang fra Nattens Mulm til Middagssolens
stærke Lys er smertelig, Straalerne blænde og
det skjærer i Øinene; men lidt efter lidt kan Overgangen
skee, og da er Lyset velgjørende. Menneskene
have af Naturen Blik for Sandhed ligesom de have
Blik for Lyset, Trang til Sandhed ligesom de have
Trang til Lyset, dog med den Forskjel, som altid
finder Sted mellem det blot Naturlige og det Aandelige.
Det aandelige Blik og den aandelige Trang maa
vækkes, renses og skjærpes, det vil sige, der maa
Opdragelse til, for at Mennesket virkelig kan blive
Menneske.

Fra dette Synspunkt er Modsætningen mellem det
Aristokratiske og det Demokratiske iøinefaldende. Begge,
Aristokraten saavel som Demokraten, paavirke Mængden,
men i forskjellig Aand.

Demokraten gaaer ud fra den Forudsætning, at
Sandheden er noget Medfødt. Ligesom ethvert levende
Menneske veed, hvad Livet er, og hvad det er at leve,
saaledes veed ogsaa Enhver i Samfundet, hvad der er
Samfundet gavnligt og skadeligt. Sandheden bestaaer
i at skjelne imellem det Gavnlige og det Skadelige.
Men hvad er saa det Gavnlige? Det er jo lige%
% --- p. 82 ---
\setcounter{footnote}{0}%
\opage{82}frem. Gavnligt er det, som forhøier vort Velvære, som
skaffer os de største Friheder og de fleste Behageligheder.
Til at forstaae det, behøves der ingen besværlig
Underviisning. Det Gavnlige er det, vi Alle
attraae; at Attraaen er tilstede, sees deraf, at Driften,
Tilbøieligheden, ja Lidenskaben er tilstede. I Drift,
Tilbøielighed og Lidenskab sidder Mængdens Villie.
Mængden veed nok hvad den vil, naar den faaer det
at vide. Under denne Forudsætning veed Demokraten
da ogsaa hvad han vil: han vil naturligviis lade
Mængden vide, hvad det egenlig er den selv vil, det
Sande, det Tjenlige, det Gavnlige, det Bedste. Hans
Tale gaaer ikke ud paa at overbevise, men paa at
overtale. At overbevise er besværligt og, naar der
tales til Mængden, yderst kjedsommeligt. Men at
overtale er baade lokkende og oplivende. Til at overbevise
behøves der Grunde; men Grunde have ofte,
ligesom Lægens Piller, en bitter Smag, og saa ligge
de tungt i Bevidstheden og ere vanskelige at fordøie.
Til at overtale behøves der ikke egenlige Grunde, men
Slagord; Slagord virke i Forsamlingen som Pauker i
et Orchester, oplivende paa Sindet, musikalsk paa
Trommehinden.

Ogsaa Aristokraten maa henvende sig til Mængden,
thi i Mængden er Folket udvortes til. Hvis han i
Taushed holder sig tilbage og bilder sig ind, at Folket
ikke er andet end Mængde, at man kan ringeagte
Folket, fordi man i en vis Forstand kan oversee
Mængden, da er han kun en forloren Aristokrat, et
Gjenfærd fra Fortiden, med hvem vi intet have
at gjøre. Altsaa --- for Folkets Skyld --- maa ogsaa
Aristokraten tale til Mængden. Men hans Tale er en
% --- p. 83 ---
\setcounter{footnote}{0}%
\opage{83}anden end Demokratens. Aristokraten vil meer end
overtale, han vil overbevise. Han agter ikke paa
hvad der behager eller mishager Mængden, men alene
paa det, som efter hans velbegrundede Antagelse er
det Sande, det Tjenlige og Bedste. Han ønsker naturligviis,
at Mange, om muligt Alle, ville slutte sig
sammen om det, han anseer for det Bedste, men det
falder ham aldrig ind at ansee Noget for det Bedste,
fordi det har vundet Bifald hos Mange. Han er ikke
saa indbildsk at ansee sin egen Synsmaade for ufeilbar,
han er tværtimod villig til at give efter for Modgrunde;
men det nytter ikke at møde ham med den Indvending,
at Fleertallet er af en anden Mening end han, thi
han spørger ikke om Tal, men om Grunde. Det Aristokratiske
er at kunne staae midt i Mængden alene med
den erkjendte Sandhed, kun bøiende sig for Grunde,
men ikke for et Tal med mange Nuller.

Hvad Fleertallet antager, er sandt og gyldigt; det
er en demokratisk Grundsætning. Hvad der er sandt
og gyldigt bør Fleertallet antage, det er en aristokratisk
Grundsætning. Demokraten sætter Fleertallet over Sandheden;
Aristokraten sætter Sandheden over Fleertallet.
Gjælder det nu, at en Person staaer desto lavere, jo
lavere han selv sætter Sandhed og Ret, og desto høiere,
jo høiere han sætter Sandhed og Ret: saa følger jo
heraf, at Demokraten maa staae lavere i Forhold til
som han er sin Grundsætning tro, og Aristokraten desto
høiere, jo troligere han gjennemfører sin Grundsætning.
Men heraf forstaaes da ogsaa, hvad der ellers kunde
synes besynderligt, at Demokraten i Reglen er bedre
end hans Grundsætning, medens Aristokraten sjelden eller
aldrig ganske formaaer at gjøre sin Grundsætning Fyldest.

% --- p. 84 ---
\setcounter{footnote}{0}%
\opage{84}Folkelivets Fremgang vil altsaa beroe paa, at
Aristokratismen stadig voxer, og at Demokratismen
stadig aftager. Forvirringerne i det statsborgerlige
Samfund have, naar vi see nærmere til, altid deres Udspring
fra oven; fra Forvanskning og Misbrug af
Aristokratismen, og hvad mere er, fra Mangel paa Evne
og Villie til at faae de indsnegne Misbrug afskaffede,
før det er for sildigt, stammer Fordærvelsen. Den
franske Revolution afgiver lærerige Exempler.

Det var ikke for Folkets, men for Statens Skyld,
at Enevælden blev hævet til sit Høieste. Skjøndt det
menige Folk led usigelig derunder, saa fandt det sig
dog deri, saa længe Statens Magt og Storhed gav dets
Selvfølelse Næring. Men da hele Statsmagten var gaaet
op i Kongemagten, og den enevældige Kongemagt havde
gjort Sit til at ødelægge Staten, fulgte det jo af sig
selv, at Kongemagtens Grundvold maatte blive undergravet,
og Kongen personlig komme til at bære Skylden
for de offenlige Ulykker. Til Kongevælden støttede sig
Adelsvælden med dens Forrettigheder, dens æreløse
Hovmod og bestandig stigende Anmasselser; disse bleve
desto utaaleligere, jo dybere Staten sank. Kongen var
Aristokraternes Aristokrat, saa at man ikke kan finde
det forunderligt, at den offenlige Mening ansaae „Konge
og arvelig Adel for en af Guds værste Svøber.“ Under
det absolute Kongedømme blev Kirken og dens Hierarchie
i Grunden forvandlet til en Statsinstitution, thi Kongen
var i Virkeligheden den øverste Hierarch: den kongelige
Samvittighed var Undersaatternes Normalsamvittighed.
Var nu Staten „uden agtede Love, dens Vaaben uden
Glands, Hoffet uden Værdighed og Sæderne uden Ærbarhed“,
var desuden den høiere Geistligheds Umoralitet
% --- p. 85 ---
\setcounter{footnote}{0}%
\opage{85}--- man tænke paa en Cardinal Dubois, en Cardinal
Rohan og Flere --- i Gjennemsnit lige saa ryggesløs
som Adelens, dens Vantro lige saa fremtrædende som
Materialisternes og dens officielle Troesbekjendelse en
Maske, hvorunder den kun holdt fast paa de store
kirkelige Indkomster: hvad Under, om den sunde Forstand
ansaae Religionen for en mørk Overtro og dens
Præster for Bedragere, der burde hænges. Saaledes var
Fordærvelsen udgaaet fra oven.

Faren nærmede sig, og Forblindelsen voxte med
Faren. Allerede Madame Pompadour, Ludvig XV. alt
formaaende Herskerinde, skal have sagt: \emph{efter os
kommer Syndfloden}; og Syndfloden kom, og Vandene
steg over de høie Bjerge. Der var Mangel paa
Villie, fordi der var Mangel paa Forstaaelse, en forvirret
Brydning mellem Gammelt og Nyt gjorde det
Onde værre. „Til at tage Stødet af for alt dette havde
Frankrig kun Ludvig XVI. en rigtig brav, men lidet
duelig Mand, som, af Mistillid til sig selv, sædvanligviis
søgte Hjælp hos endnu mindre dygtige Personer, der
navnlig vare endnu mindre redelige. Han var nødt til
at skifte Ministre og med dem System, og naar de
slette skadede ham, saa gavnede de gode ham ikke.
Det kortsynede Hof, der var fulgt efter den grændseløse
Fordærvelse hos Ludvig XV. Hofmænd, forstod
ikke at stille Kongen i Spidsen for Røret og forlangte
derfor, at han skulde standse det, men uden at kunne
indgyde ham den nødvendige Kraft. Derfor viste Regjeringen
den Blanding af Vilkaarligked og Svaghed, % print: „Vilkaarligked“ (misprint for Vilkaarlighed)
som ophidser uden at kue Lysten til Modstand, men
tværtimod gjør denne populær og giver den Udsigt til
% --- p. 86 ---
\setcounter{footnote}{0}%
\opage{86}at lykkes. Omtumlet mellem Ministre, Hofmænd,
Kvinder, historiske Overleveringer og Philosophie, lod
Ludvig XVI. Statens Skib løbe for en halv Vind, og
han fængsler først Interessen, naar han hører op med
at handle og begynder at lide“ (Cantù).

Hvad Adel og Geistlighed angaaer var Stemningen
deelt mellem haardnakket Fastholden ved det Gamle
og letfærdig Leflen med det Nye. Det er saa fristende
at lade sig henrive af nye Ideer, saa længe de ikke
gribe forstyrrende ind i Ens egne medfødte Rettigheder
og indgroede Vaner. Man morede sig med Voltaire,
lod sig begeistre af Rousseau, loe ad „Figaros
Giftermaal“ og studerede Encyclopædisterne uden at
betænke, hvor farligt det er at lege med Ild og Lys
ved et aabent Krudttaarn. „Lafayette prædikede Korstog
til Amerika i Frihedens Navn og drog selv did for
at udgyde sit adelige Brod for Friheden. Unge Mennesker, % print: „Brod“ (misprint for Blod)
det franske Aristokraties Fremtidsstøtter, ilede
for at kæmpe for Kuldkastelsen af de Særrettigheder,
som bestode uforandrede i deres Fædreland, og for
at lære Grundsætninger om Lighed, om Had til al
Kongevælde, Ministervælde og Præstevælde.“

Parlamentet stred med Hoffet, og Autoriteterne
forraadte hverandre gjensidigt. Borgerstanden var for
en Deel kommet med; men Ingen, med Undtagelse af
nogle Philantroper, bekymrede sig om Menigmand.
Frihedens litterære Forkæmpere tænkte lidet eller
intet paa Folkets moralske Gjenfødelse; deres Lærdomme
vare beregnede paa den dannede Klasse og
angik ikke, hvad man dengang kaldte Bærmen (la
canaille).

% --- p. 87 ---
\setcounter{footnote}{0}%
\opage{87}Ikke desto mindre var det netop den lavere Befolkning,
den af Frihedsmændene foragtede Pøbel, som
ved de nye Lærdomme skulde komme til Magten. De
nye Lærdomme havde fornemmelig deres Styrke i det
blot Nedrivende; Kritiken, Tvivlen, Fornegtelsen vare
de virksomste Vaaben; det Positive og Nye indskrænkede
sig væsenlig til bittre Theorier. Forvovne Paradoxer,
dristig udslyngede Slagord, Tankelyn, der i et
Øieblik opklarede Mørket, fore hen over Masserne og
satte den store Hob i lidenskabeligt Røre. Under disse
Forudsætninger maatte Bevægelsen nødvendigviis blive
demokratisk.

Den radicale Demokratisme er kjendelig paa den
elementære Naturkraft, hvormed den angriber og nedriver
bestaaende Skranker, og paa den totale Mangel
paa Evne til at organisere noget Nyt ved at udvikle
de Spirer, der indeholdes i det Gamle. Den formaaer
ikke at skille Kjernen fra Skallen. Hvad der er behæftet
med Misbrug, var det end i sig selv det ædleste og
bedste, skal afskaffes. Kongemagten var misbrugt, og
saa afskaffede man Kongemagten; Adelsmagten var
misbrugt, og saa afskaffede man Adelen; Kirkemagten
var misbrugt, og saa afskaffede man Religionen; Intelligensen
var misbrugt, og saa bagvaskede man Videnskaben
og henrettede en Lavoisier. --- Terrorismen er
den radicale Demokratisme; Robespierre er den radicaleste
og consekventeste af alle Demokrater, hans Dyd
guillotinerer alle Misbrug.

Fødselsaristokratiet har i sin Tid misbrugt sin
Stilling, Embedsaristokratiet har misbrugt sin Stilling
og Underviisningsaristokratiet forsømt at gjøre
% --- p. 88 ---
\setcounter{footnote}{0}%
\opage{88}sin Pligt. Ved at oprulle Aristokratismens gamle
Synderegister giver Historien os en tilstrækkelig Belæring
om Demokratismens Opkomst. Det er ikke politiske
Slagord, men en besindig Opfattelse af Samfundets
historiske Udvikling, der kan lære os, hvorledes
det Gamle forholder sig til det Nye, hvad der under
Omdannelsen maa være det Bestaaende, hvad der maa
gjenfødes og hvad der maa bortdøe og forsvinde.

% --- p. 89 ---
\setcounter{footnote}{0}%
\lecture{VI.} {Politik og Samvittighed: Macchiavellisme.}
\opage{89}Politik og Samvittighed! Det er to store Magter i vor
Tilværelse. Kunde disse to Magter enes om at indgaae
„en hellig Alliance“ og holde den ubrødelig, hvad vilde
derved ikke være vundet for Menneskelighedens gode
Sag? Og hvad vilde der ikke være sparet? Midler
og Kapitaler, smertelige Offre, dyre Kræfter og Blod
og Taarer! Men desværre, de to Magter staae endnu
saa vidt adskilte, at det næsten seer ud, som hørte de
til to hinanden uvedkommende Riger Politiken tør % print: no stop after „Riger“
ikke spørge Samvittigheden tilraads, og Samvittigheden
kan sjelden eller aldrig godkjende Politikens Planer.
Dette maa dog ikke forstaaes, som om de kloge Mænd,
der sidde i Fyrsternes Raad, de Travle og Snilde, der
arbeide i Kabinetterne, vare samvittighedsløse Mennesker
fremfor os andre Dødelige. Langtfra! Det Samvittighedsløse
hidrører ikke umiddelbart fra Enkeltes
onde Sindelag. Det stammer oprindelig fra en Forvanksning % print: „Forvanksning“ (misprint for Forvanskning)
af Politikens Væsen. Den samvittighedsløse
Politik er et satanisk Væv, hvori Menneskeheden ---
ganske vist ved Tusindernes Brøde --- i Tidernes Løb
meer og meer har hildet sig selv.

% --- p. 90 ---
\setcounter{footnote}{0}%
\opage{90}Det nytter kun lidt at moralisere for Statsmændene
om at afvæbne, at være aaben og ærlig, altid at ville
det Rette o. s. v. De smile og trække paa Skuldrene
og deres talende Taushed er let at tyde. „Afvæbne,
naar Naboer og Modstandere ruste sig til det Yderste?“
„Være aaben og ærlig, naar Venner og Fjender anvende
alle Snildhedens og Snedighedens Omsvøb?“ „Ville det
Rette, det i og for sig Sande og Gode, naar alle Magter,
venskabelige og uvenskabelige, kun forstaae ved det
Rette, hvad Egennytten tilsiger dem om deres egen
Interesse og Ære?“ Politikeren veed, at „Verden ikke
styres med Fadervor“. Den Bitterhed, hvormed han
afviser Moralisten, er blandet med Veemod. Naturligviis
ville Alle hellere end gjerne ud af de piinlige Forviklinger;
men der er Masker i Skjebnens Net, som
ikke lade sig sønderrive; og Nettet omspænder Alle,
selv de Mægtigste sprælle i Garnet.

Er der da ingen Udvei? Jo ganske vist. Men til
at finde den og vælge den og følge den udkræves der
en Selvbesindelse, som kun kan opnaaes igjennem
smertelige Erfaringer.

Det første Spørgsmaal, der foreligger til Besvarelse,
er dette: hvorledes er det egenlig gaaet til, at Politik
og Samvittighed have fjernet sig saa vidt fra hinanden?
Det er et bekjendt Ord, at den, der vil Øiemedet, ogsaa
maa ville Midlerne. Heraf har man udledet, at Øiemedet
helliger Midlerne, og heraf igjen udledet, at slette Midler
kunne anvendes til Opnaaelse af et godt Øiemed. Dette
Sidste er aabenbart jesuitisk. Dog maa det bemærkes,
at Jesuiterne ikke ere første Opfindere; Misforholdet er
fra ældre Dato og af dybere Udspring.

% --- p. 91 ---
\setcounter{footnote}{0}%
\opage{91}Misforholdet mellem Midler og Øiemed kan først
erkjendes i sin Dybde, naar Øiemedet er evigt og Midlerne
timelige, Øiemedet helligt, men Midlerne verdslige.
Den gamle Verdens, navnlig Romernes, Politik var i
mange Henseender umoralsk, vilkaarlig, oprørende og
grusom; og dog er det Samvittighedsløse i al denne
Raahed og selviske Ubændighed kun Børneværk i
Sammenligning med den kirkelige Politik i Middelalderen.
Før Samvittigheden ret kan krænkes, maa den
være vaagnet, dens ideale Fordring anerkjendt, og en
almeengjældende Maalestok ved Bedømmelsen af de
virkelige Handlinger fastsat. Men den gamle Verdens
Samvittighed var indhyllet i naturlig Selvviskhed og
rørte sig kun i mythiske Drømme. Det er ikke blot
en Nero, der er samvittighedsløs; men selv om Ethiken
vender sig til en Cato den Ældre, sprætter ham op og
seer ham indvendig efter i Sømmene, vil den neppe
finde Spor af det vi maae kalde Samvittighed. Det
eneste Oldtidsfolk, som i strengere Forstand kunde siges
at have Samvittighed, var Jøderne; og dog er Jødefolket
bekjendt for at have anvendt en god Deel af sin
naturlige Forstand paa at omgaae Loven og snyde
Jehova. Først med Kristendommen vaagner Samvittigheden
i Historien.

Det Karakteristiske ved den katholske Kirkepolitik
er nu let at belyse. Kirken maatte naturligviis indskærpe
Samvittighedens Fordringer; det Formaal, den
satte sig, var oververdsligt, dens Øiemed helligt.
Ikke desto mindre skulde dette hellige Øiemed
dog, ifølge den katholske Synsmaade, fremmes og naaes
ved verdslige Midler. Midler og Øiemed ere da saa
ueensartede, som det vel er muligt. Jo mere Kirken
% --- p. 92 ---
\setcounter{footnote}{0}%
\opage{92}blev indviklet i verdslige Handeler, desto mere udviklede
sig dens Politik. Christus var ingen Politiker; de
eneste Verdensmidler han anvendte, vare Lidelse og
Selvopoffrelse; og Midlerne svarede til Øiemedet, thi hans
Rige var ikke af denne Verden. Men da Kirken havde
sat sig det høie Maal at verdsliggjøre det Oververdslige,
kunde den ikke lade sig nøie med saa ringe Midler.
Til Verdsliggjørelse af det Oververdslige maatte den i
Alvor være betænkt paa at bruge verdslige Midler, den
maatte drive Politik. Naar nu Samvittigheden krympede
sig ved Midlerne, men dog bifaldt Øiemedet, „det
hellige Øiemed“, da laa det nær at tiltale den praktisk
og foreholde den, at Enhver, som vil Øiemedet, nødvendigt
maa ville Midlerne. Fandt Samvittigheden sig oprørt
ved de Midler, som anvendtes, hvor trøsterigt dulmende
var da ikke Forvisningen om, at Kirken havde Soningsmidler,
at Øiemedet helliger Midlerne, at man med god
Samvittighed kan anvende slette Midler i gode Øiemed.
Den pavelige Politik skulde, efter sit Ideal, være en
Samvittighedspolitik, dens driftigste Midler Samvittighedsmidler;
men i Virkeligheden var den saa egennyttig,
saa vilkaarlig og rænkefuld, at den ordinære Fyrstepolitik
havde ondt nok ved at maale sig med den i
Samvittighedsløshed: det var en skrigende Modsigelse.

Det næste Skridt var, at Samvittigheden gik ud, og
Politiken blev ene tilbage. Dette Vendepunkt er med
klar Bevidsthed angivet af Florentineren Nicòlo
Macchiavelli (1469—1527). Han fik ved at deeltage i
det offenlige Liv, tidlig Smag paa Politiken og indøvede,
som det synes, selv i private Forhold den Grundsætning,
at man til Opnaaelse af et Øiemed ikke maa tage
det saa nøie med Midlerne. Han førte en fordomsfri,
% --- p. 93 ---
\setcounter{footnote}{0}%
\opage{93}ubunden Tilværelse, var forelsket ikke blot i sin Ungdom,
men selv i en fremrykket Alder, sang Serenader,
skrev uanstændige Komedier, og var i den Grad Kammeraternes
gode Hoved, at de, under hans Fraværelse som
landsforviist, skreve til ham: „Nu da I ikke er her
mere, forstaaer man sig hverken paa Spil, Drik eller
anden Morskab.“ Af Kirken lærte han, at den, der
ved List og Magt vil gjøre Lykke paa det store Theater,
maa være en god Skuespiller. For at prøve sine Kræfter
i andægtige Roller, sluttede han sig til et af de
fromme Broderskaber og holdt, da Raden kom til ham,
en Bodsprædiken, der endte med en hjerteknusende Opfordring,
fuld af Salvelse: „Vi bør efterligne St. Franciscus
og St. Hieronymus, som for at spæge Kjødet og berøve det
Kraft til at friste dem, væltede sig i Torne eller sloge sig
for Brystet med en Steen; vi derimod ere bedragne af
Vellyst, indlullede i Vildfarelser, hildede i Syndens
Snarer og ere i Djævelens Hænder. Altsaa maae vi,
for at slippe ud deraf, gjøre Bod og raabe med David:
miserere mei Deus, eller græde med St. Peder.“

Var Religionen et Kirkebedrag, et Skuespil, hvad
var saa Sædeligheden? I sin lige saa berømte som berygtede
Bog \emph{Fyrsten} (il principe) gjør Macchiavelli Rede for
den politiske Retfærdighed. „Herskeren skal stedse føre
i Munden Ord som Retfærd, Lovlighed, Mildhed og
Religion, men ikke bekymre sig om disse Ting, naar
han har Fordeel af det Modsatte. Grusomhed er nødvendig
for en ny Regjering, og man skal hellere gjøre
sig frygtet end elsket, naar man ikke kan blive begge
Dele paa een Gang. En Regjerings Maal er dens Varighed,
og det kan ikke naaes uden Haardhed, fordi
Menneskene sædvanligviis ere utaknemmelige, falske og
% --- p. 94 ---
\setcounter{footnote}{0}%
\opage{94}kivagtige, saa at de maae holdes i Tømme ved Frygt
for Straf. Fyrsten maa desaarsag forstaae at bruge
baade Gift og Dolk, baade List og Vold, han maa kunne
spille baade Ræven og Løven.“ Hvor alvorlig Macchiavelli
mener det, og hvor afvigende hans Opfattelse af
Retfærdigheden er fra den sokratiske i Platos Stat, sees
blandt Andet af „Samtalerne over Livius's Decader“,
hvori han ofte henviser til \emph{Fyrsten}. Han erklærer
aabenlyst, at Retfærdigheden er opstaaet ved, at Menneskene
have seet, hvor nyttigt det Gode og hvor skadeligt
det Onde var, at Menneskene kun lade sig bevæge
til det Gode af Nødvendighed og Egennytte“, o. s. v.
Mærkværdigt nok, at de Grundsætninger, som „Hedningen“
Plato lægger i Munden paa en uforskammet
Sophist, netop i en katholsk-christelig Tidsalder blive
lyste i Kuld og Kjøn af „Bodsprædikanten“ Macchiavelli!
Det er en stor Bevidsthedsforvandling, et verdenshistorisk
Udbytte, vi her have for os.

Macchiavellismen har det tilfælles med den gamle
Verdens politiske Anskuelse, at den sætter Staten som
det Høieste, det Eneste, der har Gyldighed i sig selv.
Den græske og romerske Historie har macchiavellistiske
Træk i Ord saa vel som i Handling. At Macchiavellismen
ikke desto mindre er forskjellig fra den hedenske
Statsbetragtning, følger af, at den har Christendommen,
særlig Katholicismen, til Forudsætning. Naar en hedensk
Sophist paastaaer, at Intet i og for sig er godt
eller ondt, kan han dog paa en Maade undskyldes med,
at den Samvittighed, han kjender, er et forblommet
Væsen; men naar en Macchiavelli aabenlyst fremsætter
den samme Paaatand, da skeer det med Haan % print: „Paaatand“ (misprint for Paastand)
og Fornegtelse af den tydelig erkjendte Samvittighed.
% --- p. 95 ---
\setcounter{footnote}{0}%
\opage{95}Naar en Statsmand i Oldtiden betragtede Samfundets
Individer som retløse over for Staten og desaarsag ikke
tog i Betænkning at offre dem i Tusindviis for en eller
anden politisk Hensigt, da bliver det Oprørende heri
formildet ved Forestillingen om, at det Personlige i
Menneskeværd og Menneskeret endnu var hyllet i mythisk
Taage; men naar Macchiavellismen med utilsløret
Anprisning af Falskhed, Forstillelse, Edsbrud og Brodermord
koldsindigt og følgestrengt udvikler den samme
Grundtanke, da betyder det, at Idealet af Menneskets
personlige Værd kun har viist sig i Dagens Lys, for
med skamløs Frækhed at blive traadt under Fødder.

Man kan indvende, at hverken den kirkelige eller
den verdslige Politik nogensinde har vedkjendt sig
Macchiavellismen endnu mindre gjennemført dens Conseqvens.
Men hvad betyder en saadan Indvending? At
vedkjende sig Macchiavellismen ærligt og aabent, vilde
jo være det Dummeste en Macchiavellist kunde gjøre.
Hvad Gjennemførelsen angaaer, kunne en Ludvig XI i
Frankrig og en Ferdinand den Katholske i Spanien
vel betragtes som talentfulde Forløbere i Praxis,
og hvad Theorien selv angaaer, er det en Kjendsgjerning,
at Katharina af Medicis, da hun blev formælet
med Henrik II., bragte Bogen om „Fyrsten“
med sig fra Italien. At hun har studeret den
grundigt, og at Frugterne af hendes Studier have paa
en glimrende Maade lagt sig for Dagen, særlig ved
Blodbrylluppet i Paris (Bartholomæusnatten 1572), er
bekjendt. Saa længe den verdslige Politik havde et
religiøst Tilsnit, og Fanatismens Brintflamme forstærkede
Heden, var Temperaturen i høi Grad gunstig for
Macchiavellismens Blomstring.

% --- p. 96 ---
\setcounter{footnote}{0}%
\opage{96}Det Kirkelige bødede paa det Verdslige, naar Fyrsten
øvede en tyrannisk Udaad, og Pavens hellige
Haand afvaskede Blodpletten (det stokholmske Blodbad).

I hvilken Grad Katholikerne havde Sands for
Macchiavellis Lære, bliver os klart, naar vi see, hvor
skarpsindigt, fromt og snildt den blev opfattet og
videre udviklet af Jesuiterne: Jesuitismen er, politisk
forstaaet, en kirkelig anvendt Macchiavellisme. At
den religiøse Macchiavellisme, endog uden Hensyn til
det Kirkelige, er blevet tilegnet af Protestanterne, beviser,
hvor dybe Rødder den har i den menneskelige
Natur. Da Karl X (Gustav) havde brudt Roskildefreden,
og i den Hensigt at tilrane sig Danmark-Norge
var landet med sin Flaade ved Korsør, istemte han
Psalmen: „Paa dig jeg haaber, o Herre kjær!“ Kan
det lykkes en Fyrste at krænge sin protestantiske Samvittighed
i den Grad, at Gud i Himlen selv bliver
Macchiavellist, da kan han ogsaa med Karl Gustav
uden at blues tale om sine Erobringer som „Gaver, han
skylder Guds Naade.“ I denne Politiks Aand har Karl
Gustavs fortrolige Ven Schlippenbach i Fortrolighed
meddeelt Kurfyrste Frederik Vilhelm af Brandenborg en
opbyggelig Fortolkning af Naadens Aabenbaring i Politiken:
„Gud aabenbarer ikke mere sin Villie for Fyrsterne
igjennem Propheter eller Drømme; men naar der
viser sig en god Leilighed til at skade Naboen, bør
det holdes for et guddommeligt Kald.“ --- Den store
Kurfyrste har uden al Tvivl følt sig saa opbygget ved
Fortolkningen, at han ikke har seet nogen Hindring
for en oprigtig Union mellem Lutheranere og Reformeerte.


% --- p. 97 ---
\setcounter{footnote}{0}%
\opage{97}Vel fremkom der i Litteraturen efterhaanden en
Deel moralske Gjendrivelser af Macchiavellis Lære; men
Praxis var stærkere end Theorien. Det gjælder ikke
mindst om den brandenborgske Praxis. Frederik II. af
Preussen skrev i sin Ungdom en Antimacchiavelli, et
Skrift fuldt af dydig Indignation over de umoralske Grundsætninger;
men Preusserkongen, hvis politiske Samvittighed
i det Hele synes at have befundet sig i en
noget mythisk Tilstand, kom netop ved dette Ungdomsarbeide
til at bryde Staven over Grundsætninger, ved
hvis snilde Anvendelse han siden som Schlesiens Erobrer
og Syvaarskrigens Helt skulde glimre i Historien. Denne
Inconseqvens er fra et macchiavellistisk Synspunkt for
Øvrigt ret conseqvent.

I vort Aarhundrede kan vel Macchiavellismen
gjennemsnitlig siges at være i Aftagende. Kun nu og
da er det Umoralske i Politiken kommet til Gjennembrud
paa oprørende Maade. For ikke at komme Nutiden
for nær og oprive de Saar, som endnu bløde,
ville vi indskrænke os til at minde om Engellændernes
Overfald paa Kjøbenhavn 1807. Dette Skridt var for
den naturlige Følelse saa oprørende, at en dansk Digter
endog har villet gjøre os anskueligt, hvilken Misbilligelse
det maatte vække hos selve det engelske Folk.
Bekjendt er Sangen om den engelske Kaptain: „med
Lyst jeg raabte: Gud og Kong Georg! Hvergang jeg
fyred paa Napoleons Brigger, Men her jeg næsten
græde maa af Sorg: Paa Røveri jo nu vor Flaade ligger.
--- Seil I! jeg lyder ikke Kongens Bud, Men gaaer
paa andre Eventyr og Reiser, Der staaer i Skriften:
Du skal lyde Gud Meer selv, end Du skal lyde Romas
Keiser!“

% --- p. 98 ---
\setcounter{footnote}{0}%
\opage{98}Det være med Sagnet, som det vil; vist er det, at
den engelske Politik ikke viste sig Macchiavellismen
uværdig. Til et saa stort Hverv udkræves der Mænd
af den rette Skole. Det engelske Ministerium havde i
Sir Brook Taylor og Sir Francis Jackson netop to saadanne
Hædersmænd. Sir Brook Taylor havde allerede
et Navn i den store Verden. „Han havde nemlig tidligere,
da han repræsenterede England i Kassel, deeltaget
i Forsøgene paa at myrde Førsteconsul Bonaparte“;
Øiemedet helliger Midlet. Hvad Sir Francis
Jacksons fremtrædende Personlighed angaaer, har en af
vore yngre Historikere\footnote{Alex. Thorsøe Den danske Stats politiske Historie. Første Tidsrum. Kbhvn. 1873. S. 154 flg.} leveret en Karakteristik, der
fortjener at læses. Som en \emph{Deus} machina (en Gud af % print: „Deus machina“ (sic)
en Maskine) optræder denne diplomatiske Agent pludselig
i Kiel og aabenbarer Kronprinds Frederik Englands
skjulte Planer. Han „begyndte sit Budskab med
den Erklæring, at den engelske Regjering havde forskaffet
sig paalidelige Vidnesbyrd om, at Napoleon
vilde tvinge Danmark til at lukke sine Havne for Englænderne
og til at tage Deel i Fastlandsspærringen.
Saa vel Storbritanniens Interesse og Sikkerhed som
Danmarks Selvstændighed fordrede, at Hoffet i Kjøbenhavn
traadte i et inderligt Forhold til England og samtykkede
i, at dets Flaade blev ført til England, førend
Frosten havde indespærret den i Havnen, for at den
kunde være beskyttet mod Napoleons Efterstræbelser.“
Han forsikkrede høitideligt, „at Flaaden skulde blive
tilbagegiven, naar Freden var sluttet.“ Han tilbød
Kronprindsen et nøie Forbund med Storbritannien,
% --- p. 99 ---
\setcounter{footnote}{0}%
\opage{99}Garantier o. s. v., og tilføiede endelig, „at hvis den
danske Regjering ved et saadant Skridt frygtede for at
bringe sig i en farlig Stilling over for Frankrig, saa
maatte den dog ikke glemme, at den engelske Styrke,
der nu laa foran Kjøbenhavn, var saa betydelig, at det
godt kunde see ud, som om den havde givet efter for
Tvang, hvortil England iøvrigt vilde tye, dersom Danmark
ikke gik ind paa dets Planer.“ --- Forslag af et
saa fyndigt Indhold maatte naturligviis fremsættes i en
tilsvarende fyndig Form, og Sir Francis Jackson var en
Mester i den fyndige Form. Der berettes, „at Jackson
i Begyndelsen ikke sagde andet end „De“ til Kronprindsen,
hvorpaa denne spurgte, om han ikke vidste,
hvem det var, han talte med, at det var en Prinds,
Kongen af Danmarks Søn, hvorpaa Jackson blev noget
høfligere“ o. s. v.

Den yderliggaaende Brutalitet, hvormed Macchiavellismens
Apostle ved denne Anledning meldte sig,
havde desværre til Følge, at Ærlighedens og Moralitetens
værdige Repræsentant, den ædle Kronprinds, ikke
var i Stand til at komme til politisk Besindelse, før
det var for sildigt. „Kronprinds Frederik burde vistnok
have søgt at vinde Tid ved at knytte Underhandlinger
med Jackson, indtil den franske Regjerings Hensigter
vare blevne ham bekjendte. Han vilde da, naar dette
var skeet, være kommet til den Erkjendelse, at Keiserens
Planer med Hensyn til Danmark vare lige saa
stridende mod Folkeretten og lige saa lidt lode sig
forene med den danske Regjerings frie Selvbestemmelsesret
som det britiske Kabinets. Kronprindsen havde da
bragt i Erfaring, at Napoleon med Troppeconcentreringen
ved Holsteens Grændse, ganske forbandt den samme
% --- p. 100 ---
\setcounter{footnote}{0}%
\opage{100}Tanke, som det engelske Ministerium med Flaadens
Samling i Øresundet, og det vilde saaledes have staaet
aldeles klart for ham, at de to Magter i lige høi Grad
vare ledede af Egennytte, og at de begge med lige stor
Hensynsløshed og Tilsidesættelse af Danmarks Handlefrihed
kun gjennem Vold og Tvang søgte at føre deres
egen Villie igjennem. Men Kronprinds Frederik indlod
sig slet ikke i Underhandlinger med Jackson, og kan
man end ikke beundre hans „diplomatiske Kløgt“, saa
faaer man et desto mere tilfredsstillende Indtryk af hans
moralske Karakteer. „Han erklærede --- hvad der i
hiin Tid, da politiske Voldshandlinger hørte til Dagens
Orden, klang besynderligt --- at der neppe i Historien
fandtes noget Exempel paa et Angreb som det, hvormed
man vovede at true Danmark, han udtalte, at man hos
de barbariske Stater kunde vente at finde større Redelighed
end hos den engelske Regjering, paa hvis Hjælp
man ikke engang kunde stole, selv om den var Forbundsfælle.
Da Sir Francis Jackson nedlagde Indsigelse mod
denne Udtalelse, yttrede Kronprinds Frederik, at den
Mand, der havde moralsk Mod til at paatage sig en
Sendelse som hans, maatte ogsaa kunne taale at høre
Alt; da det engelske Sendebud forsikkrede, at England
med Penge vilde erstatte Danmark den lidte
Skade, blev han afbrudt med de navnkundige Ord: „Og
hvormed vil man erstatte Danmarks Ære?“

Macchiavellismen indskrænker sig ikke til visse
Statsforfatninger; den trives omtrent lige godt i dem
alle. --- Demokraterne ville gjerne bilde os ind, at naar
Flertallet kommer til Roret, saa vil de store Uretfærdigheders
sidste Time slaae, og Politik og Samvittighed
broderlig favne hinanden; thi under Fleertalsregi%
% --- p. 101 ---
\setcounter{footnote}{0}%
\opage{101}mentet skal det gaae ærligt til med Midler og Øiemed.
Hvor det er skjønt, naar man blot kunde troe det!
Men desværre, Historien lærer, at Macchiavellismen
netop i Fleertalsmagten har en gammel Forbundsfælle,
og det er Parti-Aanden. „Den, der vil Øiemedet, maa
ville Midlerne; Midlet er godt, naar det fører til
Maalet.“ Partiet gjør de Enkelte til Midler, bruger
dem, saa længe de tjene dets Øiemed, og kaster dem
bort, naar det ikke behøver dem længere. Omvendt
gjøre Deeltagerne igjen Partiet til Middel, og forlade
det, naar de have brugt det som Skammel. Det er
just det Liberale, det ægte Frisindede ved Parti-Aanden,
at dens Krav paa de Deeltagendes indbyrdes
Enighed ingenlunde ere nærgaaende. Den er langt fra
at forlange, at Partiets Medlemmer skulle være enige
om noget saa væsenligt Indvortes, at det kunde genere
Samvittigheden; det er nok at de ere enige om et eller
andet Udvortes, som at styrte et Ministerium, at hjælpe
deres Førere til Magten, at tage Hevn over Modstandere
o. desl. Ethvert saadant udvortes Øiemed er i
Reglen kun et Middel for andre og høiere Øiemed.
Enhver i Partiet kan da frit gaae omkring og kjæle
for sin Samvittighed og bilde den ind, at de Smaakneb,
Partiet nu og da tillader sig, kun ere uskyldige Fif,
hvorved man baner en Vei for Aanden. Vel er der
Mange i Partiet, som endnu ikke have Aanden, men
Partiet er jo ogsaa kun en Overgang, et Middel. Vel
ere nogle af de Midler, Partiet anvender, maaskee ikke
ganske lovlige, og vilde neppe fortjene Samvittighedens
Bifald, hvis det var Modstanderne, der brugte dem;
men man kan --- i Betragtning af den politiske Situation
tillade sig Adskilligt, naar det skeer i et godt Øiemed.
% --- p. 102 ---
\setcounter{footnote}{0}%
\opage{102}„Modstanderne have --- lykkeligviis --- intet godt Øiemed.
Der er i, hvad de sige og gjøre, vist nok Noget, vi
vilde finde rigtigt, dersom det var os, der sagde det,
os selv, der gjorde det, men fra Modstandernes Side er
det --- paa Grund af den politiske Situation --- lige
fuldt ufornuftigt og uhæderligt; thi Sagen er, at Partiets
Modstandere uden Undtagelse ere Idioter.“ Ved
slige Betragtninger kan en skikkelig Mand let faae sin
Samvittighed chloroformeret, endog i den Grad, at selv
de dristigste Operationer med den politiske Kniv ikke
volder den Smerte. Det er den allerbetænkeligste Art
af Macchiavellisme.

Da Herkules, beretter Mythen, paa sit vestlige Tog
kom til Lybien, stødte han paa Kæmpen Antæus, hvis
Grusomhed mod de Fremmede gjorde ham til et sandt
Uhyre. Antæus tvang nemlig de Ankomne til at brydes
med sig, og naar han havde overvundet dem, kvalte
han dem og plantede Hovedpanderne op omkring sin
Bolig. Det, der gjorde ham uovervindelig i Kampen,
var hans Moder, \emph{Jorden}, hvorved hans Kraft fordobledes.
Men Herkules greb ham med Armene om
Livet og holdt ham svævende i \emph{Luften}, indtil han;
kvalt ved Heltens Kraft, opgav sin Aand. --- Macchiavellismen
er som Kæmpen Antæus; den kan ikke beseires
ved blot verdslige Kræfter; thi Verdsligheden er
dens Moder, dens inderste Væsen er det selvisk Verdslige.
Skal Uhyret kvæles, da maa det løftes op i Religionens
og Sædelighedens Sphære, i den Luftkreds,
hvori Samvittigheden aander, og holdes deroppe af mægtige
Arme. Ingen diplomatisk Kunstner og intet politisk
Geni formaaer at holde Stand mod Uhyret; det
blænder de Kløgtige med deres egen Kløgt og kvæler
% --- p. 103 ---
\setcounter{footnote}{0}%
\opage{103}de Stærke i kvalme Dunster. Skal der gives nogen
Magt paa Jorden, som formaaer at gribe Uhyret om
Livet og holde det svævende i Luften, da maa det være
Folke-Aanden, naar den rigtig vaagner. Folke-Aanden er
en Herkules af moralske Kræfter, en virkelig Aand i et
aandeligt virksomt Folk. Det er en Helt, der gaaer
ærlig til Værks; i hans Bannermærke staaer Menneskelighedens
Krav med luende Skrift, og Verdslighedens
List og Politikens Underfundighed maa skamfuld bøie
sig derunder. Kommer han en Gang paa sit Tog til
Brydekampen, da vil han sikkert bruge redelige, om
just ikke lemfældige Midler: han vil kue den smaalige
Trods, lufte ud i Partiernes Leire og rense den belemrede
Stat, som hans mythiske Forgænger fordum
rensede Augias' Stald.

\end{document}
